% concatenated from arxiv.tex, packages.tex
\documentclass{article}
\PassOptionsToPackage{numbers, compress}{natbib}
\usepackage[preprint]{neurips_2026}
\usepackage[utf8]{inputenc}
\usepackage[T1]{fontenc}
\usepackage{hyperref}
\usepackage{url}
\usepackage{booktabs}
\usepackage{amsfonts}
\usepackage{nicefrac}
\usepackage{microtype}
\usepackage{xcolor}
\usepackage{titletoc}
\bibliographystyle{plainnat}
\usepackage{amsmath,mathtools,bm}
\usepackage{amsthm}
\usepackage[T1]{fontenc}
\usepackage{microtype}
\usepackage{enumitem}
\usepackage{booktabs,tabularx,array, makecell}
\usepackage{changepage}
\newcolumntype{P}[1]{>{\raggedright\arraybackslash}p{#1}}
\newcolumntype{Y}{>{\raggedright\arraybackslash}X}
\hypersetup{colorlinks=true, linkcolor=blue!50!black, urlcolor=blue!50!black, citecolor=blue!50!black}
\newtheorem{theorem}{Theorem}
\newtheorem{proposition}[theorem]{Proposition}
\newtheorem{lemma}[theorem]{Lemma}
\newtheorem{corollary}[theorem]{Corollary}
\newtheorem{remark}[theorem]{Remark}
\newtheorem{definition}[theorem]{Definition}

\newcommand{\X}{\mathcal{S}}
\newcommand{\A}{\mathcal{A}}
\newcommand{\Fcal}{\mathcal{F}}
\newcommand{\Pcal}{\mathcal{P}}
\newcommand{\R}{\mathbb{R}}
\newcommand{\E}{\mathbb{E}}
\newcommand{\Prob}{\mathbb{P}}
\newcommand{\Tc}{T^{\pi}}
\newcommand{\norm}[1]{\left\lVert #1 \right\rVert}
\newcommand{\ip}[2]{\left\langle #1, #2 \right\rangle}
\newcommand{\Ustar}{U^{\star}}
\newcommand{\Vstar}{V^{\star}}


\title{A Finite-Iteration Theory for Asynchronous Categorical Distributional Temporal-Difference Learning}

\author{Ege C.~Kaya, Abolfazl~Hashemi  \\
  Elmore Family School of Electrical and Computer Engineering\\
  Purdue University\\
  West Lafayette, IN 47906 \\
  \texttt{kayae@purdue.edu, abolfazl@purdue.edu}}

\begin{document}
\maketitle

\begin{abstract}
We study finite-iteration behavior of the exact asynchronous recursions used by categorical distributional temporal-difference methods. The analysis covers scalar categorical TD in the Cram\'er geometry and multivariate signed-categorical TD in the maximum mean discrepancy geometry. Existing statewise isometric embeddings turn both methods into single-state stochastic-approximation recursions that contract in a block-supremum norm, but the categorical operators are contractive only on invariant representation domains. We establish the required restricted-domain theory and obtain discounted bounds under i.i.d.\ sampling and under a Markovian trajectory. A Poisson-equation decomposition handles trajectory dependence without an explicit mixing-time window. For undiscounted fixed-horizon policy evaluation, we establish analogous finite-iteration guarantees for horizon-stacked categorical methods under episodic sampling. Together, these results provide a unified non-asymptotic analysis of asynchronous categorical distributional TD across scalar, multivariate, discounted, and fixed-horizon settings.
\end{abstract}

\section{Introduction}
Categorical representations are among the central approximation families in distributional reinforcement learning (RL) \citep{bellemare2017distributional,rowland2018analysis,dabney2018distributional,dabney2018implicit,yang2019fully,lyle2019comparative,rowland2019statistics,bellemare2023distributional}. In the scalar case they underlie the categorical methods initiated by C51 \citep{bellemare2017distributional}, while in the multivariate case they support signed-measure constructions for vector-valued returns \citep{wiltzer2024foundations}. These methods replace infinite-dimensional return-distribution objects with finite-dimensional ones while preserving the geometry that governs projected distributional policy evaluation \citep{howard1972risk,jaquette1973markov,jaquette1976utility,sobel1982variance,morimura2010nonparametric}. For these categorical methods, the asymptotic picture is by now well understood \citep{gurvits1994incremental, dayan1994td, tsitsiklis1994asynchronous, jaakkola1993convergence, bertsekas1996neuro, littman1996generalized, rowland2018analysis,bellemare2023distributional}. The scalar categorical temporal-difference (CTD) method admits a projected Bellman contraction in the Cram\'er metric together with asymptotic guarantees \citep{rowland2018analysis,bellemare2023distributional}. The multivariate signed-categorical temporal-difference (MTD) method admits the analogous maximum mean discrepancy (MMD)-based contraction and theory \citep{wiltzer2024foundations}. What remains less well understood is the finite-iteration behavior of these TD methods.

Standard TD learning is usually not a synchronous full-state procedure \citep{sutton1988learning,sutton1998reinforcement,tsitsiklis1996analysis,borkar1998asynchronous}. In practice, each update modifies only the sampled state, and in online RL, the sampled states are generated by a Markovian trajectory. Finite-iteration guarantees for asynchronous, trajectory-driven updates are therefore the natural benchmark for understanding how quickly CTD and MTD approach their projected fixed point in regimes relevant to practice. Several recent results address neighboring finite-iteration questions \citep{qu2020finite,bhandari2018finite,srikant2019finite,patil2023finite,dalal2018finite,peng2024statistical,boeck2022speedy, zhang2023estimation, rowland2024near, peng2025finite}. In particular, the analysis of \citet{peng2025statisticalefficiencydistributionaltemporal} covers one-state CTD with independent samples and Markovian variants that use data dropping or variance reduction. What remains open in that line is a finite-iteration analysis of the unmodified every-transition Markovian recursion, as well as corresponding guarantees for signed-categorical MTD and undiscounted fixed-horizon learning.

The present paper is organized in two parts. The discounted part sets off from the contractive stochastic approximation (SA) tools of \citet{chen2020finite,chen2022finite,chen2024lyapunov} and makes them valid on the invariant categorical domains where the CTD and MTD estimates live. For Markovian trajectory sampling, we combine the smoothed block-supremum potential with a Poisson equation rather than a mixing-time comparison. The fixed-horizon part requires an entirely different argument, involving the freezing of each episode at its boundary, aggregating phase-dependent drifts through cumulative occupancy, and controlling the error created by the online within-episode updates.

\noindent\textbf{Sampling regimes.} The i.i.d.\ model is representative of a replay-based or generative-model benchmark, where updates are formed from independent samples drawn from a buffer or simulator. Markovian sampling follows one online behavior trajectory, standard for online TD learning. In the finite-horizon undiscounted case, the episodic regime is the standard reset-based formulation, where interaction proceeds for a fixed horizon $H$ and then restarts from an initial distribution.

\noindent\textbf{Contributions.}
\begin{enumerate}[leftmargin=*,labelsep=0.2em, align=left, itemsep=2pt]
    \item We prove finite-iteration bounds for the exact discounted asynchronous CTD and MTD recursions. The i.i.d.\ result is a direct specialization of contractive SA once the categorical domain issue is resolved. The Markovian result uses a Poisson equation decomposition to retain the iteration-rate exponents obtained under i.i.d.\ sampling, without an explicit logarithmic mixing-window loss.
    \item For fixed-horizon policy evaluation, we establish finite-iteration guarantees under episodic sampling and quantify the effects of state coverage and horizon length on convergence.
    \item We isolate the method-specific perturbation geometry. CTD has a uniform support-radius bound, whereas signed-categorical MTD has affine growth. The resulting constants expose their dependence on coverage, contraction gaps, trajectory dependence, and support geometry.
\end{enumerate}

\section{Related work}
Our finite-iteration arguments build on existing contractive SA theory. \citet{chen2020finite,chen2022finite,chen2024lyapunov} develop smooth-envelope analyses for arbitrary contraction norms under i.i.d.\ and Markovian noise, extending the broader SA and TD literature \citep{robbins1951stochastic,ljung2003analysis,borkar2008stochastic,kushner2003stochastic,bhandari2018finite,qu2020finite}. For the discounted Markovian result, we use the bounded-state Poisson-equation theorem of \citet{haque2024stochastic}, which avoids the explicit mixing window appearing in earlier Markovian SA bounds. Our additional work verifies these tools on invariant categorical domains and derives the method-specific perturbation constants.

The distributional ingredients also come from established foundations. \citet{rowland2018analysis,bellemare2023distributional} provide the Cram\'er geometry, categorical projection, and asymptotic CTD theory. \citet{wiltzer2024foundations} develop the signed-measure MMD geometry and almost-sure asynchronous MTD convergence. \citet{de2020fixed} motivate the horizon-indexed stack and online full-stack update in the fixed-horizon setting, but do not provide our episode-boundary distributional rates. 

Non-asymptotic distributional results answer neighboring questions under different sampling access or estimators. The analysis of \citet{peng2024statistical} gives high-probability rates for synchronous distributional TD and CTD with a generative model. The expanded analysis in \citep{peng2025statisticalefficiencydistributionaltemporal} also treats independent one-state sampling and Markovian data, using data dropping or a variance-reduced multi-epoch algorithm in the Markovian case. \citet{boeck2022speedy} analyze an accelerated categorical method, while \citet{zhang2023estimation,rowland2024near} study model-based and direct estimators. \citet{wu2023distributional} consider offline evaluation, \citet{peng2025finite} treat linear approximation, and \citet{kastner2025categorical} study KL-based categorical dynamics asymptotically. Our i.i.d.\ result is a contractive-SA analysis of the one-state categorical recursion rather than a competing minimax result. The distinct discounted contribution is the finite-iteration analysis of unmodified CTD and MTD along every transition of one Markovian trajectory. The fixed-horizon guarantees address a separate episodic setting.

\section{Discounted categorical policy evaluation}\label{sec:discounted}
We consider a discounted Markov decision process \citep{sutton1998reinforcement} $(\X,\A,P,R,\gamma)$ with finite state and action spaces $\X$, $\A$, reward function $R(s,a)$, transition kernel $P(\cdot\mid s,a)$, and discount factor $\gamma\in(0,1)$. We assume that $R:\X \times \A \to [0, 1]$ in the scalar, and $R:\X \times \A \to[0, 1]^q, q\ge2$ in the multivariate case. A policy $\pi$ is fixed throughout. At iteration $k$, let $S_k$ denote the queried state, $A_k\sim\pi(\cdot\mid S_k)$, write $R_k:=R(S_k,A_k)$, and let $S'_k\sim P(\cdot\mid S_k,A_k)$ denote its sampled successor. In the i.i.d.\ regime, the tuples $(S_k,A_k,R_k,S'_k)_{k\ge0}$ are i.i.d.\ with query-state marginal $S_k\sim\rho$, and $S'_k$ is distinct from the next queried state $S_{k+1}$. In the Markovian regime, $(S_k)_{k\ge0}$ is an irreducible state trajectory with transition matrix $Q$, stationary distribution $\mu_{\X}$, and $S'_k=S_{k+1}$. Irreducibility on the finite state space guarantees a unique stationary distribution with full support. Periodicity is allowed. Temporal dependence enters our bounds through the Poisson condition number $A_Q$ defined in Appendix~\ref{app:common}.

For both CTD and MTD, we work with a block-supremum contraction metric $\ell_\infty$, which admits a statewise isometric embedding $I$ to a product space with contraction norm $\norm{\cdot}_{2,\infty}$, and a statewise projection $\Pi^\Theta$ onto the space of $\Theta$-supported, state-indexed representations $\mathcal F^\X_\Theta$. We can then compose the distributional Bellman operator $T^\pi$ to obtain the embedded, projected operator
\begin{equation}
\mathcal O := I \circ \Pi^\Theta T^\pi \circ I^{-1}.
\end{equation}
This operator admits a one-step sampled Bellman target $\widehat T(U_k;s,(R_k, S'_k))$ computed from a current estimate $U_k$ and a random sample $(R_k, S'_k)$, satisfying
\begin{equation}
\E \left[ \widehat T(U_k; S_k, (R_k, S'_k)) \mid U_k, S_k =s\right] = (\mathcal O U_k)(s),
\end{equation}
and a recursion that takes the form
\begin{equation}\label{eq:recursion}
U_{k+1} = U_k + \alpha_k P_{S_k} (\widehat T(U_k; S_k, (R_k, S'_k)) - U_k(S_k)),
\end{equation}
where for each $s \in \X$, $P_s$ denotes the coordinate projector onto block $s$.

Our discounted analysis expresses each categorical update as an asynchronous contractive SA step. The proofs use contraction of the stationary averaged operator in $\lVert \cdot \rVert_{2,\infty}$, a Moreau-envelope smoothing of $\lVert \cdot \rVert_{2,\infty}$ based on $\lVert \cdot \rVert_{2,p}$ \citep{borkar2008stochastic,kushner2003stochastic,beck2017first,moreau1965proximite,bauschke2017convex}, and a uniform or affine perturbation bound depending on the method. For Markovian trajectory sampling, a Poisson equation separates the stationary drift from the temporal dependence without introducing a mixing-time window \citep{haque2024stochastic}. More precisely, with $p^\star := \max\{2,\lceil \log \lvert \X \rvert\rceil\}$, the smoothing argument introduces a parameter $\vartheta>0$ through the generalized Moreau envelope $M_{\vartheta,p^\star}$ of the squared block-supremum distance. More details about these common ingredients are deferred to Appendix~\ref{app:common}, while the formal verification of the discounted finite-iteration results for CTD and MTD is deferred to Appendices~\ref{app:ctd} and~\ref{app:mtd}.

\subsection{Discounted CTD}\label{sec:ctd-new2}
For each state $s\in\X$, fix an ordered support $\Theta(s)=\{\theta_1(s)<\cdots<\theta_d(s)\}\subset\mathbb R$. $\mathcal F^{\X}_{\mathrm C,\Theta}$ is the class of state-indexed categorical laws supported on these statewise grids. The contraction metric $\ell_{\mathrm C, \infty}$ is the supremum Cram\'er metric, the embedding $I_{\mathrm C}$ is the standard cumulative-mass isometry $I_{\mathrm C, s}$ \citep{rowland2018analysis, bellemare2023distributional} applied to all states, and the statewise projection $\Pi_{\mathrm C}^{\Theta}$ is the usual linear-interpolation categorical projection $\Pi_{\mathrm C}^{\Theta(s)}$ applied to all states.

Given a sampled transition $(S_k,A_k,R_k,S'_k)$, the sampled Bellman target is
\begin{equation}
\widehat T_{\mathrm C}(U_k;S_k,(R_k,S'_k)) :=I_{\mathrm C,S_k}\Bigl(\Pi_{\mathrm C}^{\Theta(S_k)}\bigl((f_{R_k,\gamma})_\# I_{\mathrm C, S'_k}^{-1}(U_k(S'_k))\bigr)\Bigr),
\end{equation}
where $f_{r,\gamma}(z)=r+\gamma z$ and $\#$ is the measure-pushforward operator. The next theorem states the discounted CTD convergence rates for various step size regimes. Appendix~\ref{app:ctd} verifies the contraction, Lipschitz continuity and perturbation ingredients and records the explicit constants.

\begin{theorem}[Discounted CTD]\label{thm:ctd-main2}
Let $\eta^\star \in \mathcal F_{\mathrm C,\Theta}^{\X}$ denote the unique fixed point of $\Pi_{\mathrm C}^{\Theta} T^\pi$. From any $\eta_0\in\mathcal F_{\mathrm C,\Theta}^{\X}$, let $U_0=I_{\mathrm C}(\eta_0)$ and let $\eta_k := I_{\mathrm C}^{-1}(U_k)$ be generated by the asynchronous recursion \eqref{eq:recursion} with $\widehat T = \widehat T_{\mathrm C}$.

\textup{(i) i.i.d.\ sampling.} Suppose the sampled tuples $(S_k,A_k,R_k,S'_k)_{k\ge 0}$ are i.i.d.\ with query-state marginal $\rho$ on $\X$, and $\min_{s \in \X}\rho(s) > 0$. There exist explicit constants $C^{\mathrm{iid}}_\mathrm C, c^{\mathrm{iid}}_\mathrm C, \bar \alpha^{\mathrm{iid}}_\mathrm C >0$ and explicit thresholds $h^{\mathrm{iid}}_\mathrm C(\alpha)$ and $h^{\mathrm{iid}}_\mathrm C(\alpha, z)$ for which the following results hold on three step size regimes.

\textup{Constant step size.} If $\alpha_k \equiv \alpha \le \bar \alpha ^{\mathrm{iid}}_\mathrm C$, the iterates converge geometrically to an $O(\alpha)$ neighborhood:
\begin{equation}
\E\bigl[\ell_{\mathrm C, \infty}(\eta_k, \eta^\star)^2 \bigr] \le C^{\mathrm{iid}}_\mathrm C \ell_{\mathrm C, \infty}(\eta_0, \eta^\star)^2(1- c^{\mathrm{iid}}_\mathrm C\alpha)^k + C^{\mathrm{iid}}_\mathrm C \alpha.
\end{equation}
\textup{Linearly-diminishing step size.} If $\alpha_k = \alpha / (k+h)$, $\alpha > 1/c^{\mathrm{iid}}_\mathrm C$, and $h \ge h^{\mathrm{iid}}_\mathrm C(\alpha)$, then the leading residual term decays in $O(1/k)$:
\begin{equation}
\E\bigl[\ell_{\mathrm C, \infty}(\eta_k, \eta^\star)^2 \bigr] \le C^{\mathrm{iid}}_\mathrm C \ell_{\mathrm C, \infty}(\eta_0, \eta^\star)^2\left(\frac{h}{k+h} \right)^{c^{\mathrm{iid}}_\mathrm C \alpha} + \frac{C^{\mathrm{iid}}_\mathrm C \alpha^2}{c^{\mathrm{iid}}_\mathrm C\alpha -1} \cdot \frac{1}{k+h}.
\end{equation}
\textup{Polynomially-diminishing step size.} If $\alpha_k = \alpha/(k+h)^z$, $z \in (0,1)$, and $h \ge h^{\mathrm{iid}}_\mathrm C(\alpha, z)$, then the leading residual term decays in $O(1/k^z)$:
\begin{equation}
\E\bigl[\ell_{\mathrm C, \infty}(\eta_k, \eta^\star)^2 \bigr] \le C^{\mathrm{iid}}_\mathrm C \ell_{\mathrm C, \infty}(\eta_0, \eta^\star)^2\exp\left(-\frac{c^{\mathrm{iid}}_\mathrm C \alpha}{1-z} \bigl((k+h)^{1-z} - h^{1-z} \bigr) \right) + \frac{C^{\mathrm{iid}}_\mathrm C \alpha}{(k+h)^z}.
\end{equation}

\textup{(ii) Markovian trajectory sampling.} Suppose $(S_k)_{k\ge 0}$ is an irreducible chain with stationary distribution $\mu_\X$. There exist explicit constants $C^{\mathrm{mk}}_\mathrm C, c^{\mathrm{mk}}_\mathrm C, \bar \alpha^{\mathrm{mk}}_\mathrm C >0$ and explicit thresholds $h^{\mathrm{mk}}_\mathrm C(\alpha)$ and $h^{\mathrm{mk}}_\mathrm C(\alpha ,z)$ for which the following results hold for every $k\ge1$.

\textup{Constant step size.} If $\alpha_k \equiv \alpha$ with $0<\alpha \le \bar \alpha ^{\mathrm{mk}}_\mathrm C$, then the iterates converge geometrically to an $O(\alpha)$ neighborhood:
\begin{equation}
\E\bigl[\ell_{\mathrm C, \infty}(\eta_k, \eta^\star)^2 \bigr] \le C^{\mathrm{mk}}_\mathrm C \bigl(1+\ell_{\mathrm C, \infty}(\eta_0, \eta^\star)^2\bigr)\exp\left(-c^{\mathrm{mk}}_\mathrm C\alpha k\right) + C^{\mathrm{mk}}_\mathrm C \alpha.
\end{equation}
\textup{Linearly-diminishing step size.} If $\alpha_k = \alpha / (k+h)$, $\alpha >1 /c^{\mathrm{mk}}_\mathrm C$, and $h \ge h^{\mathrm{mk}}_\mathrm C(\alpha)$, then the leading residual term decays in $O(1/k)$:
\begin{equation}
\E\bigl[\ell_{\mathrm C, \infty}(\eta_k, \eta^\star)^2 \bigr] \le C^{\mathrm{mk}}_\mathrm C \bigl(1+\ell_{\mathrm C, \infty}(\eta_0, \eta^\star)^2\bigr)\left(\frac{h}{k+h}\right)^{c^{\mathrm{mk}}_\mathrm C \alpha} + \frac{C^{\mathrm{mk}}_\mathrm C \alpha^2}{c^{\mathrm{mk}}_\mathrm C \alpha -1} \cdot \frac{1}{k+h}.
\end{equation}
\textup{Polynomially-diminishing step size.} If $\alpha_k = \alpha / (k+h)^z$, $\alpha>0$, $z \in (0, 1)$, and $h \ge h^{\mathrm{mk}}_\mathrm C(\alpha, z)$, define $\chi_{\mathrm C}:=z/\alpha+2c^{\mathrm{mk}}_\mathrm C$. Then the leading residual term decays in $O(1/k^z)$:
\begin{equation}
\resizebox{\linewidth}{!}{$\displaystyle \E\bigl[\ell_{\mathrm C, \infty}(\eta_k, \eta^\star)^2 \bigr] \le C^{\mathrm{mk}}_\mathrm C \bigl(1+\ell_{\mathrm C, \infty}(\eta_0, \eta^\star)^2\bigr)\exp\left(-\frac{c^{\mathrm{mk}}_\mathrm C \alpha}{1-z} \bigl((k+h)^{1-z} - h^{1-z} \bigr) \right) + \frac{C^{\mathrm{mk}}_\mathrm C \alpha(1+\chi_{\mathrm C})}{(k+h)^z}.$}
\end{equation}
\end{theorem}
\begin{proof}[Proof sketch]
The Cram\'er embedding makes the projected Bellman map a $\sqrt{\gamma}$-contraction. Under i.i.d.\ sampling, the averaged asynchronous map contracts according to the minimum state mass. Under Markovian trajectory sampling, a Poisson equation decomposes the state-dependent drift into its stationary mean and a telescoping correction. The sampled target is Lipschitz and its centered perturbation is uniformly bounded by the support radius. A sufficiently small smoothing parameter makes the Moreau-envelope drift positive, while the step size conditions control its quadratic remainder. Appendix~\ref{app:ctd} gives the exact constants and thresholds.
\end{proof}
The following finite-sample corollary follows from the linearly-diminishing step size bounds.

\begin{corollary}\label{cor:disc-ctd}
To guarantee
$\mathbb{E}[\ell_{\mathrm C,\infty}(\eta_k,\eta^\star)] \le \varepsilon$ with the discounted CTD recursion, it suffices to take $k = O(\varepsilon^{-2})$ samples under either i.i.d.\ or Markovian sampling.
\end{corollary}

\subsection{Discounted MTD}\label{sec:mtd-new2}
Let $q \ge 2$. For each state $s\in\X$, fix a $d$-point support $\Theta(s)=\{\theta_1(s),\cdots,\theta_d(s)\}\subset\mathbb R^q$. $\mathcal F^{\X}_{\mathrm M,\Theta}$ is the class of state-indexed, mass-$1$ signed categorical laws supported on these points. The contraction metric $\ell_{\mathrm M, \infty}$ is the supremum MMD metric with a continuous characteristic kernel $\kappa$ induced by a shift-invariant, $c$-homogeneous semimetric of strong negative type \citep{smola2007hilbert,gretton2012kernel,muandet2017kernel,wiltzer2024foundations}.

The embedding $I_{\mathrm M}$ is the per-state embedding $I_{\mathrm M, s}(\eta(s)):=K_s^{1/2}p(s)$ applied to all states, $K_s:=\bigl(\kappa(\theta_i(s),\theta_j(s))\bigr)_{i,j=1}^d$ is the Gram matrix on the support $\Theta(s)$ \citep{wiltzer2024foundations}. The statewise projection $\Pi_{\mathrm M}^{\Theta}$ is
\begin{equation}
\Pi^{\Theta(s)}_{\mathrm M}\nu := \arg\inf_{p\in\mathbb R_1^d} \mathrm{MMD}_{\kappa}\Bigl(\sum_{i=1}^d p_i\delta_{\theta_i(s)},\nu\Bigr),
\end{equation}
applied to all states, where $\mathbb R_1^d$ is the affine subspace of $\mathbb{R}^d$ with unit total mass. The projection is well-defined for $\mathcal F^\X_{\mathrm M,\Theta}$, as shown by \citet{wiltzer2024foundations}.

Given a sampled transition $(S_k,A_k,R_k,S'_k)$, the sampled Bellman target is
\begin{equation}
\widehat T_{\mathrm M}(U_k;S_k,(R_k,S'_k)) := I_{\mathrm M, S_k}\Bigl(\Pi_{\mathrm M}^{\Theta(S_k)}\bigl((f_{R_k,\gamma})_\# I_{\mathrm M, S'_k}^{-1}(U_k(S'_k))\bigr)\Bigr).
\end{equation}

The next theorem gives the MTD rates analogous to Theorem~\ref{thm:ctd-main2}. The proof backbone is the same as for CTD, but the perturbation geometry is affine in the centered error.

\begin{theorem}[Discounted MTD]\label{thm:mtd-main}
Let $\eta^\star \in \mathcal F_{\mathrm M,\Theta}^{\X}$ denote the unique fixed point of $\Pi_{\mathrm M}^{\Theta} T^\pi$. From any $\eta_0\in\mathcal F_{\mathrm M,\Theta}^{\X}$, let $U_0=I_{\mathrm M}(\eta_0)$ and let $\eta_k := I_{\mathrm M}^{-1}(U_k)$ be generated by \eqref{eq:recursion} with $\widehat T=\widehat T_{\mathrm M}$.

\textup{(i) i.i.d.\ sampling.} Suppose the sampled tuples $(S_k,A_k,R_k,S'_k)_{k\ge 0}$ are i.i.d.\ with query-state marginal $\rho$ on $\X$, and $\min_{s \in \X}\rho(s) > 0$. There exist explicit constants $C^{\mathrm{iid}}_\mathrm M, c^{\mathrm{iid}}_\mathrm M, \bar \alpha^{\mathrm{iid}}_\mathrm M >0$ and explicit thresholds $h^{\mathrm{iid}}_\mathrm M(\alpha)$ and $h^{\mathrm{iid}}_\mathrm M(\alpha, z)$ for which the following results hold on three step size regimes.

\textup{Constant step size.} If $\alpha_k \equiv \alpha \le \bar \alpha ^{\mathrm{iid}}_\mathrm M$, the iterates converge geometrically to an $O(\alpha)$ neighborhood:
\begin{equation}
\E\bigl[\ell_{\mathrm M, \infty}(\eta_k, \eta^\star)^2 \bigr] \le C^{\mathrm{iid}}_\mathrm M \ell_{\mathrm M, \infty}(\eta_0, \eta^\star)^2(1- c^{\mathrm{iid}}_\mathrm M\alpha)^k + C^{\mathrm{iid}}_\mathrm M \alpha.
\end{equation}
\textup{Linearly-diminishing step size.} If $\alpha_k = \alpha / (k+h)$, $\alpha > 1/c^{\mathrm{iid}}_\mathrm M$, and $h \ge h^{\mathrm{iid}}_\mathrm M(\alpha)$, then the leading residual term decays in $O(1/k)$:
\begin{equation}
\E\bigl[\ell_{\mathrm M, \infty}(\eta_k, \eta^\star)^2 \bigr] \le C^{\mathrm{iid}}_\mathrm M \ell_{\mathrm M, \infty}(\eta_0, \eta^\star)^2\left(\frac{h}{k+h} \right)^{c^{\mathrm{iid}}_\mathrm M \alpha} + \frac{C^{\mathrm{iid}}_\mathrm M \alpha^2}{c^{\mathrm{iid}}_\mathrm M\alpha -1} \cdot \frac{1}{k+h}.
\end{equation}
\textup{Polynomially-diminishing step size.} If $\alpha_k = \alpha/(k+h)^z$, $z \in (0,1)$, and $h \ge h^{\mathrm{iid}}_\mathrm M(\alpha, z)$, then the leading residual term decays in $O(1/k^z)$:
\begin{equation}
\E\bigl[\ell_{\mathrm M, \infty}(\eta_k, \eta^\star)^2 \bigr] \le C^{\mathrm{iid}}_\mathrm M \ell_{\mathrm M, \infty}(\eta_0, \eta^\star)^2\exp\left(-\frac{c^{\mathrm{iid}}_\mathrm M \alpha}{1-z} \bigl((k+h)^{1-z} - h^{1-z} \bigr) \right) + \frac{C^{\mathrm{iid}}_\mathrm M \alpha}{(k+h)^z}.
\end{equation}

\textup{(ii) Markovian trajectory sampling.} Suppose $(S_k)_{k\ge 0}$ is an irreducible chain with stationary distribution $\mu_\X$. There exist explicit constants $C^{\mathrm{mk}}_\mathrm M, c^{\mathrm{mk}}_\mathrm M, \bar \alpha^{\mathrm{mk}}_\mathrm M >0$ and explicit thresholds $h^{\mathrm{mk}}_\mathrm M(\alpha)$ and $h^{\mathrm{mk}}_\mathrm M(\alpha ,z)$ for which the following results hold for every $k\ge1$.

\textup{Constant step size.} If $\alpha_k \equiv \alpha$ with $0<\alpha \le \bar \alpha ^{\mathrm{mk}}_\mathrm M$, then the iterates converge geometrically to an $O(\alpha)$ neighborhood:
\begin{equation}
\E\bigl[\ell_{\mathrm M, \infty}(\eta_k, \eta^\star)^2 \bigr] \le C^{\mathrm{mk}}_\mathrm M \bigl(1+\ell_{\mathrm M, \infty}(\eta_0, \eta^\star)^2\bigr)\exp\left(-c^{\mathrm{mk}}_\mathrm M\alpha k\right) + C^{\mathrm{mk}}_\mathrm M \alpha.
\end{equation}
\textup{Linearly-diminishing step size.} If $\alpha_k = \alpha / (k+h)$, $\alpha >1 /c^{\mathrm{mk}}_\mathrm M$, and $h \ge h^{\mathrm{mk}}_\mathrm M(\alpha)$, then the leading residual term decays in $O(1/k)$:
\begin{equation}
\E\bigl[\ell_{\mathrm M, \infty}(\eta_k, \eta^\star)^2 \bigr] \le C^{\mathrm{mk}}_\mathrm M \bigl(1+\ell_{\mathrm M, \infty}(\eta_0, \eta^\star)^2\bigr)\left(\frac{h}{k+h}\right)^{c^{\mathrm{mk}}_\mathrm M \alpha} + \frac{C^{\mathrm{mk}}_\mathrm M \alpha^2}{c^{\mathrm{mk}}_\mathrm M \alpha -1} \cdot \frac{1}{k+h}.
\end{equation}
\textup{Polynomially-diminishing step size.} If $\alpha_k = \alpha / (k+h)^z$, $\alpha>0$, $z \in (0, 1)$, and $h \ge h^{\mathrm{mk}}_\mathrm M(\alpha, z)$, define $\chi_{\mathrm M}:=z/\alpha+2c^{\mathrm{mk}}_\mathrm M$. Then the leading residual term decays in $O(1/k^z)$:
\begin{equation}
\resizebox{\linewidth}{!}{$\displaystyle \E\bigl[\ell_{\mathrm M, \infty}(\eta_k, \eta^\star)^2 \bigr] \le C^{\mathrm{mk}}_\mathrm M \bigl(1+\ell_{\mathrm M, \infty}(\eta_0, \eta^\star)^2\bigr)\exp\left(-\frac{c^{\mathrm{mk}}_\mathrm M \alpha}{1-z} \bigl((k+h)^{1-z} - h^{1-z} \bigr) \right) + \frac{C^{\mathrm{mk}}_\mathrm M \alpha(1+\chi_{\mathrm M})}{(k+h)^z}.$}
\end{equation}
\end{theorem}
\begin{proof}[Proof sketch]
The Gram-matrix embedding turns the projected MTD Bellman map into a $\gamma^{c/2}$-contraction in the block-supremum MMD geometry. The signed-categorical projection gives an affine centered perturbation bound unlike the uniform CTD bound. Under Markovian trajectory sampling, the Poisson decomposition used for CTD applies, with the affine coefficient entering the admissible step size. The smoothing parameter makes the Moreau-envelope drift positive, and the schedule thresholds control the quadratic remainder. Appendix~\ref{app:mtd} gives the exact constants and thresholds.
\end{proof}

Once again, the following finite-sample result is a simple consequence of the previous theorem.
\begin{corollary}\label{cor:disc-mtd}
To guarantee
$\mathbb{E}[\ell_{\mathrm M,\infty}(\eta_k,\eta^\star)] \le \varepsilon$ with the discounted MTD recursion, it suffices to take $k = O(\varepsilon^{-2})$ samples under either i.i.d.\ or Markovian sampling.
\end{corollary}

\subsection{Explicit discounted dependence}
The constants above can be read directly in natural problem parameters. Let $\beta_{\mathrm C}=\sqrt\gamma$ and $\beta_{\mathrm M}=\gamma^{c/2}$. For method $\star\in\{\mathrm C,\mathrm M\}$, set
\begin{equation}
\delta_{\star,\rho}:=\rho_{\min}(1-\beta_\star),\qquad \delta_{\star,\mu}:=\mu_{\min}(1-\beta_\star).
\end{equation}
With the explicit smoothing choice in Proposition~\ref{prop:explicit-discounted-constants}, the constant-step i.i.d.\ residual is
\begin{equation}
O\left(\frac{\log\lvert\X\rvert\,A_\star^{\mathrm{iid}}}{\delta_{\star,\rho}^2}\alpha\right),
\end{equation}
where $A_{\mathrm C}^{\mathrm{iid}}=4B_{\mathrm C}^2$ and $A_{\mathrm M}^{\mathrm{iid}}=8\beta_{\mathrm M}^2+2(B_{\mathrm M}^\star)^2$. The corresponding Markovian residual is
\begin{equation}
O\left(\frac{\log\lvert\X\rvert\,A_Q(B_{\star}^{\mathrm{mk}})^2}{\delta_{\star,\mu}^2}\alpha\right),
\end{equation}
where $B_{\mathrm C}^{\mathrm{mk}}=2B_{\mathrm C}$ and $B_{\mathrm M}^{\mathrm{mk}}=B_{\mathrm M}^\star$. The factor $A_Q$ measures trajectory dependence through the Poisson equation and replaces the usual logarithmic mixing window in Markovian analyses \citep{chen2022finite, chen2024lyapunov}. Choosing the linear-schedule coefficient as a fixed multiple of the inverse drift gap gives, after the explicit offset and initialization transient,
\begin{equation}
k_{\mathrm{iid}}=O\left(\frac{\log\lvert\X\rvert\,A_\star^{\mathrm{iid}}}{\delta_{\star,\rho}^3\varepsilon^2}\right),\qquad k_{\mathrm{mk}}=O\left(\frac{\log\lvert\X\rvert\,A_Q(B_\star^{\mathrm{mk}})^2}{\delta_{\star,\mu}^3\varepsilon^2}\right)
\end{equation}
asynchronous updates for expected metric error at most $\varepsilon$. The support size $d$ and reward dimension $q$ do not enter through the block-supremum smoothing step. Their influence is localized in $B_{\mathrm C}$ or in the MTD support-geometry constant $B_{\mathrm M}^\star$. For CTD with uniform queries and a standard return support of width $O((1-\gamma)^{-1})$, the conservative i.i.d.\ bound scales as $\tilde O(\lvert\X\rvert^3(1-\gamma)^{-4}\varepsilon^{-2})$ updates.

Table~\ref{tab:discounted-comparison} compares the closest finite-sample results for tabular scalar distributional policy evaluation. Offline evaluation, function approximation, and asymptotic analyses are not included in the table because they are not directly comparable. For compactness, we write $\mathfrak D_{\mathrm C}^{\mathrm{iid}}:=\log\lvert\X\rvert B_{\mathrm C}^2\delta_{\mathrm C,\rho}^{-3}$ and $\mathfrak D_{\mathrm C}^{\mathrm{mk}}:=\log\lvert\X\rvert A_QB_{\mathrm C}^2\delta_{\mathrm C,\mu}^{-3}$.
\begin{table}[t]
\centering
\caption{Discounted tabular scalar distributional policy-evaluation guarantees. The final column counts individual environment transitions after converting synchronous sweeps and per-state sample counts. The $\tilde O$ rates suppress both logarithmic factors and method-specific support conditions.}
\label{tab:discounted-comparison}
\scriptsize
\begin{tabularx}{\linewidth}{@{}>{\raggedright\arraybackslash}p{0.16\linewidth}>{\raggedright\arraybackslash}p{0.19\linewidth}>{\raggedright\arraybackslash}X>{\raggedright\arraybackslash}p{0.31\linewidth}@{}}
\toprule
Work & Method and sampling & Error criterion & Leading transition complexity \\
\midrule
\citet{boeck2022speedy} & accelerated SCPE, synchronous generative model & Cram\'er to projected fixed point, high probability & $\tilde O(\lvert\X\rvert\lvert\A\rvert\varepsilon^{-2}(1-\gamma)^{-3})$ \\
\citet{zhang2023estimation} & plug-in estimator, generative model & supremum $1$-Wasserstein to true return law, high probability & $\tilde O(\lvert\X\rvert\lvert\A\rvert\varepsilon^{-2}(1-\gamma)^{-4})$ \\
\citet{rowland2024near} & DCFP, generative model & supremum $1$-Wasserstein to true return law, high probability & $\tilde O(\lvert\X\rvert\varepsilon^{-2}(1-\gamma)^{-3})$ \\
\citet{peng2025statisticalefficiencydistributionaltemporal} & CTD, one-state i.i.d.\ queries & Cram\'er to true return law, high probability & $\tilde O(\rho_{\min}^{-1}\varepsilon^{-2}(1-\gamma)^{-5/2})$ \\
\citet{peng2025statisticalefficiencydistributionaltemporal} & VR-CTD, one Markovian trajectory & supremum $1$-Wasserstein to true return law, high probability & \makecell[l]{$\tilde O\!\left(\mu_{\X,\min}^{-1}\varepsilon^{-2}(1-\gamma)^{-3}\right)$\\${}+\tilde O\left(\mu_{\X,\min}^{-1}t_{\mathrm{mix}}(1-\gamma)^{-1}\right)$} \\
This work$^\dagger$ & CTD, one-state i.i.d.\ queries & Cram\'er to projected fixed point, expectation & $O(\mathfrak D_{\mathrm C}^{\mathrm{iid}}\varepsilon^{-2})$ one-state updates \\
This work$^\dagger$ & CTD, one Markovian trajectory & Cram\'er to projected fixed point, expectation & $O(\mathfrak D_{\mathrm C}^{\mathrm{mk}}\varepsilon^{-2})$ transitions \\
\bottomrule
\multicolumn{4}{@{}l@{}}{$^\dagger$The paper also proves analogous guarantees for multivariate MTD. The displayed rows report the scalar CTD specialization.}
\end{tabularx}
\end{table}

\section{Undiscounted fixed-horizon categorical policy evaluation}\label{sec:undisc}

We now turn to the undiscounted fixed-horizon setting, where the value object is indexed by the remaining horizon \citep{sutton1998reinforcement,azar2017minimax,de2020fixed}. The aim of this section is to recover a comparable finite-iteration picture without relying on discounting, which means the learned object must now be treated as a horizon-indexed stack rather than a single per-state quantity. Fix a horizon $H \in \mathbb N$ and a stationary policy $\pi(\cdot \mid s)$. For each $h \in \{0,1,\dots,H\}$ and each $s \in \X$, let $\eta^h(s)$ denote the $h$-step return-distribution object under $\pi$, with terminal layer $\eta^0(s)=\delta_0$.

For $h \ge 1$, we define the undiscounted fixed-horizon distributional Bellman operator by
\begin{equation}
(T_H^\pi \eta)^h(s) := \sum_{a\in \A}\pi(a \mid s)\sum_{s'\in \X} P(s' \mid s,a)\, (f_{R(s,a),1})_\# \eta^{h-1}(s').
\end{equation}
Thus the $h$th horizon bootstraps from the $(h-1)$st horizon, as in the fixed-horizon TD formulation of \citet{de2020fixed}.

For both CTD and MTD, we use exactly the same local categorical projections and local statewise isometric embeddings as in Section~\ref{sec:discounted}, now applied separately at each horizon and then stacked over $h=1,\dots,H$. For each $s \in \X$, the corresponding embedded iterate is
\begin{equation}
U(s) := \bigl(U^1(s), \ldots, U^H(s)\bigr), \qquad U := \bigl(U^h(s)\bigr)_{1 \le h \le H,\ s \in \X}.
\end{equation}

The contraction of $\mathcal O$ in the discounted setting is entirely due to the discount factor $\gamma \in (0,1)$ and hence does not carry over to the undiscounted case. However, after introducing a suitable weighted block metric, we recover a contraction. Set $ \lambda:=H/(H+1)$ and define the weighted block metric and weighted embedded norm by
\begin{equation}
\ell_{H,\infty}(\eta,\eta') := \max_{1\le h\le H,\ s\in \X} \lambda^h \ell\bigl(\eta^h(s),{\eta'}^h(s)\bigr), \qquad \lVert U \rVert_{H,2,\infty} := \max_{1\le h\le H,\ s\in \X} \lambda^h \lVert U^h(s) \rVert_2.
\end{equation}
The natural unweighted error is
\begin{equation}
\ell_{H,\infty}^{\mathrm{unw}}(\eta,\eta'):=\max_{1\le h\le H,\ s\in\X}\ell\bigl(\eta^h(s),{\eta'}^h(s)\bigr).
\end{equation}
Since $(1+1/H)^H\le e$, this choice gives
\begin{equation}
\ell_{H,\infty}(\eta,\eta')\le\ell_{H,\infty}^{\mathrm{unw}}(\eta,\eta')\le\lambda^{-H}\ell_{H,\infty}(\eta,\eta')\le e\,\ell_{H,\infty}(\eta,\eta').
\end{equation}
Let $\Pi_H^\Theta$ and $I_H$ denote the horizonwise projection and horizonwise embedding, and define
\begin{equation}
\mathcal O_H:=I_H\circ \Pi_H^\Theta T_H^\pi \circ I_H^{-1}.
\end{equation}

\begin{proposition}\label{prop:undisc-contract}
For both CTD and MTD,
\begin{equation}
\ell_{H,\infty}(\Pi_H^\Theta T_H^\pi \eta,\Pi_H^\Theta T_H^\pi \eta') \le \lambda \ell_{H,\infty}(\eta,\eta').
\end{equation}
By construction, $I_H$ is an isometric embedding from the fixed-horizon distribution space equipped with $\ell_{H,\infty}$ into the embedding space equipped with $\lVert \cdot \rVert_{H,2,\infty}$. Hence the preceding statement is equivalent to
\begin{equation}
\lVert \mathcal O_H U-\mathcal O_H U'\rVert_{H,2,\infty} \le \lambda \lVert U-U'\rVert_{H,2,\infty}.
\end{equation}
\end{proposition}
\begin{proof}[Proof sketch] 
The Bellman update moves horizon $h$ to $h-1$, and the weight $\lambda^h$ therefore contributes exactly one factor $\lambda$. The local categorical projections are nonexpansive in the underlying statewise metric, so the weighted stack inherits a contraction. Details are recorded in Appendices~\ref{app:fh-ctd} and~\ref{app:fh-mtd}.
\end{proof}

For each sampled transition $(S_k,A_k,R_k,S_{k+1})$, let
\begin{equation}
\widehat T_H\bigl(U_k; S_k, (R_k,S_{k+1})\bigr) := \bigl(\widehat T_H^1\bigl(U_k; S_k, (R_k,S_{k+1})\bigr), \ldots, \widehat T_H^H\bigl(U_k; S_k, (R_k,S_{k+1})\bigr)\bigr)
\end{equation}
denote the stacked one-step sampled Bellman target, where the horizonwise components are defined in the CTD and MTD subsections below. Then, for every $s \in \X$,
\begin{equation}
\mathbb{E}\bigl[\widehat T_H\bigl(U_k; S_k, (R_k,S_{k+1})\bigr)\mid U_k,\ S_k = s \bigr] = (\mathcal O_H U_k)(s).
\end{equation}
Thus a single sampled transition supplies a Bellman sample for the full stack at the visited state: the observed first reward is shared across all horizons, while the continuation term for horizon $h$ is the current $(h-1)$-step estimate at the next state. If $P_s^H$ denotes the coordinate projector onto the full horizon stack at state $s$, the online fixed-horizon recursion takes the form
\begin{equation}\label{eq:undisc-recur}
U_{k+1} = U_k + \alpha_k P_{S_k}^H \Bigl(\widehat T_H\bigl(U_k; S_k, (R_k,S_{k+1})\bigr) - U_k(S_k) \Bigr).
\end{equation}

We consider only the standard episodic sampling model in which episodes have fixed horizon $H$, reset distribution $\nu_0$, are i.i.d.\ across resets, and follow the stationary policy $\pi$ within each episode. The recursion in \eqref{eq:undisc-recur} is still performed after each transition, but its averaged drift depends on the within-episode phase. We therefore analyze the episode-boundary sequence $(U_{mH})_{m \ge 0}$. To formalize this phase dependence, define the phase distributions
\begin{equation}
\rho_t(s) := \Pr(S_t = s), \qquad 0 \le t \le H-1, \qquad s \in \X,
\end{equation}
and the average episode occupancy
\begin{equation}
\bar\rho_H(s):=\frac{1}{H}\sum_{t=0}^{H-1}\rho_t(s),\qquad \rho_{H,\min}:=\min_{s\in\X}\bar\rho_H(s)>0.
\end{equation}
The cumulative occupancy condition permits the phase distributions to have different supports while ensuring that every state is visited with positive probability over a complete episode. For each phase $t \in \{0,\ldots,H-1\}$, define the phasewise averaged increment and phasewise averaged map by
\begin{equation}
\Gamma_t(U) := \sum_{s \in \X} \rho_t(s)\, P_s^H \bigl((\mathcal O_H U)(s) - U(s) \bigr), \qquad G_t(U) := U + \Gamma_t(U).
\end{equation}
For each episode $m \ge 0$, define the episodewise first- and second-order step size masses
\begin{equation}
\bar\alpha_m := \sum_{u=0}^{H-1}\alpha_{mH+u}, \qquad \bar\alpha_m^{(2)} := \sum_{u=0}^{H-1}\alpha_{mH+u}^2.
\end{equation}
These are the episodewise analogues of the single-step quantities $\alpha_k$ and $\alpha_k^2$ in the discounted analysis. The phase dependence enters through the averaged maps $G_t$, while the weighted block-supremum geometry from Appendix~\ref{app:common} remains unchanged.
The maps $G_t$ are auxiliary proof objects for the episodewise drift analysis rather than additional algorithmic iterates.

\subsection{Undiscounted fixed-horizon CTD}

For CTD, each state-horizon pair $(h,s)$ carries an ordered, possibly horizon-dependent scalar support $\Theta_h(s)=\{\theta_{h,1}(s)<\cdots<\theta_{h,d}(s)\}\subset\mathbb R$. The statewise Cram\'er metric $\ell_{H, \mathrm C, \infty}$, cumulative-mass embedding $I_{H, \mathrm C, h, s}$, and linear-interpolation projection $\Pi^{\Theta_h(s)}_{H, \mathrm C}$ are defined exactly as in Section~\ref{sec:ctd-new2}, only now applied separately at each horizon and stacked over $h=1,\dots,H$ \citep{rowland2018analysis,bellemare2023distributional}. We write
\begin{equation}
B_{H,\mathrm C}:=\max_{1\le h\le H,\ s\in\X}\lambda^h\sqrt{\theta_{h,d}(s)-\theta_{h,1}(s)}
\end{equation}
for the weighted support radius. Given a sampled transition $(S_k,A_k,R_k,S_{k+1})$, the sampled Bellman target is computed at every horizon $h=1,\dots,H$ by
\begin{equation}
\widehat T_{H,\mathrm C}^h(U_k;S_k,(R_k,S_{k+1})) := I_{H,\mathrm C,h,S_k} \Bigl(\Pi_{H,\mathrm C}^{\Theta_h(S_k)}\bigl((f_{R_k,1})_\# I_{H,\mathrm C,h-1,S_{k+1}}^{-1}(U_k^{h-1}(S_{k+1}))\bigr)\Bigr),
\end{equation}
with the terminal convention $I_{H,\mathrm C,0,s}^{-1}(U_k^0(s)):=\delta_0$ for every $s$ and $k$. The learned horizon stack therefore consists of layers $1,\dots,H$, with $\delta_0$ providing the fixed base case. Stacking these horizonwise targets gives
\begin{equation}
\widehat T_{H,\mathrm C}(U_k;S_k,(R_k,S_{k+1})) := \bigl(\widehat T_{H,\mathrm C}^1(U_k;S_k,(R_k,S_{k+1})), \dots, \widehat T_{H,\mathrm C}^H(U_k;S_k,(R_k,S_{k+1}))\bigr).
\end{equation}
One sampled transition therefore updates the full block at $S_k$: all horizons share the same first step and differ only in the bootstrapped tail. The next theorem states the resulting episode-boundary rates.

\begin{theorem}[Undiscounted fixed-horizon episodic CTD]\label{thm:undisc-ctd}
Assume the episodic sampling model described above and $\rho_{H,\min}>0$. Let $\eta_{H,\mathrm C}^\star$ denote the unique fixed point of $\Pi_{H,\mathrm C}^{\Theta} T_H^\pi$. From any $\eta_0\in\mathcal F_{H,\mathrm C,\Theta}^{\X}$, let $U_0=I_{H,\mathrm C}(\eta_0)$ and let $\eta_k := I_{H,\mathrm C}^{-1}(U_k)$ be generated by \eqref{eq:undisc-recur} with $\widehat T_{H} = \widehat T_{H,\mathrm C}$. There exist explicit constants $C^H_\mathrm C, c^H_\mathrm C, \bar \alpha ^H_\mathrm C > 0$ and explicit thresholds $g^H_\mathrm C(\alpha)$ and $g^H_\mathrm C(\alpha, z)$ for which the following episode-boundary results hold in the natural unweighted error.

\textup{Constant step size.} If $\alpha_k \equiv \alpha \le \bar \alpha^H_\mathrm C$, then the episode-boundary error decays geometrically to an $O(\alpha)$ neighborhood:
\begin{equation}
\E\bigl[\ell_{H, \mathrm C, \infty}^{\mathrm{unw}}(\eta_{mH}, \eta_{H,\mathrm C}^\star)^2 \bigr] \le C^{H}_\mathrm C \ell_{H,\mathrm C, \infty}^{\mathrm{unw}}(\eta_0, \eta_{H,\mathrm C}^\star)^2(1- c^{H}_\mathrm C H\alpha)^m + C^{H}_\mathrm C \alpha.
\end{equation}
For the diminishing step size regimes below, write $\tau_m := mH+g+H-1$, so $\tau_0 = g+H-1$, where $g$ is the step-size offset.

\textup{Linearly-diminishing step size.} If $\alpha_k = \alpha / (k+g)$, $\alpha > 1/c^{H}_\mathrm C$, and $g \ge \max\{H-1,g^{H}_\mathrm C(\alpha)\}$, then the leading episode-boundary residual term decays in $O(1/m)$:
\begin{equation}
\E\bigl[\ell_{H,\mathrm C, \infty}^{\mathrm{unw}}(\eta_{mH}, \eta_{H,\mathrm C}^\star)^2 \bigr] \le C^{H}_\mathrm C \ell_{H, \mathrm C, \infty}^{\mathrm{unw}}(\eta_0, \eta_{H,\mathrm C}^\star)^2\left(\frac{\tau_0}{\tau_m} \right)^{c^{H}_\mathrm C\alpha} + \frac{C^{H}_\mathrm C \alpha^2}{c^{H}_\mathrm C\alpha -1} \cdot \frac{1}{\tau_m}.
\end{equation}
\textup{Polynomially-diminishing step size.} If $\alpha_k = \alpha/(k+g)^z$, $z \in (0,1)$, and $g \ge \max\{H-1,g^{H}_\mathrm C(\alpha, z)\}$, then the leading episode-boundary residual term decays in $O(1/m^z)$:
\begin{equation}
\E\bigl[\ell_{H, \mathrm C, \infty}^{\mathrm{unw}}(\eta_{mH}, \eta_{H,\mathrm C}^\star)^2 \bigr] \le C^{H}_\mathrm C \ell_{H, \mathrm C, \infty}^{\mathrm{unw}}(\eta_0, \eta_{H,\mathrm C}^\star)^2\exp\left(-\frac{c^{H}_\mathrm C \alpha}{1-z} \bigl(\tau^{1-z}_m - \tau^{1-z}_0 \bigr) \right) + \frac{C^{H}_\mathrm C \alpha}{\tau_m^z}.
\end{equation}
\end{theorem}
\begin{proof}[Proof sketch]
Weighting horizon $h$ by $\lambda^h$ turns the horizon-to-horizon bootstrap into a contraction, and $\lambda=H/(H+1)$ converts the resulting bound to the unweighted metric with a universal factor. The proof freezes the iterate at the episode boundary and combines all phasewise means into one averaged episode map. The contraction in this case is governed by the average occupancy $\rho_{H,\min}$, while the difference from the online trajectory is controlled by a mean-zero frozen term and within-episode movement. Details are given in Appendix~\ref{app:fh-ctd}.
\end{proof}

\begin{corollary}\label{cor:undisc-ctd}
To guarantee $\mathbb{E}[\ell_{H, \mathrm C,\infty}^{\mathrm{unw}}(\eta_{mH},\eta_{H, \mathrm C}^\star)] \le \varepsilon$ with the undiscounted fixed-horizon CTD recursion, it suffices to take $m = O(\varepsilon^{-2})$ episodes. The constant is polynomial in $H$, $1/\rho_{H,\min}$, $1+\log(H\lvert\X\rvert)$, and the support radius $B_{H,\mathrm C}$. Equivalently, the transition count is $k=mH$.
\end{corollary}

\subsection{Undiscounted fixed-horizon MTD}
For MTD, each state-horizon pair $(h,s)$ carries a possibly horizon-dependent $d$-point support $\Theta_h(s)=\{\theta_{h,1}(s),\dots,\theta_{h,d}(s)\}\subset\mathbb R^q$. The statewise MMD metric $\ell_{H, \mathrm M, \infty}$ with a characteristic kernel $\kappa$ induced by a shift-invariant, $c$-homogeneous semimetric of strong negative type, the signed-categorical projection $\Pi^{\Theta_h(s)}_{H, \mathrm M}$, and the Gram-matrix embedding $I_{H, \mathrm M, h, s}$ are defined exactly as in Section~\ref{sec:mtd-new2}, now applied separately at each horizon and stacked over $h=1,\dots,H$.

Given a sampled transition $(S_k,A_k,R_k,S_{k+1})$, the sampled Bellman target is computed at every horizon $h=1,\dots,H$ by
\begin{equation}
\widehat T_{H,\mathrm M}^h(U_k;S_k,(R_k,S_{k+1})) := I_{H,\mathrm M,h,S_k} \Bigl(\Pi_{H,\mathrm M}^{\Theta_h(S_k)}\bigl((f_{R_k,1})_\# I_{H,\mathrm M,h-1,S_{k+1}}^{-1}(U_k^{h-1}(S_{k+1}))\bigr)\Bigr),
\end{equation}
again with the terminal convention $I_{H,\mathrm M,0,s}^{-1}(U_k^0(s)):=\delta_0$ for every $s$ and $k$. Stacking these horizonwise targets gives
\begin{equation}
\widehat T_{H,\mathrm M}(U_k;S_k,(R_k,S_{k+1})) := \bigl(\widehat T_{H,\mathrm M}^1(U_k;S_k,(R_k,S_{k+1})),\dots,\widehat T_{H,\mathrm M}^H(U_k;S_k,(R_k,S_{k+1}))\bigr).
\end{equation}
The next theorem follows the  same episode-boundary viewpoint as CTD. The weighted contraction mechanism is unchanged, while the perturbation term is now affine in the current boundary iterate.

\begin{theorem}[Undiscounted fixed-horizon episodic MTD]\label{thm:undisc-mtd}
Assume the episodic sampling model described above and $\rho_{H,\min}>0$. Let $\eta_{H,\mathrm M}^\star$ denote the unique fixed point of $\Pi_{H,\mathrm M}^{\Theta} T_H^\pi$. From any $\eta_0\in\mathcal F_{H,\mathrm M,\Theta}^{\X}$, let $U_0=I_{H,\mathrm M}(\eta_0)$ and let $\eta_k := I_{H,\mathrm M}^{-1}(U_k)$ be generated by \eqref{eq:undisc-recur} with $\widehat T_{H} = \widehat T_{H,\mathrm M}$. There exist explicit constants $C^H_\mathrm M, c^H_\mathrm M, \bar \alpha ^H_\mathrm M > 0$ and explicit thresholds $g^H_\mathrm M(\alpha)$ and $g^H_\mathrm M(\alpha, z)$ for which the following episode-boundary results hold in the natural unweighted error.

\textup{Constant step size.} If $\alpha_k \equiv \alpha \le \bar \alpha^H_\mathrm M$, then the episode-boundary error decays geometrically to an $O(\alpha)$ neighborhood:
\begin{equation}
\E\bigl[\ell_{H, \mathrm M, \infty}^{\mathrm{unw}}(\eta_{mH}, \eta_{H,\mathrm M}^\star)^2 \bigr] \le C^{H}_\mathrm M \ell_{H,\mathrm M, \infty}^{\mathrm{unw}}(\eta_0, \eta_{H,\mathrm M}^\star)^2(1- c^{H}_\mathrm M H\alpha)^m + C^{H}_\mathrm M \alpha.
\end{equation}
For the diminishing step size regimes below, write $\tau_m := mH+g+H-1$, so $\tau_0 = g+H-1$, where $g$ is the step-size offset.

\textup{Linearly-diminishing step size.} If $\alpha_k = \alpha / (k+g)$, $\alpha > 1/c^{H}_\mathrm M$, and $g \ge \max\{H-1,g^{H}_\mathrm M(\alpha)\}$, then the leading episode-boundary residual term decays in $O(1/m)$:
\begin{equation}
\E\bigl[\ell_{H,\mathrm M, \infty}^{\mathrm{unw}}(\eta_{mH}, \eta_{H,\mathrm M}^\star)^2 \bigr] \le C^{H}_\mathrm M \ell_{H, \mathrm M, \infty}^{\mathrm{unw}}(\eta_0, \eta_{H,\mathrm M}^\star)^2\left(\frac{\tau_0}{\tau_m} \right)^{c^{H}_\mathrm M\alpha} + \frac{C^{H}_\mathrm M \alpha^2}{c^{H}_\mathrm M\alpha -1} \cdot \frac{1}{\tau_m}.
\end{equation}
\textup{Polynomially-diminishing step size.} If $\alpha_k = \alpha/(k+g)^z$, $z \in (0,1)$, and $g \ge \max\{H-1,g^{H}_\mathrm M(\alpha, z)\}$, then the leading episode-boundary residual term decays in $O(1/m^z)$:
\begin{equation}
\E\bigl[\ell_{H, \mathrm M, \infty}^{\mathrm{unw}}(\eta_{mH}, \eta_{H,\mathrm M}^\star)^2 \bigr] \le C^{H}_\mathrm M \ell_{H, \mathrm M, \infty}^{\mathrm{unw}}(\eta_0, \eta_{H,\mathrm M}^\star)^2\exp\left(-\frac{c^{H}_\mathrm M \alpha}{1-z} \bigl(\tau^{1-z}_m - \tau^{1-z}_0 \bigr) \right) + \frac{C^{H}_\mathrm M \alpha}{\tau_m^z}.
\end{equation}
\end{theorem}
\begin{proof}[Proof sketch]
This follows the frozen episode-boundary CTD argument after replacing the Cram\'er embedding by the Gram-matrix embedding. The cumulative-occupancy contraction and universal weighted-to-unweighted conversion remain unchanged, while the signed-categorical perturbation bound is affine in the current boundary iterate. Details are given in Appendix~\ref{app:fh-mtd}.
\end{proof}

\begin{corollary}\label{cor:undisc-mtd}
To guarantee $\mathbb{E}[\ell_{H, \mathrm M,\infty}^{\mathrm{unw}}(\eta_{mH},\eta_{H, \mathrm M}^\star)] \le \varepsilon$ with the undiscounted fixed-horizon MTD recursion, it suffices to take $m = O(\varepsilon^{-2})$ episodes. The constant is polynomial in $H$, $1/\rho_{H,\min}$, $1+\log(H\lvert\X\rvert)$, and the support-geometry constants. Equivalently, the transition count is $k=mH$.
\end{corollary}

\subsection{Explicit horizon dependence}
The choice $\lambda=H/(H+1)$ gives the episode-level contraction gap $\kappa_H=\rho_{H,\min}/(2(H+1))$ and the universal metric conversion $\lambda^{-H}\le e$. Let $p_H^\star:=\max\{2,\lceil\log(H\lvert\X\rvert)\rceil\}$, $q_H:=(H\lvert\X\rvert)^{2/p_H^\star}$, and $L_H^{\mathrm{sm}}:=(p_H^\star-1)q_H^2=O(1+\log(H\lvert\X\rvert))$. Appendix~\ref{sec:fh-ctd-drift-app} establishes
\begin{equation}
c_{\mathrm C}^H\ge\frac{\kappa_H}{4},\qquad \bar\alpha_{\mathrm C}^H\ge\min\left\{\frac{1}{H},\frac{\kappa_H^2}{32HL_H^{\mathrm{sm}}}\right\},\qquad \frac{a_{H,\mathrm C,4}^{\mathrm{fh}}}{a_{H,\mathrm C,2}^{\mathrm{fh}}}\le\frac{1536e^2HL_H^{\mathrm{sm}}B_{H,\mathrm C}^2}{\kappa_H^3}.
\end{equation}
Thus the constant-step residual is $O(HL_H^{\mathrm{sm}}B_{H,\mathrm C}^2\kappa_H^{-3}\alpha)$. For the linear schedule with $\alpha=2/a_{H,\mathrm C,2}^{\mathrm{fh}}$, its leading episode-boundary residual is
\begin{equation}
O\left(\frac{H^2L_H^{\mathrm{sm}}B_{H,\mathrm C}^2}{\kappa_H^4}\cdot\frac{1}{m}\right)=O\left(\frac{H^2(H+1)^4L_H^{\mathrm{sm}}B_{H,\mathrm C}^2}{\rho_{H,\min}^4}\cdot\frac{1}{m}\right).
\end{equation}
For MTD, let $G_{H,\mathrm M}:=\overline C_{H,\mathrm M}^{\mathrm{pot}}(1+2\lVert U_{H,\mathrm M}^\star\rVert_{H,2,\infty}^2)$, where Appendix~\ref{sec:fh-mtd-drift-app} defines $\overline C_{H,\mathrm M}^{\mathrm{pot}}$ explicitly in the support-geometry constants. The corresponding constant-step and linear-schedule residuals are
\begin{equation}
O\left(\frac{HG_{H,\mathrm M}}{\kappa_H}\alpha\right),\qquad O\left(\frac{H^2G_{H,\mathrm M}}{\kappa_H^2}\cdot\frac{1}{m}\right),
\end{equation}
respectively. Moreover, $\overline C_{H,\mathrm M}^{\mathrm{pot}}=O(L_H^{\mathrm{sm}}((C_{H,\mathrm M}^{\mathrm{move}})^2/\kappa_H^2+(C_{H,\mathrm M}^{\mathrm{end}})^2/\kappa_H))$. Thus no factor of the form $\lambda^{-H}$ or $\exp(1/\kappa_H)$ remains, and the displayed dependence outside the named support-geometry quantities is polynomial in $H$.

\section{Experiments}\label{sec:experiments}
We include a small tabular study to check that the behavior predicted by the finite-iteration bounds is visible in practice. For the discounted experiments, we use a five-state sticky ring with self-transition probability $0.94$, discount factor $\gamma=0.9$, deterministic state-dependent rewards, and $51$ uniformly spaced categorical atoms. We compare i.i.d. uniform state queries with a single Markovian trajectory. The fixed-horizon experiment uses $H=5$, two states at each phase, and cumulative occupancy $\rho_{H,\min}=0.1$. We report the mean squared statewise supremum Cram\'er error over $32$ independent runs, with interquartile bands.

In the discounted experiments, the constant, linear, and polynomial schedules are $0.04$, $100/(k+400)$, and $0.25(400/(k+400))^{0.7}$. Their fixed-horizon counterparts are $0.05$, $150/(k+600)$, and $0.25(600/(k+600))^{0.7}$. Figure~\ref{fig:synthetic-rates} shows a persistent constant-step neighborhood and decay under both diminishing schedules. Log-log fits over the final $40\%$ of each run give slopes $-1.01$, $-1.00$, and $-0.97$ for the linear schedule in the i.i.d., Markovian, and fixed-horizon settings. The corresponding polynomial-schedule slopes are $-0.70$, $-0.71$, and $-0.70$. These observations agree with the $O(1/k)$, $O(1/k^z)$, $O(1/m)$, and $O(1/m^z)$ squared-error rates, including under cumulative rather than phasewise support.

\begin{figure}[t]
\centering
\includegraphics[width=\linewidth]{figures/synthetic_rates.pdf}
\caption{Finite-iteration behavior of tabular CTD. Solid colored curves show means over $32$ runs, shaded regions show interquartile ranges, and black dashed segments show log-log fits over the final $40\%$ of the linear and polynomial runs. The fixed-horizon example has phase-dependent state occupancy.}
\label{fig:synthetic-rates}
\end{figure}

\section{Discussion of the results}
Theorems~\ref{thm:ctd-main2} and~\ref{thm:mtd-main} show that the standard categorical recursions are structurally simpler than they may initially appear. After suitable statewise isometric embeddings, both CTD and MTD become asynchronous SA schemes that update one state block at a time and contract in a block-supremum norm. This common structure is what makes a unified finite-iteration analysis possible. The real point of departure is the perturbation geometry. CTD lives in a bounded-noise regime because bounded categorical supports and the Cram\'er geometry give uniform samplewise control, whereas the signed-categorical MTD perturbation is affine in the centered error $\lVert U-U^\star \rVert_{2,\infty}$.

In the undiscounted fixed-horizon half, the learned object is a horizon-indexed stack, and one sampled transition updates the whole stack block at the visited state because every horizon shares the same first reward and differs only in the bootstrapped tail. There is no discount-driven contraction to exploit. A weighted horizon norm supplies the proof metric, while the final guarantees are converted to the natural unweighted metric with a universal factor because $\lambda=H/(H+1)$. The theory is stated at episode boundaries because the averaged drift changes with phase. Across all settings, constant steps give a controllable $O(\alpha)$ neighborhood, linearly-diminishing steps give $O(1/k)$ squared error, and polynomially-diminishing steps give $O(1/k^z)$. Further discussion is deferred to Appendix~\ref{app:discus}.

\section{Conclusion}
We established finite-iteration guarantees for asynchronous categorical distributional temporal-difference learning in the scalar Cramér and multivariate signed-categorical MMD settings. The scope of the present theory is tabular finite-state policy evaluation. We do not treat control or policy improvement scenarios, and our guarantees are expectation bounds rather than high-probability bounds. Within that scope, the discounted results cover i.i.d.\ sampled transition tuples and Markovian trajectories, while the undiscounted results treat the fixed-horizon episodic regime under the exact online recursions.

\section*{Acknowledgments}

The authors are grateful to Prof. Zaiwei Chen for helpful feedback during the preparation of this manuscript.

\newpage
\bibliography{refs}

\newpage
\appendix
\section*{Appendix Table of Contents}
\addcontentsline{toc}{section}{Appendix}
\markboth{Appendix}{Appendix}
\startcontents[appendix]
\printcontents[appendix]{l}{1}{\setcounter{tocdepth}{3}}

\section{Common finite-iteration ingredients in the block-supremum geometry}
\label{app:common}

We collect the abstract ingredients \citep{robbins1951stochastic, ljung2003analysis, borkar2008stochastic, kushner2003stochastic} used in the proofs of Theorems~\ref{thm:ctd-main2} and~\ref{thm:mtd-main}. Throughout this appendix, let
\begin{equation}
V := \prod_{s\in \X} E(s)
\end{equation}
be a finite product of Euclidean state blocks. For $U\in V$ and $p\in[2,\infty)$, define
\begin{equation}
\lVert U \rVert_{2,\infty} := \max_{s\in \X}\lVert U(s) \rVert_2, \qquad \lVert U \rVert_{2,p} := \left(\sum_{s\in \X}\lVert U(s) \rVert_2^p\right)^{1/p}.
\end{equation}

Because $V$ is finite-dimensional and each block $E(s)$ is Euclidean, all gradients, smoothness statements, and Moreau-envelope constructions below are taken with respect to the product Euclidean inner product on $V$. In particular, $\lVert \cdot \rVert_{2,p}$ is a genuine norm on a finite-dimensional Euclidean space, so the standard smoothness results for squared $\ell_p$-type norms apply directly.

For each state $s\in \X$, let $P_s:V\to V$ denote the coordinate projector onto block $s$.

\subsection{Norm comparison and smoothing}

\begin{lemma}[Block-norm comparison]
\label{lem:block-norm-comparison}
For every $U\in V$ and every $p\in[2,\infty)$,
\begin{equation}
\lVert U \rVert_{2,\infty} \le \lVert U \rVert_{2,p} \le \lvert \X \rvert^{1/p}\lVert U \rVert_{2,\infty}.
\end{equation}
\end{lemma}

\begin{proof}
The left inequality is immediate because the maximum of finitely many nonnegative numbers is bounded above by their $\ell_p$ norm. For the right inequality,
\begin{equation}
\lVert U \rVert_{2,p}^p = \sum_{s\in \X}\lVert U(s) \rVert_2^p \le \sum_{s\in \X}\lVert U \rVert_{2,\infty}^p = \lvert \X \rvert \lVert U \rVert_{2,\infty}^p.
\end{equation}
Taking $p$th roots proves the claim.
\end{proof}

\begin{proposition}[Choice of the smoothing exponent]
\label{prop:pstar-choice}
Let
\begin{equation}
p^\star := \max\{2,\lceil \log \lvert \X \rvert\rceil\}.
\end{equation}
Then
\begin{equation}
\lvert \X \rvert^{2/p^\star} \le e^2.
\end{equation}
Consequently,
\begin{equation}
(p^\star-1)\lvert \X \rvert^{2/p^\star} \le e^2 \log \lvert \X \rvert \qquad\text{for }\lvert \X \rvert\ge 2.
\end{equation}
\end{proposition}

\begin{proof}
If $\lvert \X \rvert\le e^2$, then $p^\star=2$ and the claim is immediate. If $\lvert \X \rvert>e^2$, then $p^\star\ge \log \lvert \X \rvert$, hence
\begin{equation}
\lvert \X \rvert^{2/p^\star} = \exp\left(\frac{2\log\lvert \X \rvert}{p^\star}\right) \le e^2.
\end{equation}
For the second claim, first suppose $2\le\lvert \X\rvert\le e^2$. Then $p^\star=2$, so its left-hand side equals $\lvert \X\rvert$. The function $x\mapsto e^2\log x-x$ is increasing on $[2,e^2]$ and is positive at $x=2$, hence $\lvert \X\rvert\le e^2\log\lvert \X\rvert$. If $\lvert \X\rvert>e^2$, then $p^\star-1\le\log\lvert \X\rvert$, and multiplying this inequality by $\lvert \X\rvert^{2/p^\star}\le e^2$ gives the result.
\end{proof}

\begin{proposition}[Smooth block potential]
\label{prop:block-smooth-potential}
For $p\in[2,\infty)$, define
\begin{equation}
g_p(U) := \frac{1}{2}\lVert U \rVert_{2,p}^2.
\end{equation}
Then $g_p$ is convex and $(p-1)$-smooth \citep{beck2017first} with respect to $\lVert \cdot \rVert_{2,p}$.
\end{proposition}

\begin{proof}
This is the standard smoothness of the squared $\ell_p$ norm for $p\ge 2$, applied to the mixed norm $\lVert \cdot \rVert_{2,p}$ on the finite-dimensional product space $V$. Equivalently, after choosing orthonormal bases in each block $E(s)$, the space $V$ identifies with $\mathbb R^N$ for some finite $N$, and $\lVert \cdot \rVert_{2,p}$ becomes a norm of $\ell_p$ type on grouped coordinates. The standard $\ell_p$ smoothness estimate therefore gives
\begin{equation}
g_p(U') \le g_p(U) + \langle \nabla g_p(U), U'-U\rangle + \frac{p-1}{2}\lVert U'-U \rVert_{2,p}^2
\end{equation}
for all $U,U'\in V$.
\end{proof}

\begin{proposition}[Generalized Moreau envelope \citep{moreau1965proximite, bauschke2017convex} for block-supremum norms]
\label{prop:moreau-envelope}
Fix $U^\star\in V$, $\vartheta>0$, and $p\in[2,\infty)$. Define
\begin{equation}
M_{\vartheta,p}(U) := \inf_{W\in V}\Bigl\{\frac{1}{2}\lVert W-U^\star \rVert_{2,\infty}^2 + \frac{1}{2\vartheta\lvert\X\rvert^{2/p}}\lVert U-W \rVert_{2,p}^2\Bigr\}.
\end{equation}
Then $M_{\vartheta,p}$ is convex and $(p-1)/\vartheta$-smooth with respect to $\lVert \cdot \rVert_{2,p}$. Moreover, for every $U\in V$,
\begin{equation}
(1+\vartheta)M_{\vartheta,p}(U) \le \frac{1}{2}\lVert U-U^\star \rVert_{2,\infty}^2 \le (1+\vartheta \lvert \X \rvert^{2/p})M_{\vartheta,p}(U).
\end{equation}
\end{proposition}

\begin{proof}
This is the generalized Moreau-envelope approximation theorem used by \citet{chen2020finite, chen2024lyapunov}, specialized to
\begin{equation}
h_1(U) := \frac{1}{2}\lVert U-U^\star \rVert_{2,\infty}^2, \qquad h_2(U) := \frac{1}{2}\lVert U \rVert_{2,p}^2.
\end{equation}
The proximal parameter is $\vartheta\lvert\X\rvert^{2/p}$, so Proposition~\ref{prop:block-smooth-potential} gives the exact smoothness constant $(p-1)/(\vartheta\lvert\X\rvert^{2/p})$. Since $\lvert\X\rvert^{2/p}\ge1$, the larger constant $(p-1)/\vartheta$ is also valid. To verify the comparison directly, set $E:=U-U^\star$, $Z:=W-U^\star$, and $q:=\lvert\X\rvert^{2/p}$. Lemma~\ref{lem:block-norm-comparison} gives
\begin{equation}
\inf_Z\left\{\frac12\lVert Z\rVert_{2,\infty}^2+\frac{1}{2\vartheta q}\lVert E-Z\rVert_{2,\infty}^2\right\}\le M_{\vartheta,p}(U)\le\inf_Z\left\{\frac12\lVert Z\rVert_{2,\infty}^2+\frac{1}{2\vartheta}\lVert E-Z\rVert_{2,\infty}^2\right\}.
\end{equation}
The two infima equal $\lVert E\rVert_{2,\infty}^2/(2(1+\vartheta q))$ and $\lVert E\rVert_{2,\infty}^2/(2(1+\vartheta))$, respectively, which proves the stated inequalities.
\end{proof}
\begin{remark}[Use in the fixed-horizon appendices]
In Appendices~\ref{app:fh-ctd} and~\ref{app:fh-mtd}, the fixed-horizon recursions are first rewritten on a flattened weighted product space indexed by horizon-state pairs $(h,s)$. After this flattening, the weighted fixed-horizon norm becomes an ordinary block-supremum norm on a finite product of Euclidean blocks, so the present appendix applies verbatim with the number of blocks equal to the cardinality of that flattened index set.
\end{remark}
The categorical recursions do not range over all of the ambient product space $V$. For CTD, valid iterates lie in the embedded product of probability simplices. For MTD, they lie in the embedded affine space of mass-$1$ signed categorical measures. The general finite-iteration theorems are stated on a vector space and might therefore appear to require contraction and perturbation estimates at points outside these categorical domains, where the estimates used below need not hold. The following lemma closes this gap, by showing that samplewise preservation and convexity keep every iterate in the categorical domain, while the proof of the general theorem uses operator assumptions only at such iterates and their comparisons with the fixed point. The Moreau envelope and its smoothness remain defined on $V$, so restricting the operator assumptions does not alter the Lyapunov calculation or its constants.
\begin{lemma}[Invariant-domain restriction]
\label{lem:invariant-domain}
Let $\mathcal D\subseteq V$ be nonempty, closed, and convex. Suppose that $U_0\in\mathcal D$ and that a full-sample map $\widehat H$ satisfies
\begin{equation}
\widehat H(U;y)\in\mathcal D \qquad\text{for every }U\in\mathcal D\text{ and every admissible sample }y.
\end{equation}
If $0\le\alpha_k\le1$ and
\begin{equation}
U_{k+1}=(1-\alpha_k)U_k+\alpha_k\widehat H(U_k;Y_k),
\end{equation}
then $U_k\in\mathcal D$ for every $k$. Consequently, the contractive SA arguments of \citet{chen2020finite,chen2024lyapunov} remain valid when their contraction, Lipschitz, and perturbation assumptions are imposed only on $\mathcal D$, provided that the fixed point belongs to $\mathcal D$ and the sampled recursion preserves $\mathcal D$.
\end{lemma}

\begin{proof}
The invariance claim follows by induction because $U_{k+1}$ is a convex combination of two points in $\mathcal D$. For the second claim, translate the recursion by its fixed point $U^\star$ and set $\mathcal D^\star:=\mathcal D-U^\star$. Then $0\in\mathcal D^\star$, every centered iterate belongs to $\mathcal D^\star$, and every comparison of two iterates or of an iterate with the fixed point remains inside $\mathcal D^\star$. The proofs in \citet{chen2020finite,chen2024lyapunov} use contractivity and samplewise Lipschitz continuity only for such comparisons. Their smooth Lyapunov function is defined on the ambient Euclidean space $V$, but it is evaluated only at centered iterates and their increments. Thus the Moreau-envelope smoothness and norm-comparison steps remain ambient-space statements, while every operator assumption is needed only on $\mathcal D^\star$. Repeating the same one-step inequalities gives the same constants and recursions under the restricted-domain assumptions.
\end{proof}

\subsection{Abstract i.i.d.\ finite-iteration framework}

Assume first that $(S_k,\xi_k)_{k\ge 0}$ are i.i.d., that the query-state marginal is $S_k\sim\rho$, and let
\begin{equation}
\rho_{\min} := \min_{s\in \X}\rho(s) > 0.
\end{equation}

Let $\mathcal D\subseteq V$ be nonempty, closed, and convex. Let $\mathcal O:\mathcal D\to\mathcal D$ be a contraction in $\lVert \cdot \rVert_{2,\infty}$ with modulus $\beta\in(0,1)$ and fixed point $U^\star\in\mathcal D$:
\begin{equation}
\lVert \mathcal O U-\mathcal O U'\rVert_{2,\infty} \le \beta \lVert U-U'\rVert_{2,\infty} \quad\text{ for all } U,U'\in \mathcal D.
\end{equation}
Let $\widehat T(U;s,\xi)$ be a sampled target satisfying
\begin{equation}
\mathbb E\left[ \widehat T(U_k;S_k,\xi_k) \mid U_k, S_k=s\right] = (\mathcal O U_k)(s).
\end{equation}
Consider the asynchronous recursion
\begin{equation}
U_{k+1} = U_k + \alpha_k P_{S_k}\left(\widehat T(U_k;S_k,\xi_k)-U_k(S_k)\right).
\end{equation}
Define
\begin{equation}
\widehat H(U;s,\xi):=U+P_s\bigl(\widehat T(U;s,\xi)-U(s)\bigr),
\end{equation}
and assume that $U_0\in\mathcal D$ and $\widehat H(U;s,\xi)\in\mathcal D$ for every $U\in\mathcal D$ and every admissible $(s,\xi)$.

For the centered recursion, write
\begin{equation}
W_k := U_k-U^\star.
\end{equation}
Assume there exists a finite constant $A^{\mathrm{iid}}\ge 0$ such that
\begin{equation}
\mathbb E\left[\left\lVert\widehat T(U_k;S_k,\xi_k)-(\mathcal O U_k)(S_k)\right\rVert_2^2\mid U_k, S_k\right] \le A^{\mathrm{iid}}\bigl(1+\lVert W_k \rVert_{2,\infty}^2\bigr) \qquad\text{a.s.}
\end{equation}

\begin{proposition}[Abstract i.i.d.\ finite-iteration bound]
\label{prop:abstract-iid}
Fix any $\vartheta > 0$ and define
\begin{equation}
\beta_{\rho} := 1 - \rho_{\min}(1-\beta), \qquad a_1 := \frac{1 + \vartheta \lvert \X \rvert^{2/p^\star}}{1+\vartheta}, \qquad a_2 := 1 - \beta_{\rho}\sqrt{\frac{1 + \vartheta \lvert \X \rvert^{2/p^\star}}{1+\vartheta}},
\end{equation}
and
\begin{equation}
a_3 := \frac{4(p^\star-1)\lvert \X \rvert^{2/p^\star}(A^{\mathrm{iid}}+2)(1+\vartheta)}{\vartheta}, \qquad a_4 := \frac{2(p^\star-1)\lvert \X \rvert^{2/p^\star}A^{\mathrm{iid}}(1+\vartheta)}{\vartheta}.
\end{equation}
Since $a_2 \to 1-\beta_{\rho}>0$ as $\vartheta\downarrow 0$, the condition $a_2 >0$ holds for all sufficiently small $\vartheta$. If $a_2 > 0$, $(\alpha_k)_{k \ge 0}$ is nonincreasing, and $\alpha_0 \le \min\{1,a_2/a_3\}$, then for all $k \ge 0$,
\begin{equation}
\mathbb{E}\left[\lVert W_k \rVert_{2,\infty}^2\right] \le a_1 \lVert W_0 \rVert_{2,\infty}^2 \prod_{j=0}^{k-1}(1-a_2\alpha_j) + a_4 \sum_{i=0}^{k-1} \alpha_i^2 \prod_{j=i+1}^{k-1}(1-a_2\alpha_j).
\end{equation}
In particular:
\begin{enumerate}
\item if $\alpha_k \equiv \alpha$ and $\alpha \le \min\{1,a_2/a_3\}$, then for all $k \ge 0$,
\begin{equation}
\mathbb{E}\left[\lVert W_k \rVert_{2,\infty}^2\right] \le a_1 \lVert W_0 \rVert_{2,\infty}^2 (1-a_2\alpha)^k + \frac{a_4}{a_2}\alpha,
\end{equation}

\item if $\alpha_k = \alpha/(k+h)$, $\alpha > 1/a_2$, and
\begin{equation}
h \ge \max\left\{1,\alpha,\frac{\alpha a_3}{a_2}\right\},
\end{equation}
then for all $k \ge 0$,
\begin{equation}
\mathbb{E}\left[\lVert W_k \rVert_{2,\infty}^2\right] \le a_1 \lVert W_0 \rVert_{2,\infty}^2 \left(\frac{h}{k+h}\right)^{a_2\alpha} + \frac{4e\,\alpha^2 a_4}{a_2\alpha-1}\cdot \frac{1}{k+h},
\end{equation}

\item if $\alpha_k = \alpha/(k+h)^z$ with $z \in (0,1)$ and
\begin{equation}
h \ge \max\left\{1,\alpha^{1/z}, \left(\frac{\alpha a_3}{a_2}\right)^{1/z}, \left(\frac{2z}{a_2\alpha}\right)^{1/(1-z)} \right\},
\end{equation}
then for all $k \ge 0$,
\begin{equation}
\mathbb{E}\left[\lVert W_k \rVert_{2,\infty}^2\right] \le a_1 \lVert W_0 \rVert_{2,\infty}^2 \exp\left(-\frac{a_2\alpha}{1-z}\bigl((k+h)^{1-z} - h^{1-z}\bigr)\right) + \frac{2\alpha a_4}{a_2}\cdot \frac{1}{(k+h)^z}.
\end{equation}
\end{enumerate}
\end{proposition}

\begin{proof}
Lemma~\ref{lem:invariant-domain} shows that every iterate belongs to $\mathcal D$.
Define the averaged asynchronous operator
\begin{equation}
H_{\rho}(U) := U + \sum_{s\in \X}\rho(s)P_s\bigl((\mathcal O U)(s)-U(s)\bigr).
\end{equation}
For each state block $s$,
\begin{equation}
\bigl(H_\rho(U)-H_\rho(V)\bigr)(s) =(1-\rho(s))(U(s)-V(s)) + \rho(s)\bigl((\mathcal O U)(s)-(\mathcal O V)(s)\bigr).
\end{equation}
Since $\mathcal O$ is a $\beta$-contraction in $\lVert \cdot \rVert_{2,\infty}$,
\begin{equation}
\begin{aligned}
\left\lVert\bigl(H_\rho(U)-H_\rho(V)\bigr)(s) \right\rVert_2 &\le \bigl((1-\rho(s))+\rho(s)\beta\bigr)\lVert U-V \rVert_{2,\infty}\\
&= \bigl(1-\rho(s)(1-\beta)\bigr)\lVert U-V \rVert_{2,\infty}\\
&\le \beta_\rho \lVert U-V \rVert_{2,\infty}.
\end{aligned}
\end{equation}
Taking the maximum over $s$ on both sides gives
\begin{equation}
\lVert H_\rho(U)-H_\rho(V) \rVert_{2,\infty} \le \max_{s\in\X}\bigl(1-\rho(s)(1-\beta)\bigr)\lVert U-V \rVert_{2,\infty} = \beta_\rho \lVert U-V \rVert_{2,\infty},
\end{equation}
since the coefficient $1-\rho(s)(1-\beta)$ is decreasing in $\rho(s)$. By Lemma~\ref{lem:invariant-domain}, Corollary~2.1 of \citet{chen2020finite} applies to the centered recursion on $\mathcal D-U^\star$, with contraction norm $\lVert \cdot \rVert_{2,\infty}$, smoothing norm $\lVert \cdot \rVert_{2,p^\star}$, Lyapunov function $M_{\vartheta,p^\star}$ from Proposition~\ref{prop:moreau-envelope}, and noise constant $A^{\mathrm{iid}}$. The approximation factors in Proposition~\ref{prop:moreau-envelope} and the smoothness constant in Proposition~\ref{prop:block-smooth-potential} produce exactly the displayed constants. The constant-step, linearly-diminishing, and polynomially-diminishing step size bounds follow by specializing the displayed recursion.
\end{proof}

\subsection{Poisson-equation Markovian finite-iteration framework}

Let $(S_k)_{k\ge0}$ be an irreducible Markov chain on $\X$ with transition matrix $Q$ and stationary distribution $\mu_{\X}$. Aperiodicity is not needed in this subsection. Define
\begin{equation}
\mu_{\min}:=\min_{s\in\X}\mu_{\X}(s)>0,\qquad
Z_Q:=\bigl(I-Q+\bm 1\mu_{\X}^{\mathsf T}\bigr)^{-1},\qquad
C_Q^{\mathrm{Pois}}:=\max_{s\in\X}\sum_{x\in\X}\lvert (Z_Q)_{sx}\rvert.
\end{equation}
The inverse exists for every finite irreducible chain, including periodic chains. The quantity $C_Q^{\mathrm{Pois}}$ is an explicit Poisson condition number that records the temporal dependence of the trajectory.

We retain the invariant domain $\mathcal D$, the fixed point $U^\star$, and the sampled full-replacement map $\widehat H$ from the i.i.d.\ framework. Let us write $W_k:=U_k-U^\star$ and define
\begin{equation}
h(W,s):=P_s\bigl((\mathcal O(U^\star+W))(s)-U^\star(s)-W(s)\bigr),\qquad
\bar h(W):=\sum_{s\in\X}\mu_{\X}(s)h(W,s).
\end{equation}
The centered recursion is
\begin{equation}
W_{k+1}=W_k+\alpha_k\bigl(h(W_k,S_k)+\Delta_k\bigr),
\end{equation}
where
\begin{equation}
\Delta_k:=\widehat H(U_k;S_k,\xi_k)-U_k-h(W_k,S_k).
\end{equation}
Assume
\begin{equation}
\mathbb E[\Delta_k\mid\mathcal G_k]=0,\qquad
\lVert\Delta_k\rVert_{2,\infty}\le A_3\lVert W_k\rVert_{2,\infty}+B_3
\end{equation}
almost surely, where $\mathcal G_k$ contains the history through $U_k$ and $S_k$ before the transition sample at time $k$ is revealed. For each $W\in\mathcal D-U^\star$, define the vector-valued Poisson solution
\begin{equation}
V_W:=Z_Q\bigl(h(W,\cdot)-\bar h(W)\bigr),
\end{equation}
which satisfies
\begin{equation}
V_W(s)-\sum_{x\in\X}Q(s,x)V_W(x)=h(W,s)-\bar h(W)
\end{equation}
and
\begin{equation}
\max_{s\in\X}\lVert V_W(s)-V_{W'}(s)\rVert_{2,\infty}\le A_Q\lVert W-W'\rVert_{2,\infty},\qquad
A_Q:=\max\{1,4C_Q^{\mathrm{Pois}}\}.
\end{equation}
Indeed, both $h(\cdot,s)$ and its stationary average are $2$-Lipschitz because $\mathcal O$ is contractive, and $h(0,s)=V_0(s)=0$.

\begin{proposition}[Poisson-equation Markovian finite-iteration bound]
\label{prop:abstract-markov}
Let
\begin{equation}
q_{\X}:=\lvert\X\rvert^{2/p^\star},\qquad
\beta_\mu:=1-\mu_{\min}(1-\beta).
\end{equation}
Fix $\vartheta>0$ and define
\begin{equation}
r_\mu:=\frac{1+\vartheta q_{\X}}{1+\vartheta},\qquad
\omega_\mu:=1-\beta_\mu\sqrt{r_\mu}.
\end{equation}
Assume $\omega_\mu>0$. Set
\begin{equation}
\begin{gathered}
\eta_\mu:=2\omega_\mu,\qquad
L_\mu:=\frac{p^\star-1}{\vartheta},\qquad
l_\mu:=2(1+\vartheta),\qquad
u_\mu:=2(1+\vartheta q_{\X}),\\
\mathsf A:=(3+A_3)^2,\qquad
\mathsf B:=2B_3^2,\qquad
K_{\mu,1}:=u_\mu L_\mu q_{\X}A_Q,\\
K_{\mu,0}:=\frac{2u_\mu(1+2\mathsf A K_{\mu,1})}{l_\mu}\lVert W_0\rVert_{2,\infty}^2+4\mathsf B K_{\mu,1}.
\end{gathered}
\end{equation}
Suppose $U_0\in\mathcal D$, every full-sample replacement preserves $\mathcal D$, and the step sizes $\alpha_k\le 1$ for all $k$. Then the following bounds hold.

If $\alpha_k\equiv\alpha$ with $\alpha>0$ and
\begin{equation}
\alpha\le\min\left\{1,\frac{\eta_\mu}{\mathsf A(5+2\eta_\mu)K_{\mu,1}}\right\},
\end{equation}
then, for all $k\ge0$,
\begin{equation}
\mathbb E\left[\lVert W_{k+1}\rVert_{2,\infty}^2\right]\le
K_{\mu,0}\exp\left(-\frac{\eta_\mu\alpha k}{2}\right)
+18\mathsf B K_{\mu,1}\alpha+\frac{40\mathsf B K_{\mu,1}}{\eta_\mu}\alpha.
\end{equation}

If $\alpha_k=\alpha/(k+h)$, $\alpha>2/\eta_\mu$, and
\begin{equation}
h\ge\max\left\{\mathsf A\alpha(5\alpha+8)K_{\mu,1},2\right\},
\end{equation}
then, for all $k\ge0$,
\begin{equation}
\begin{gathered}
\mathbb E\left[\lVert W_{k+1}\rVert_{2,\infty}^2\right]\le
K_{\mu,0}\left(\frac{h}{k+h}\right)^{\eta_\mu\alpha/2}
+\frac{2\mathsf B K_{\mu,1}\alpha}{k+h}
+\frac{8e\mathsf B(5+4\eta_\mu)K_{\mu,1}\alpha^2}{(\eta_\mu\alpha/2-1)(k+h)}.
\end{gathered}
\end{equation}

If $\alpha_k=\alpha/(k+h)^z$ with $\alpha>0$ and $z\in(0,1)$, define $\chi_\mu:=z/\alpha+\eta_\mu$. If
\begin{equation}
h\ge\max\left\{\left(\frac{2\mathsf A\alpha(5+2\chi_\mu)K_{\mu,1}}{\eta_\mu}\right)^{1/z},2\right\},
\end{equation}
then, for all $k\ge0$,
\begin{equation}
\begin{gathered}
\mathbb E\left[\lVert W_{k+1}\rVert_{2,\infty}^2\right]\le
K_{\mu,0}\exp\left(-\frac{\eta_\mu\alpha}{2(1-z)}
\bigl((k+h)^{1-z}-h^{1-z}\bigr)\right)\\
+\frac{2\mathsf B K_{\mu,1}\alpha}{(k+h)^z}
+\frac{8\mathsf B(5+2\chi_\mu)K_{\mu,1}\alpha}{\eta_\mu(k+h)^z}.
\end{gathered}
\end{equation}
\end{proposition}

\begin{proof}
The stationary asynchronous map $W\mapsto W+\bar h(W)$ is a $\beta_\mu$-contraction. The same dual-norm argument used in the proof of Proposition~\ref{prop:abstract-iid} gives
\begin{equation}
\langle\nabla M_{\vartheta,p^\star}(U^\star+W),\bar h(W)\rangle\le-\eta_\mu M_{\vartheta,p^\star}(U^\star+W).
\end{equation}
The envelope is $L_\mu$-smooth in $\lVert\cdot\rVert_{2,p^\star}$ and satisfies
\begin{equation}
l_\mu M_{\vartheta,p^\star}(U^\star+W)\le\lVert W\rVert_{2,\infty}^2\le u_\mu M_{\vartheta,p^\star}(U^\star+W).
\end{equation}
Moreover,
\begin{equation}
\lVert h(W,s)\rVert_{2,\infty}\le2\lVert W\rVert_{2,\infty},
\end{equation}
the Poisson solution is $A_Q$-Lipschitz, and the martingale term has coefficients $A_3,B_3$. The proof of \citet[Theorem A.2]{haque2024stochastic} therefore applies with their constants (from Assumptions 3.1–3.3) $A_1=2$, $B_1=0$, $A_2=A_Q$, $B_2=0$, $A=(A_1+A_3+1)^2=\mathsf A$, and $B=2(B_1+B_3+B_2/A_2)^2=\mathsf B$. Lemma~\ref{lem:invariant-domain} permits the proof to be restricted to $\mathcal D-U^\star$ because every operator and Poisson comparison is evaluated only at iterates in the invariant domain. Substituting the three step-size schedules into that theorem gives the displayed bounds.
\end{proof}

For compact theorem statements, define $c_\mu:=\eta_\mu/2$ and
\begin{equation}
\begin{gathered}
C_{\mu,0}:=\max\left\{\frac{2u_\mu(1+2\mathsf A K_{\mu,1})}{l_\mu},4\mathsf B K_{\mu,1}\right\}, \\
C_{\mu,\mathrm{cs}}:=18\mathsf B K_{\mu,1}+\frac{40\mathsf B K_{\mu,1}}{\eta_\mu}, \\
C_{\mu,\mathrm{lin}}:=2\mathsf B K_{\mu,1}+8e\mathsf B(5+4\eta_\mu)K_{\mu,1}, \\
C_{\mu,\mathrm{poly}}:=2\mathsf B K_{\mu,1}+\frac{40\mathsf B K_{\mu,1}}{\eta_\mu}, \\
C_\mu^{\mathrm{mk}}:=2e\max\{C_{\mu,0},C_{\mu,\mathrm{cs}},C_{\mu,\mathrm{lin}},C_{\mu,\mathrm{poly}}\}.
\end{gathered}
\end{equation}
The bounds in Proposition~\ref{prop:abstract-markov} then have the compact forms used in Theorems~\ref{thm:ctd-main2} and~\ref{thm:mtd-main}. Those theorems are written for the $k$th iterate rather than the $(k+1)$st iterate. Replacing $k$ in Proposition~\ref{prop:abstract-markov} by $k-1$ introduces factors of at most $e$ in the transient terms and at most $2$ in the residual terms because $c_\mu<1$, the admissible schedules have $\alpha_0\le1$, and $h\ge2$. These factors are included in $C_\mu^{\mathrm{mk}}$. For the linear schedule, $c_\mu\alpha>1$ also implies $\alpha\le\alpha^2/(c_\mu\alpha-1)$. For the polynomial schedule, $5+2\chi_\mu\le5(1+\chi_\mu)$. These estimates give the residual terms displayed in the main text with all dependence on $k$ explicit.


\subsection{Explicit dependence of the discounted constants}

The preceding propositions retain exact constants for arbitrary smoothing parameters. The following specialization makes their dependence on the statistically relevant quantities explicit.

\begin{proposition}[Explicit smoothing choice]
\label{prop:explicit-discounted-constants}
Let
\begin{equation}
q_{\X}:=\lvert\X\rvert^{2/p^\star},\qquad \delta_\rho:=\rho_{\min}(1-\beta),\qquad \delta_\mu:=\mu_{\min}(1-\beta).
\end{equation}
For the i.i.d. result, choose $\vartheta_\rho:=\delta_\rho/q_{\X}$. Then the constants in Proposition~\ref{prop:abstract-iid} satisfy
\begin{equation}
\begin{gathered}
a_1\le2,\qquad a_2\ge\frac{\delta_\rho}{2},\qquad a_3\le\frac{8(p^\star-1)q_{\X}^2(A^{\mathrm{iid}}+2)}{\delta_\rho},\qquad a_4\le\frac{4(p^\star-1)q_{\X}^2A^{\mathrm{iid}}}{\delta_\rho}.
\end{gathered}
\end{equation}
In particular, the sufficient constant-step condition
\begin{equation}
\alpha\le\min\left\{1,\frac{\delta_\rho^2}{16(p^\star-1)q_{\X}^2(A^{\mathrm{iid}}+2)}\right\}
\end{equation}
gives
\begin{equation}
\mathbb E\left[\lVert W_k\rVert_{2,\infty}^2\right]\le2\lVert W_0\rVert_{2,\infty}^2\exp\left(-\frac{\delta_\rho\alpha k}{2}\right)+\frac{8(p^\star-1)q_{\X}^2A^{\mathrm{iid}}}{\delta_\rho^2}\alpha.
\end{equation}

For the Markovian result, choose $\vartheta_\mu:=\delta_\mu/q_{\X}$. Then the constants in Proposition~\ref{prop:abstract-markov} satisfy
\begin{equation}
\eta_\mu\ge\delta_\mu,\qquad L_\mu=\frac{(p^\star-1)q_{\X}}{\delta_\mu},\qquad u_\mu\le4,\qquad
K_{\mu,1}\le\frac{4(p^\star-1)q_{\X}^2A_Q}{\delta_\mu}.
\end{equation}
Consequently, the sufficient constant-step condition
\begin{equation}
\alpha\le\min\left\{1,\frac{\delta_\mu^2}{36\mathsf A(p^\star-1)q_{\X}^2A_Q}\right\}
\end{equation}
implies that the constant-step residual obeys
\begin{equation}
18\mathsf B K_{\mu,1}\alpha+\frac{40\mathsf B K_{\mu,1}}{\eta_\mu}\alpha
\le\frac{232(p^\star-1)q_{\X}^2A_Q\mathsf B}{\delta_\mu^2}\alpha.
\end{equation}
Thus the Poisson analysis removes the explicit $t_\alpha$ factor. The trajectory dependence is instead isolated in $A_Q=\max\{1,4C_Q^{\mathrm{Pois}}\}$. Moreover, $q_{\X}\le e^2$ and $(p^\star-1)q_{\X}^2=O(\log\lvert\X\rvert)$ by Proposition~\ref{prop:pstar-choice}.
\end{proposition}

\begin{proof}
For either $\delta\in\{\delta_\rho,\delta_\mu\}$ and $\vartheta=\delta/q_{\X}$,
\begin{equation}
\frac{1+\vartheta q_{\X}}{1+\vartheta}\le1+\delta\le2
\end{equation}
and
\begin{equation}
\sqrt{\frac{1+\vartheta q_{\X}}{1+\vartheta}}\le\sqrt{1+\vartheta q_{\X}}\le1+\frac{\delta}{2}.
\end{equation}
Consequently,
\begin{equation}
1-(1-\delta)\sqrt{\frac{1+\vartheta q_{\X}}{1+\vartheta}}\ge1-(1-\delta)\left(1+\frac{\delta}{2}\right)\ge\frac{\delta}{2}.
\end{equation}
Substituting $\vartheta=\delta_\rho/q_{\X}$ and $1+\vartheta q_{\X}=1+\delta_\rho\le2$ into $a_3$ and $a_4$ gives the i.i.d. bounds. For the Markovian bounds, $1+\vartheta_\mu q_{\X}=1+\delta_\mu\le2$ gives $u_\mu\le4$, while the preceding lower bound gives $\eta_\mu\ge\delta_\mu$. Direct substitution into Proposition~\ref{prop:abstract-markov}, together with $\eta_\mu\le2$ and $\delta_\mu\le1$, gives the stated threshold and residual estimates.
\end{proof}

\section{CTD details and proof of Theorem~\ref{thm:ctd-main2}}
\label{app:ctd}

We verify the hypotheses of Appendix~\ref{app:common} for the discounted scalar categorical temporal-difference recursion and thereby prove Theorem~\ref{thm:ctd-main2}. The proof has three parts. First, we record the statewise Cram\'er isometry and the induced block-supremum contraction of the projected Bellman operator. Second, we verify the samplewise Lipschitz property of the sampled target map. Third, we verify the centered perturbation bounds required by the i.i.d.\ and Markovian finite-iteration arguments.

\subsection{Statewise Cram\'er isometry and contraction}

For each state $s\in \X$, let
\begin{equation}
\Theta(s)=\{\theta_1(s)<\cdots<\theta_d(s)\}\subset\mathbb R
\end{equation}
and let $\mathcal F^{\X}_{\mathrm C,\Theta}$ denote the class of state-indexed categorical laws supported on these statewise grids. For
\begin{equation}
\eta(s)=\sum_{i=1}^d p_i(s)\delta_{\theta_i(s)},
\end{equation}
define the cumulative masses and the grid separations
\begin{equation}
c^\eta_j(s) := \sum_{i=1}^j p_i(s),
\qquad
\Delta_j(s) := \theta_{j+1}(s)-\theta_j(s),
\qquad
j=1,\dots,d-1.
\end{equation}
The statewise Cram\'er embedding is
\begin{equation}
I_s(\eta(s)) := \bigl(\sqrt{\Delta_1(s)}\,c^\eta_1(s), \dots, \sqrt{\Delta_{d-1}(s)}\,c^\eta_{d-1}(s),0\bigr)^\top \in\mathbb R^d.
\end{equation}

\begin{proposition}[Statewise Cram\'er isometry]
\label{prop:ctd-isometry}
For every state $s\in \X$ and all $\mu,\nu\in \mathcal F_{\mathrm C,\Theta(s)}$,
\begin{equation}
\ell_{\mathrm C}(\mu,\nu)^2 = \sum_{j=1}^{d-1}\Delta_j(s)\bigl(c^\mu_j(s)-c^\nu_j(s)\bigr)^2 = \lVert I_s(\mu)-I_s(\nu) \rVert_2^2.
\end{equation}
Consequently, with the product embedding
\begin{equation}
I_{\mathrm C}(\eta):=(I_s(\eta(s)))_{s\in \X},
\end{equation}
we have
\begin{equation}
\lVert I_{\mathrm C}(\eta)-I_{\mathrm C}(\eta') \rVert_{2,\infty} = \ell_{\mathrm C,\infty}(\eta,\eta').
\end{equation}
\end{proposition}

\begin{proof}
For laws supported on the ordered grid $\Theta(s)$, the cumulative distribution functions are piecewise
constant on the intervals
\begin{equation}
[\theta_j(s),\theta_{j+1}(s)), \qquad j=1,\dots,d-1,
\end{equation}
with values $c^\eta_j(s)$ and $c^\nu_j(s)$ respectively. Therefore,
\begin{equation}
\ell_{\mathrm C}(\mu,\nu)^2 = \int_{\mathbb R}\bigl(F_\mu(z)-F_\nu(z)\bigr)^2\,\mathrm dz = \sum_{j=1}^{d-1}\Delta_j(s)\bigl(c^\mu_j(s)-c^\nu_j(s)\bigr)^2,
\end{equation}
which is exactly $\lVert I_s(\mu)-I_s(\nu) \rVert_2^2$. Taking the maximum over $s\in \X$ gives the product-space identity.
\end{proof}

Recall from Section~\ref{sec:ctd-new2} that
\begin{equation}
\mathcal O_{\mathrm C} := I_{\mathrm C}\circ \Pi_{\mathrm C}^{\Theta}T^\pi \circ I_{\mathrm C}^{-1}.
\end{equation}

Define the embedded categorical domain
\begin{equation}
\mathcal D_{\mathrm C}:=I_{\mathrm C}(\mathcal F^{\X}_{\mathrm C,\Theta})\subseteq V.
\end{equation}

\begin{proposition}[Invariance of the CTD domain]
\label{prop:ctd-domain}
The set $\mathcal D_{\mathrm C}$ is compact and convex, $\mathcal O_{\mathrm C}(\mathcal D_{\mathrm C})\subseteq\mathcal D_{\mathrm C}$, and the full-sample replacement map
\begin{equation}
\widehat H_{\mathrm C}(U;s,(r,s')):=U+P_s\bigl(\widehat T_{\mathrm C}(U;s,(r,s'))-U(s)\bigr)
\end{equation}
satisfies
\begin{equation}
\widehat H_{\mathrm C}(U;s,(r,s'))\in\mathcal D_{\mathrm C}
\end{equation}
for every $U\in\mathcal D_{\mathrm C}$ and every admissible $(s,r,s')$. Consequently, if $U_0\in\mathcal D_{\mathrm C}$ and $0\le\alpha_k\le1$, then every CTD iterate belongs to $\mathcal D_{\mathrm C}$.
\end{proposition}

\begin{proof}
The class $\mathcal F^{\X}_{\mathrm C,\Theta}$ is a finite product of probability simplices, and the cumulative-mass embedding is linear on each simplex. Its image $\mathcal D_{\mathrm C}$ is therefore compact and convex. The projected Bellman operator and every sampled projected target are categorical laws on the prescribed support. Hence $\mathcal O_{\mathrm C}$ maps $\mathcal D_{\mathrm C}$ into itself, and replacing one state block by a sampled embedded target gives $\widehat H_{\mathrm C}(U;s,(r,s'))\in\mathcal D_{\mathrm C}$. The final claim follows from Lemma~\ref{lem:invariant-domain}.
\end{proof}

\begin{proposition}[Embedded CTD contraction]
\label{prop:ctd-contraction-app}
For all $U,U'\in I_{\mathrm C}(\mathcal F^{\X}_{\mathrm C,\Theta})$,
\begin{equation}
\lVert \mathcal O_{\mathrm C}U-\mathcal O_{\mathrm C}U'\rVert_{2,\infty} \le \sqrt{\gamma}\,\lVert U-U'\rVert_{2,\infty}.
\end{equation}
\end{proposition}

\begin{proof}
Let $\eta:=I_{\mathrm C}^{-1}(U)$ and $\eta':=I_{\mathrm C}^{-1}(U')$. By Proposition~\ref{prop:ctd-isometry} and the contraction of $\mathcal O_{\mathrm C}$ in the supremum Cram\'er metric \citep{bellemare2023distributional},
\begin{equation}
\lVert \mathcal O_{\mathrm C}U-\mathcal O_{\mathrm C}U'\rVert_{2,\infty} = \ell_{\mathrm C,\infty}(\Pi_{\mathrm C}^{\Theta}T^\pi\eta,\Pi_{\mathrm C}^{\Theta}T^\pi\eta') \le \sqrt{\gamma}\,\ell_{\mathrm C,\infty}(\eta,\eta') = \sqrt{\gamma}\,\lVert U-U'\rVert_{2,\infty}.
\end{equation}
\end{proof}

\subsection{Unbiased sampled target and samplewise Lipschitz continuity}

Recall the sampled CTD target from Section~\ref{sec:ctd-new2}:
\begin{equation}
\widehat T_{\mathrm C}(U_k;S_k,(R_k,S'_k)) := I_{S_k} \Bigl(\Pi_{\mathrm C}^{\Theta(S_k)}\bigl((f_{R_k,\gamma})_\# I_{S'_k}^{-1}(U_k(S'_k))\bigr)\Bigr),
\end{equation}
where $f_{r,\gamma}(z)=r+\gamma z$.

\begin{proposition}[Conditional unbiasedness of the sampled CTD target]
\label{prop:ctd-unbiased}
For every $U\in I_{\mathrm C}(\mathcal F^{\X}_{\mathrm C,\Theta})$ and every state $s\in \X$,
\begin{equation}
\mathbb E\left[\widehat T_{\mathrm C}(U;s,(R_k,S'_k))\mid U_k=U, S_k=s\right] = (\mathcal O_{\mathrm C}U)(s).
\end{equation}
\end{proposition}

\begin{proof}
Condition on $U_k=U$ and $S_k=s$. By definition,
\begin{equation}
\widehat T_{\mathrm C}(U;s,(R_k,S'_k)) = I_s\Bigl(\Pi_{\mathrm C}^{\Theta(s)}\bigl((f_{R_k,\gamma})_\# I_{S'_k}^{-1}(U(S'_k))\bigr)\Bigr).
\end{equation}
Taking conditional expectation with respect to the one-step transition and reward law under $\pi$ gives exactly the Bellman expectation at state $s$. Applying the deterministic statewise projection and then the embedding $I_s$ yields
\begin{equation}
\mathbb E\left[\widehat T_{\mathrm C}(U;s,(R_k,S'_k))
\mid U_k=U, S_k=s\right] = I_s\bigl((\Pi_{\mathrm C}^{\Theta}T^\pi I_{\mathrm C}^{-1}U)(s)\bigr) = (\mathcal O_{\mathrm C}U)(s).
\end{equation}
Here the expectation commutes with the categorical projection and the cumulative-mass embedding because both maps are linear in the input measure.
\end{proof}

\begin{proposition}[CTD samplewise Lipschitz continuity]
\label{prop:ctd-lipschitz}
For every $s\in \X$, every $(r,s')\in[0,1]\times \X$, and every $U,U'\in I_{\mathrm C}(\mathcal F^{\X}_{\mathrm C,\Theta})$,
\begin{equation}
\left\lVert \widehat T_{\mathrm C}(U;s,(r,s')) - \widehat T_{\mathrm C}(U';s,(r,s'))\right\rVert_2 \le \lVert U-U'\rVert_{2,\infty}.
\end{equation}
Consequently, if
\begin{equation}
\widehat H_{\mathrm C}(U;s,(r, s')) := U+P_s\bigl(\widehat T_{\mathrm C}(U;s,(r, s'))-U(s)\bigr),
\end{equation}
then
\begin{equation}
\lVert\widehat H_{\mathrm C}(U;s,(r, s'))-\widehat H_{\mathrm C}(U';s,(r, s'))\rVert_{2,\infty} \le \lVert U-U'\rVert_{2,\infty}.
\end{equation}
\end{proposition}

\begin{proof}
Fix $s\in \X$ and $(r,s')\in[0,1]\times \X$. Let
\begin{equation}
\mu := I_{s'}^{-1}(U(s')), \qquad \nu := I_{s'}^{-1}(U'(s')).
\end{equation}
Then
\begin{equation}
\widehat T_{\mathrm C}(U;s,(r,s')) = I_s\bigl(\Pi_{\mathrm C}^{\Theta(s)}((f_{r,\gamma})_\# \mu)\bigr)
\end{equation}
and similarly for $U'$.
Using the statewise Cram\'er isometry, the nonexpansiveness of the categorical projection in the Cram\'er metric, and the deterministic pushforward contraction by $\sqrt{\gamma}$,
\begin{equation}
\begin{aligned}
\left\lVert\widehat T_{\mathrm C}(U;s,(r,s')) - \widehat T_{\mathrm C}(U';s,(r,s'))\right\rVert_2 &= \ell_{\mathrm C,s}\bigl(\Pi_{\mathrm C}^{\Theta(s)}((f_{r,\gamma})_\# \mu), \Pi_{\mathrm C}^{\Theta(s)}((f_{r,\gamma})_\# \nu)\bigr)\\
&\le \ell_{\mathrm C,s}((f_{r,\gamma})_\# \mu,(f_{r,\gamma})_\# \nu)\\
&\le\sqrt{\gamma}\,\ell_{\mathrm C,s'}(\mu,\nu)\\
&=\sqrt{\gamma}\,\lVert U(s')-U'(s') \rVert_2\\
&\le\lVert U-U'\rVert_{2,\infty}.
\end{aligned}
\end{equation}
For the full-sample map, every block other than $s$ is unchanged, while the $s$th block is replaced by the sampled target. Therefore,
\begin{equation}
\resizebox{\linewidth}{!}{$\displaystyle
\lVert \widehat H_{\mathrm C}(U;s,(r, s'))-\widehat H_{\mathrm C}(U';s,(r, s'))\rVert_{2,\infty} = \max\Bigl\{\max_{x\ne s}\lVert U(x)-U'(x) \rVert_2,\lVert\widehat T_{\mathrm C}(U;s,(r, s'))-\widehat T_{\mathrm C}(U';s,(r, s'))\rVert_2\Bigr\},
$}
\end{equation}
which is bounded by $\lVert U-U'\rVert_{2,\infty}$ by the previous result.
\end{proof}

\subsection{Uniform perturbation bound}

\begin{proposition}[Uniform bound on embedded categorical blocks]
\label{prop:ctd-block-bound}
For every state $s\in \X$ and every embedded categorical law $U(s)\in I_s(\mathcal F_{\mathrm C,\Theta(s)})$,
\begin{equation}
\lVert U(s) \rVert_2^2 \le \theta_d(s)-\theta_1(s).
\end{equation}
\end{proposition}

\begin{proof}
Write
\begin{equation}
U(s)=I_s(\eta(s)) = \bigl( \sqrt{\Delta_1(s)}\,c^\eta_1(s), \dots, \sqrt{\Delta_{d-1}(s)}\,c^\eta_{d-1}(s),0\bigr)^\top.
\end{equation}
Because $\eta(s)$ is a probability law, every cumulative mass $c^\eta_j(s)$ lies in $[0,1]$. Hence
\begin{equation}
\lVert U(s) \rVert_2^2 = \sum_{j=1}^{d-1}\Delta_j(s)\bigl(c^\eta_j(s)\bigr)^2 \le \sum_{j=1}^{d-1}\Delta_j(s) = \theta_d(s)-\theta_1(s).
\end{equation}
\end{proof}

\begin{proposition}[CTD centered perturbation bound]
\label{prop:ctd-perturbation}
Let the CTD expected full-sample map be
\begin{equation}
F_{\mathrm C}(U,s) := U + P_s\bigl((\mathcal O_{\mathrm C}U)(s)-U(s)\bigr)
\end{equation}
and
\begin{equation}
\Delta_k^{\mathrm C} := \widehat H_{\mathrm C}(U_k;S_k,(R_k,S'_k)) - F_{\mathrm C}(U_k,S_k).
\end{equation}
Then
\begin{equation}
\mathbb E\left[ \Delta_k^{\mathrm C} \,\mid\, \mathcal G_k \right] = 0
\end{equation}
and
\begin{equation}
\lVert \Delta_k^{\mathrm C} \rVert_{2,\infty} \le 2B_{\mathrm C} \qquad\text{a.s. for all }k,
\end{equation}
where
\begin{equation}
B_{\mathrm C} := \max_{s\in \X}\sqrt{\theta_d(s)-\theta_1(s)}.
\end{equation}
Consequently,
\begin{equation}
\mathbb E\left[ \left\lVert \widehat T_{\mathrm C}(U_k;S_k,(R_k,S'_k))-(\mathcal O_{\mathrm C}U_k)(S_k) \right\rVert_2^2 \,\mid\, U_k,\ S_k \right] \le 4B_{\mathrm C}^2.
\end{equation}
\end{proposition}

\begin{proof}
The conditional mean identity follows from Proposition~\ref{prop:ctd-unbiased}. For the pathwise bound, only the $S_k$th block is nonzero, so
\begin{equation}
\lVert \Delta_k^{\mathrm C} \rVert_{2,\infty} = \left\lVert \widehat T_{\mathrm C}(U_k;S_k,(R_k,S'_k)) - (\mathcal O_{\mathrm C}U_k)(S_k) \right\rVert_2.
\end{equation}
Both terms on the right are embedded categorical laws at state $S_k$. By Proposition~\ref{prop:ctd-block-bound}, each has Euclidean norm at most $B_{\mathrm C}$. Therefore, the triangle inequality yields
\begin{equation}
\lVert \Delta_k^{\mathrm C} \rVert_{2,\infty} \le 2B_{\mathrm C}.
\end{equation}
Squaring gives the conditional second-moment bound.
\end{proof}

\subsection{Completion of the proof of Theorem~\ref{thm:ctd-main2}}

\begin{proof}[Proof of Theorem~\ref{thm:ctd-main2}]
We now verify the hypotheses of Appendix~\ref{app:common} for the discounted CTD recursion.

Proposition~\ref{prop:ctd-domain} verifies the restricted-domain hypotheses with $\mathcal D=\mathcal D_{\mathrm C}$. In particular, the sampled recursion remains in $\mathcal D_{\mathrm C}$ whenever $U_0\in\mathcal D_{\mathrm C}$ and the step sizes are at most one. All contraction, Lipschitz, and perturbation estimates below are therefore required only on $\mathcal D_{\mathrm C}$.

First, Proposition~\ref{prop:ctd-contraction-app} gives the block-supremum contraction
\begin{equation}
\lVert \mathcal O_{\mathrm C} U - \mathcal O_{\mathrm C} U' \rVert_{2,\infty} \le \beta_{\mathrm C} \lVert U-U' \rVert_{2,\infty}, \qquad \beta_{\mathrm C} := \sqrt{\gamma}.
\end{equation}

Second, Proposition~\ref{prop:ctd-lipschitz} gives the target-level estimate
\begin{equation}
\lVert \widehat T_{\mathrm C}(U;s,(r,s')) - \widehat T_{\mathrm C}(U';s,(r,s')) \rVert_2 \le \beta_{\mathrm C} \lVert U-U' \rVert_{2,\infty} \le \lVert U-U' \rVert_{2,\infty}.
\end{equation}

Third, Proposition~\ref{prop:ctd-perturbation} gives the centered pathwise perturbation bound
\begin{equation}
\mathbb{E}[\Delta_k^{\mathrm C} \mid \mathcal G_k] = 0, \qquad \lVert \Delta_k^{\mathrm C} \rVert_{2,\infty} \le 2B_{\mathrm C} \quad \text{a.s.},
\end{equation}
so Proposition~\ref{prop:abstract-markov} applies with
\begin{equation}
A_{3,\mathrm C}=0,\qquad B_{3,\mathrm C}=2B_{\mathrm C}.
\end{equation}
In particular, the i.i.d. conditional second-moment condition in Proposition~\ref{prop:abstract-iid} holds with
\begin{equation}
A_{\mathrm C}^{\mathrm{iid}} := 4B_{\mathrm C}^2.
\end{equation}

Let $q_{\X}:=\lvert\X\rvert^{2/p^\star}$ and choose
\begin{equation}
\vartheta_{\mathrm C,\rho}:=\frac{\rho_{\min}(1-\beta_{\mathrm C})}{q_{\X}}.
\end{equation}
Applying Proposition~\ref{prop:abstract-iid} with
\begin{equation}
\beta = \beta_{\mathrm C}, \qquad A^{\mathrm{iid}} = A_{\mathrm C}^{\mathrm{iid}}, \qquad \vartheta = \vartheta_{\mathrm C,\rho}, \qquad p^\star = \max\{2,\lceil \log \lvert \X \rvert \rceil\},
\end{equation}
yields Theorem~\ref{thm:ctd-main2}(i) with
\begin{equation}
a_{\mathrm C,1}^{\mathrm{iid}} := \frac{1 + \vartheta_{\mathrm C,\rho}\lvert \X \rvert^{2/p^\star}}{1+\vartheta_{\mathrm C,\rho}}, \qquad a_{\mathrm C,2}^{\mathrm{iid}} := 1-\bar\beta_{\mathrm C,\rho}\sqrt{\frac{1 + \vartheta_{\mathrm C,\rho}\lvert \X \rvert^{2/p^\star}}{1+\vartheta_{\mathrm C,\rho}}},
\end{equation}
where
\begin{equation}
\bar\beta_{\mathrm C,\rho} := 1-\rho_{\min}(1-\beta_{\mathrm C}),
\end{equation}
and
\begin{equation}
\begin{gathered}
a_{\mathrm C,3}^{\mathrm{iid}} := \frac{4(p^\star-1)\lvert \X \rvert^{2/p^\star}(4B_{\mathrm C}^2+2)(1+\vartheta_{\mathrm C,\rho}\lvert \X \rvert^{2/p^\star})}{\vartheta_{\mathrm C,\rho}},\\
a_{\mathrm C,4}^{\mathrm{iid}} := \frac{8(p^\star-1)\lvert \X \rvert^{2/p^\star}B_{\mathrm C}^2(1+\vartheta_{\mathrm C,\rho}\lvert \X \rvert^{2/p^\star})}{\vartheta_{\mathrm C,\rho}}.
\end{gathered}
\end{equation}
The compact constants and thresholds in Theorem~\ref{thm:ctd-main2}(i) can therefore be chosen as
\begin{equation}
\begin{gathered}
c_{\mathrm C}^{\mathrm{iid}}:=a_{\mathrm C,2}^{\mathrm{iid}},\qquad
\bar\alpha_{\mathrm C}^{\mathrm{iid}}:=\min\left\{1,\frac{a_{\mathrm C,2}^{\mathrm{iid}}}{a_{\mathrm C,3}^{\mathrm{iid}}}\right\},\\
C_{\mathrm C}^{\mathrm{iid}}:=\max\left\{a_{\mathrm C,1}^{\mathrm{iid}},4e a_{\mathrm C,4}^{\mathrm{iid}},\frac{2a_{\mathrm C,4}^{\mathrm{iid}}}{a_{\mathrm C,2}^{\mathrm{iid}}}\right\},\\
h_{\mathrm C}^{\mathrm{iid}}(\alpha):=\left\lceil\max\left\{1,\alpha,\frac{\alpha a_{\mathrm C,3}^{\mathrm{iid}}}{a_{\mathrm C,2}^{\mathrm{iid}}}\right\}\right\rceil,\\
h_{\mathrm C}^{\mathrm{iid}}(\alpha,z):=\left\lceil\max\left\{1,\alpha^{1/z},\left(\frac{\alpha a_{\mathrm C,3}^{\mathrm{iid}}}{a_{\mathrm C,2}^{\mathrm{iid}}}\right)^{1/z},\left(\frac{2z}{a_{\mathrm C,2}^{\mathrm{iid}}\alpha}\right)^{1/(1-z)}\right\}\right\rceil.
\end{gathered}
\end{equation}

For the Markovian case, Proposition~\ref{prop:abstract-markov} applies with
\begin{equation}
A_{3,\mathrm C}=0,\qquad B_{3,\mathrm C}=2B_{\mathrm C},\qquad
\mathsf A_{\mathrm C}=9,\qquad \mathsf B_{\mathrm C}=8B_{\mathrm C}^2.
\end{equation}
Choose
\begin{equation}
\vartheta_{\mathrm C,\mu}:=\frac{\mu_{\min}(1-\beta_{\mathrm C})}{q_{\X}},
\end{equation}
and let $\eta_{\mathrm C,\mu}$, $K_{\mathrm C,0}$, and $K_{\mathrm C,1}$ be the explicit constants in Proposition~\ref{prop:abstract-markov} under these substitutions. The theorem drift constant and admissible constant step size can be taken as
\begin{equation}
c_{\mathrm C}^{\mathrm{mk}}:=\frac{\eta_{\mathrm C,\mu}}{2},\qquad
\bar\alpha_{\mathrm C}^{\mathrm{mk}}:=\min\left\{1,\frac{\eta_{\mathrm C,\mu}}{9(5+2\eta_{\mathrm C,\mu})K_{\mathrm C,1}}\right\}.
\end{equation}
For the diminishing schedules, one may take
\begin{equation}
h_{\mathrm C}^{\mathrm{mk}}(\alpha):=\left\lceil\max\left\{9\alpha(5\alpha+8)K_{\mathrm C,1},2\right\}\right\rceil
\end{equation}
and, with $\chi_{\mathrm C,\mu}:=z/\alpha+\eta_{\mathrm C,\mu}$,
\begin{equation}
h_{\mathrm C}^{\mathrm{mk}}(\alpha,z):=\left\lceil\max\left\{\left(\frac{18\alpha(5+2\chi_{\mathrm C,\mu})K_{\mathrm C,1}}{\eta_{\mathrm C,\mu}}\right)^{1/z},2\right\}\right\rceil.
\end{equation}
Set $C_{\mathrm C}^{\mathrm{mk}}:=C_\mu^{\mathrm{mk}}$ with $\mathsf A=\mathsf A_{\mathrm C}$, $\mathsf B=\mathsf B_{\mathrm C}$, $\eta_\mu=\eta_{\mathrm C,\mu}$, and $K_{\mu,1}=K_{\mathrm C,1}$ in the explicit definition following Proposition~\ref{prop:abstract-markov}. This coefficient is independent of $k$ and of the step size.
The constant, linearly-diminishing, and polynomially-diminishing step size bounds in Theorem~\ref{thm:ctd-main2} are therefore obtained by the corresponding specializations of Proposition~\ref{prop:abstract-iid} and Proposition~\ref{prop:abstract-markov}.
\end{proof}
\subsection{Proof of Corollary~\ref{cor:disc-ctd}}

\begin{proof}
We use the linearly-diminishing step size bounds in Theorem~\ref{thm:ctd-main2}.

For part \textup{(i)}, the linearly-diminishing i.i.d.\ bound in Theorem~\ref{thm:ctd-main2}(i) gives
\begin{equation}
\mathbb E\left[\ell_{\mathrm C,\infty}(\eta_k,\eta^\star)^2\right] \le a_{\mathrm C,1}^{\mathrm{iid}}\ell_{\mathrm C,\infty}(\eta_0,\eta^\star)^2 \left(\frac{h}{k+h}\right)^{a_{\mathrm C,2}^{\mathrm{iid}}\alpha} + \frac{4e\alpha^2 a_{\mathrm C,4}^{\mathrm{iid}}}{a_{\mathrm C,2}^{\mathrm{iid}}\alpha-1}\cdot \frac{1}{k+h}.
\end{equation}
Since $\alpha>1/a_{\mathrm C,2}^{\mathrm{iid}}$, one has $a_{\mathrm C,2}^{\mathrm{iid}}\alpha>1$, and therefore the first term is also $O((k+h)^{-1})$. Hence
\begin{equation}
\mathbb E\left[\ell_{\mathrm C,\infty}(\eta_k,\eta^\star)^2\right] = O\left(\frac{1}{k+h}\right).
\end{equation}
Thus $\mathbb E[\ell_{\mathrm C,\infty}(\eta_k,\eta^\star)]\le \varepsilon$ is guaranteed once
\begin{equation}
\mathbb E\left[\ell_{\mathrm C,\infty}(\eta_k,\eta^\star)^2\right]\le \varepsilon^2,
\end{equation}
which holds for $k=O(\varepsilon^{-2})$.

For part \textup{(ii)}, the linearly-diminishing Markovian bound in Theorem~\ref{thm:ctd-main2}, part \textup{(ii)}, gives
\begin{equation}
\mathbb E\left[\ell_{\mathrm C,\infty}(\eta_k,\eta^\star)^2\right] \le C_{\mathrm C}^{\mathrm{mk}}\bigl(1+\ell_{\mathrm C,\infty}(\eta_0,\eta^\star)^2\bigr)\left(\frac{h}{k+h}\right)^{c_{\mathrm C}^{\mathrm{mk}}\alpha}+\frac{C_{\mathrm C}^{\mathrm{mk}}\alpha^2}{c_{\mathrm C}^{\mathrm{mk}}\alpha-1}\cdot\frac{1}{k+h}.
\end{equation}
Since $\alpha>1/c_{\mathrm C}^{\mathrm{mk}}$, the first term is $O((k+h)^{-1})$. Therefore
\begin{equation}
\mathbb E\left[\ell_{\mathrm C,\infty}(\eta_k,\eta^\star)^2\right]=O\left(\frac{1}{k+h}\right),
\end{equation}
and hence $\mathbb E[\ell_{\mathrm C,\infty}(\eta_k,\eta^\star)]\le\varepsilon$ is ensured for $k=O(\varepsilon^{-2})$.
\end{proof}

\section{MTD details and proof of Theorem~\ref{thm:mtd-main}}
\label{app:mtd}

We now verify the abstract hypotheses of Appendix~\ref{app:common} for the discounted multivariate signed-categorical recursion.

\subsection{Statewise MMD isometry and contraction}

For each state $s\in \X$, let
\begin{equation}
\Theta(s)=\{\theta_1(s),\dots,\theta_d(s)\}\subset\mathbb R^q, \qquad \mathbb R_1^d := \Bigl\{ p\in\mathbb R^d : \sum_{i=1}^d p_i=1 \Bigr\}.
\end{equation}
Let $\mathcal M(\R^q)$ denote the space of finite signed Borel measures on $\R^q$. Define
\begin{equation}
\mathcal F^{\X}_{\mathrm M,\Theta} := \Bigl\{ \eta : \X\to \mathcal M(\mathbb R^q) : \eta(s)=\sum_{i=1}^d p_i(s)\delta_{\theta_i(s)}, \ p(s)\in \mathbb R_1^d \text{ for all }s\in \X \Bigr\}.
\end{equation}

Let $\kappa:\mathbb R^q\times\mathbb R^q\to\mathbb R$ be the characteristic kernel from Section~\ref{sec:mtd-new2}. For each state $s$, let
\begin{equation}
K_s := \bigl( \kappa(\theta_i(s),\theta_j(s)) \bigr)_{i,j=1}^d
\end{equation}
be the Gram matrix on $\Theta(s)$. For
\begin{equation}
\eta(s)=\sum_{i=1}^d p_i(s)\delta_{\theta_i(s)},
\end{equation}
define
\begin{equation}
I_s(\eta(s)) := K_s^{1/2}p(s)\in\mathbb R^d.
\end{equation}

\begin{proposition}[Statewise MMD isometry]
\label{prop:mtd-isometry}
For every state $s\in \X$ and all signed-categorical laws $\mu,\nu$ supported on $\Theta(s)$,
\begin{equation}
\mathrm{MMD}_\kappa(\mu,\nu)^2 = \lVert I_s(\mu)-I_s(\nu) \rVert_2^2.
\end{equation}
Consequently, with
\begin{equation}
I_{\mathrm M}(\eta):=(I_s(\eta(s)))_{s\in \X},
\end{equation}
one has
\begin{equation}
\lVert I_{\mathrm M}(\eta)-I_{\mathrm M}(\eta') \rVert_{2,\infty} = \ell_{\mathrm M,\infty}(\eta,\eta').
\end{equation}
\end{proposition}

\begin{proof}
Write
\begin{equation}
\mu=\sum_{i=1}^d p_i\delta_{\theta_i(s)}, \qquad \nu=\sum_{i=1}^d q_i\delta_{\theta_i(s)}.
\end{equation}
Then
\begin{equation}
\mu-\nu=\sum_{i=1}^d (p_i-q_i)\delta_{\theta_i(s)}.
\end{equation}
By the definition of MMD,
\begin{equation}
\mathrm{MMD}_\kappa(\mu,\nu)^2 = \sum_{i=1}^d\sum_{j=1}^d (p_i-q_i)(p_j-q_j)\kappa(\theta_i(s),\theta_j(s)) = (p-q)^\top K_s(p-q).
\end{equation}
Since $K_s$ is positive semidefinite,
\begin{equation}
(p-q)^\top K_s(p-q) = \lVert K_s^{1/2}(p-q) \rVert_2^2 = \lVert I_s(\mu)-I_s(\nu) \rVert_2^2.
\end{equation}
Taking the maximum over $s\in \X$ gives the product isometry.
\end{proof}

Recall from Section~\ref{sec:mtd-new2} that
\begin{equation}
\mathcal O_{\mathrm M} := I_{\mathrm M}\circ \Pi_{\mathrm M}^{\Theta}T^\pi \circ I_{\mathrm M}^{-1}.
\end{equation}

Define the embedded signed-categorical domain
\begin{equation}
\mathcal D_{\mathrm M}:=I_{\mathrm M}(\mathcal F^{\X}_{\mathrm M,\Theta})\subseteq V.
\end{equation}

\begin{proposition}[Invariance of the MTD domain]
\label{prop:mtd-domain}
The set $\mathcal D_{\mathrm M}$ is a closed affine subset of $V$, $\mathcal O_{\mathrm M}(\mathcal D_{\mathrm M})\subseteq\mathcal D_{\mathrm M}$, and the full-sample replacement map
\begin{equation}
\widehat H_{\mathrm M}(U;s,(r,s')):=U+P_s\bigl(\widehat T_{\mathrm M}(U;s,(r,s'))-U(s)\bigr)
\end{equation}
satisfies
\begin{equation}
\widehat H_{\mathrm M}(U;s,(r,s'))\in\mathcal D_{\mathrm M}
\end{equation}
for every $U\in\mathcal D_{\mathrm M}$ and every admissible $(s,r,s')$. Consequently, if $U_0\in\mathcal D_{\mathrm M}$ and $0\le\alpha_k\le1$, then every MTD iterate belongs to $\mathcal D_{\mathrm M}$.
\end{proposition}

\begin{proof}
At each state, the coefficient vectors of mass-$1$ signed categorical laws form the affine hyperplane $\mathbb R_1^d$. Its image under the linear Gram-matrix embedding is affine and closed, so their finite product $\mathcal D_{\mathrm M}$ is closed, affine, and therefore convex. The signed-categorical projection has mass $1$, both for the expected Bellman target and for each sampled Bellman target. Thus $\mathcal O_{\mathrm M}$ and every full-sample replacement preserve $\mathcal D_{\mathrm M}$. Lemma~\ref{lem:invariant-domain} gives the final claim.
\end{proof}

\begin{proposition}[Embedded MTD contraction]
\label{prop:mtd-contraction-app}
For all $U,U'\in I_{\mathrm M}(\mathcal F^{\X}_{\mathrm M,\Theta})$,
\begin{equation}
\lVert \mathcal O_{\mathrm M}U-\mathcal O_{\mathrm M}U'\rVert_{2,\infty} \le \gamma^{c/2}\lVert U-U'\rVert_{2,\infty}.
\end{equation}
\end{proposition}

\begin{proof}
Let $\eta:=I_{\mathrm M}^{-1}(U)$ and $\eta':=I_{\mathrm M}^{-1}(U')$. By Proposition~\ref{prop:mtd-isometry} and the contraction of $\mathcal O_{\mathrm M}$ in the supremum MMD metric \citep{wiltzer2024foundations},
\begin{equation}
\lVert \mathcal O_{\mathrm M}U-\mathcal O_{\mathrm M}U'\rVert_{2,\infty} = \ell_{\mathrm M,\infty}(\Pi_{\mathrm M}^{\Theta}T^\pi\eta,\Pi_{\mathrm M}^{\Theta}T^\pi\eta') \le \gamma^{c/2}\ell_{\mathrm M,\infty}(\eta,\eta') = \gamma^{c/2}\lVert U-U'\rVert_{2,\infty}.
\end{equation}
\end{proof}

\subsection{Unbiased sampled target, Lipschitz continuity, and perturbation control}

Recall the sampled MTD target from Section~\ref{sec:mtd-new2}:
\begin{equation}
\widehat T_{\mathrm M}(U_k;S_k,(R_k,S'_k)) := I_{S_k} \Bigl( \Pi_{\mathrm M}^{\Theta(S_k)} \bigl( (f_{R_k,\gamma})_\# I_{S'_k}^{-1}(U_k(S'_k)) \bigr) \Bigr).
\end{equation}

\begin{proposition}[Embedded MTD projection as an affine Euclidean projector]
\label{prop:mtd-proj-affine}
For each state $s \in \X$, let
\begin{equation}
\mathcal A_s := K_s^{1/2}\mathbb R_1^d \subset \mathbb R^d .
\end{equation}
For every mass-$1$ finite signed measure $\nu$ on $\mathbb R^q$, define
\begin{equation}
b_s(\nu) := \left( \int \kappa(\theta_i(s),y)\,d\nu(y) \right)_{i=1}^d \in \mathbb R^d
\end{equation}
and
\begin{equation}
\zeta_s(\nu) := K_s^{\dagger/2} b_s(\nu),
\end{equation}
where $K_s^\dagger$ is the Moore-Penrose pseudoinverse of $K_s$. Then $b_s(\nu) \in \operatorname{range}(K_s)$ and
\begin{equation}
I_{\mathrm M,s}\bigl(\Pi_{\mathrm M}^{\Theta(s)}\nu\bigr) = \operatorname{proj}_{\mathcal A_s}\bigl(\zeta_s(\nu)\bigr),
\end{equation}
where $\operatorname{proj}_{\mathcal A_s}$ is the Euclidean orthogonal projection onto $\mathcal A_s$. Consequently, the map
\begin{equation}
J_s(\nu) := I_{\mathrm M,s}\bigl(\Pi_{\mathrm M}^{\Theta(s)}\nu\bigr)
\end{equation}
is affine on the affine space of mass-$1$ finite signed measures.
\end{proposition}

\begin{proof}
Let $\mathcal H_\kappa$ be the RKHS of $\kappa$, let $\varphi : \mathbb R^q \to \mathcal H_\kappa$ denote the feature map, and define the linear operator
\begin{equation}
\Phi_s : \mathbb R^d \to \mathcal H_\kappa, \qquad \Phi_s p := \sum_{i=1}^d p_i \varphi(\theta_i(s)).
\end{equation}
Then
\begin{equation}
K_s = \Phi_s^*\Phi_s .
\end{equation}
For a mass-$1$ finite signed measure $\nu$, define its kernel mean embedding by
\begin{equation}
m_\nu := \int \varphi(y)\,d\nu(y) \in \mathcal H_\kappa .
\end{equation}
Let $P_s$ denote the orthogonal projection in $\mathcal H_\kappa$ onto $\operatorname{range}(\Phi_s)$. Since $\Phi_s^*(I-P_s)=0$, we have
\begin{equation}
b_s(\nu) = \Phi_s^* m_\nu = \Phi_s^* P_s m_\nu \in \operatorname{range}(\Phi_s^*\Phi_s) = \operatorname{range}(K_s).
\end{equation}

Now fix $p \in \mathbb R_1^d$ and write
\begin{equation}
\mu_p := \sum_{i=1}^d p_i \delta_{\theta_i(s)} .
\end{equation}
By the reproducing-kernel representation of MMD,
\begin{equation}
\mathrm{MMD}_\kappa(\mu_p,\nu)^2 = \lVert \Phi_s p - m_\nu \rVert_{\mathcal H_\kappa}^2 = \lVert \Phi_s p - P_s m_\nu \rVert_{\mathcal H_\kappa}^2 + \lVert (I-P_s)m_\nu \rVert_{\mathcal H_\kappa}^2 .
\end{equation}

We next identify $P_s m_\nu$ in coefficient form. Since $K_s=\Phi_s^*\Phi_s$, the standard Moore-Penrose identity gives
\begin{equation}
\Phi_s K_s^\dagger \Phi_s^* = P_s .
\end{equation}
Therefore,
\begin{equation}
\Phi_s\bigl(K_s^\dagger b_s(\nu)\bigr) = \Phi_s K_s^\dagger \Phi_s^* m_\nu = P_s m_\nu .
\end{equation}
Substituting this into the previous display yields
\begin{equation}
\lVert \Phi_s p - P_s m_\nu \rVert_{\mathcal H_\kappa} = \lVert \Phi_s(p-K_s^\dagger b_s(\nu)) \rVert_{\mathcal H_\kappa}.
\end{equation}
Since $\lVert \Phi_s x \rVert_{\mathcal H_\kappa}^2 = x^\top K_s x = \lVert K_s^{1/2}x \rVert_2^2$ for every $x \in \mathbb R^d$, we obtain
\begin{equation}
\lVert \Phi_s p - P_s m_\nu \rVert_{\mathcal H_\kappa} = \lVert K_s^{1/2}(p-K_s^\dagger b_s(\nu)) \rVert_2 = \lVert K_s^{1/2}p - K_s^{1/2}K_s^\dagger b_s(\nu) \rVert_2 .
\end{equation}
Because $K_s$ is symmetric positive semidefinite, spectral calculus gives
\begin{equation}
K_s^{1/2}K_s^\dagger = K_s^{\dagger/2}.
\end{equation}
Hence
\begin{equation}
\lVert \Phi_s p - P_s m_\nu \rVert_{\mathcal H_\kappa} = \lVert K_s^{1/2}p - K_s^{\dagger/2} b_s(\nu) \rVert_2 = \lVert K_s^{1/2}p - \zeta_s(\nu) \rVert_2 .
\end{equation}
Therefore
\begin{equation}
\mathrm{MMD}_\kappa(\mu_p,\nu)^2 = \lVert K_s^{1/2}p - \zeta_s(\nu) \rVert_2^2 + c_s(\nu),
\end{equation}
where
\begin{equation}
c_s(\nu) := \lVert (I-P_s)m_\nu \rVert_{\mathcal H_\kappa}^2
\end{equation}
does not depend on $p$.

Minimizing over $p \in \mathbb R_1^d$ is therefore equivalent to Euclidean projection of $\zeta_s(\nu)$ onto
\begin{equation}
\mathcal A_s = K_s^{1/2}\mathbb R_1^d .
\end{equation}
Since $I_{\mathrm M,s}\bigl(\sum_{i=1}^d p_i\delta_{\theta_i(s)}\bigr)=K_s^{1/2}p$, it follows that
\begin{equation}
I_{\mathrm M,s}\bigl(\Pi_{\mathrm M}^{\Theta(s)}\nu\bigr) = \operatorname{proj}_{\mathcal A_s}\bigl(\zeta_s(\nu)\bigr).
\end{equation}

Finally, $\nu \mapsto b_s(\nu)$ is linear, hence so is $\nu \mapsto \zeta_s(\nu)=K_s^{\dagger/2}b_s(\nu)$. Since orthogonal projection onto an affine subspace is an affine map, the composition
\begin{equation}
J_s(\nu) = \operatorname{proj}_{\mathcal A_s}\bigl(\zeta_s(\nu)\bigr)
\end{equation}
is affine on the affine space of mass-$1$ finite signed measures.
\end{proof}

\begin{proposition}[Conditional unbiasedness of the sampled MTD target]
\label{prop:mtd-unbiased}
For every admissible $U\in I_{\mathrm M}(\mathcal F^{\X}_{\mathrm M,\Theta})$ and every state $s\in \X$,
\begin{equation}
\mathbb E\left[ \widehat T_{\mathrm M}(U;s,(R_k,S'_k)) \,\mid\, U_k=U,\ S_k=s \right] = (\mathcal O_{\mathrm M}U)(s).
\end{equation}
\end{proposition}

\begin{proof}
Let $\eta:=I_{\mathrm M}^{-1}(U)$ and condition on $U_k=U$ and $S_k=s$. By Proposition~\ref{prop:mtd-proj-affine}, the statewise map
\begin{equation}
J_s(\nu):=I_s\bigl(\Pi_{\mathrm M}^{\Theta(s)}\nu\bigr)
\end{equation}
is affine on mass-$1$ finite signed measures. Hence
\begin{equation}
\begin{aligned}
\mathbb E\left[ \widehat T_{\mathrm M}(U;s,(R_k,S'_k)) \,\mid\, U_k=U,\ S_k=s \right] &= \mathbb E\left[ J_s\bigl((f_{R_k,\gamma})_\#\eta(S'_k)\bigr) \,\mid\, U_k=U,\ S_k=s \right] \\ 
&= J_s\left( \mathbb E\left[ (f_{R_k,\gamma})_\#\eta(S'_k) \,\mid\, U_k=U,\ S_k=s \right] \right)\\ 
&= J_s\bigl((T^\pi\eta)(s)\bigr)\\ 
&= I_s\bigl((\Pi_{\mathrm M}^{\Theta}T^\pi\eta)(s)\bigr)\\ 
&= (\mathcal O_{\mathrm M}U)(s).
\end{aligned}
\end{equation}
\end{proof}

\begin{proposition}[MTD samplewise Lipschitz continuity]
\label{prop:mtd-lipschitz}
For every $s\in \X$, every $(r,s')\in[0,1]^q\times \X$, and all $U,U'\in I_{\mathrm M}(\mathcal F^{\X}_{\mathrm M,\Theta})$,
\begin{equation}
\left\lVert \widehat T_{\mathrm M}(U;s,(r,s')) - \widehat T_{\mathrm M}(U';s,(r,s')) \right\rVert_2 \le \lVert U-U'\rVert_{2,\infty}.
\end{equation}
Consequently, if
\begin{equation}
\widehat H_{\mathrm M}(U;s,(r,s')) := U+P_s\bigl(\widehat T_{\mathrm M}(U;s,(r,s'))-U(s)\bigr),
\end{equation}
then
\begin{equation}
\lVert \widehat H_{\mathrm M}(U;s,(r,s'))-\widehat H_{\mathrm M}(U';s,(r,s')) \rVert_{2,\infty} \le \lVert U-U'\rVert_{2,\infty}.
\end{equation}
\end{proposition}

\begin{proof}
Fix $s\in \X$ and $(r,s')\in[0,1]^q\times \X$. Let
\begin{equation}
\mu:=I_{s'}^{-1}(U(s')), \qquad \nu:=I_{s'}^{-1}(U'(s')).
\end{equation}
Then
\begin{equation}
\widehat T_{\mathrm M}(U;s,(r,s')) = I_s\bigl(\Pi_{\mathrm M}^{\Theta(s)}((f_{r,\gamma})_\# \mu)\bigr)
\end{equation}
and likewise for $U'$. Using the statewise MMD isometry, the Euclidean projection representation from Proposition~\ref{prop:mtd-proj-affine}, and the $\gamma^{c/2}$ contraction of the pushforward in MMD,
\begin{equation}
\begin{aligned}
\left\lVert \widehat T_{\mathrm M}(U;s,(r,s')) - \widehat T_{\mathrm M}(U';s,(r,s')) \right\rVert_2 &= \mathrm{MMD}_\kappa\bigl( \Pi_{\mathrm M}^{\Theta(s)}((f_{r,\gamma})_\# \mu), \Pi_{\mathrm M}^{\Theta(s)}((f_{r,\gamma})_\# \nu) \bigr) \\ 
&\le \mathrm{MMD}_\kappa((f_{r,\gamma})_\# \mu,(f_{r,\gamma})_\# \nu) \\ 
&\le \gamma^{c/2}\,\mathrm{MMD}_\kappa(\mu,\nu) \\ 
&= \gamma^{c/2}\,\lVert U(s')-U'(s') \rVert_2 \\ 
&\le \lVert U-U'\rVert_{2,\infty}.
\end{aligned}
\end{equation}
For the full-sample map, blocks other than $s$ are unchanged, while block $s$ is replaced by the sampled target. Hence, for every $x \neq s$,
\begin{equation}
\bigl(\widehat H_{\mathrm M}(U;s,(r,s'))-\widehat H_{\mathrm M}(U';s,(r,s'))\bigr)(x) = U(x)-U'(x),
\end{equation}
and for the sampled block,
\begin{equation}
\bigl(\widehat H_{\mathrm M}(U;s,(r,s'))-\widehat H_{\mathrm M}(U';s,(r,s'))\bigr)(s) = \widehat T_{\mathrm M}(U;s,(r,s'))-\widehat T_{\mathrm M}(U';s,(r,s')).
\end{equation}
Taking the block supremum and using the target-level bound proves the claim.
\end{proof}

\begin{proposition}[Discounted MTD affine perturbation bound]
\label{prop:mtd-perturbation}
Let
\begin{equation}
U^\star := I_{\mathrm M}(\eta^\star), \qquad W_k:=U_k-U^\star, \qquad \beta_{\mathrm M} := \gamma^{c/2}.
\end{equation}
Here $\eta^\star$ denotes the unique fixed point of $\Pi_{\mathrm M}^{\Theta}T^\pi$. Define
\begin{equation}
B_{\mathrm M}^\star := \max_{s,s'\in \X} \sup_{r\in[0,1]^q} \left\lVert \widehat T_{\mathrm M}(U^\star;s,(r,s'))-(\mathcal O_{\mathrm M}U^\star)(s) \right\rVert_2
\end{equation}
Then $B_{\mathrm M}^\star<\infty$, and for every $s\in\X$, every $(r,s')\in[0,1]^q\times\X$, and every $U\in I_{\mathrm M}(\mathcal F^{\X}_{\mathrm M,\Theta})$,
\begin{equation}
\left\lVert \widehat T_{\mathrm M}(U;s,(r,s'))-(\mathcal O_{\mathrm M}U)(s) \right\rVert_2 \le 2\beta_{\mathrm M}\lVert U-U^\star \rVert_{2,\infty}+B_{\mathrm M}^\star.
\end{equation}
Consequently, for every $k\ge 0$,
\begin{equation}\label{eq:2ndmoment-mtd}
\mathbb E\left[ \left\lVert \widehat T_{\mathrm M}(U_k;S_k,(R_k,S'_k))-(\mathcal O_{\mathrm M}U_k)(S_k) \right\rVert_2^2 \,\mid\, U_k,S_k \right] \le 8\beta_{\mathrm M}^2\lVert W_k \rVert_{2,\infty}^2+2\bigl(B_{\mathrm M}^\star\bigr)^2.
\end{equation}
In particular, the i.i.d.\ conditional second-moment assumption in Proposition~\ref{prop:abstract-iid} holds with
\begin{equation}
A_{\mathrm M}^{\mathrm{iid}}:=8\beta_{\mathrm M}^2+2\bigl(B_{\mathrm M}^\star\bigr)^2.
\end{equation}
Moreover, if
\begin{equation}
F_{\mathrm M}(U,s) := U + P_s\bigl((\mathcal O_{\mathrm M}U)(s)-U(s)\bigr),
\end{equation}
\begin{equation}
\widehat H_{\mathrm M}(U;s,(r, s')) := U + P_s\bigl(\widehat T_{\mathrm M}(U;s,(r,s'))-U(s)\bigr),
\end{equation}
and
\begin{equation}
\Delta_k^{\mathrm M} := \widehat H_{\mathrm M}(U_k;S_k,(R_k,S'_k))-F_{\mathrm M}(U_k,S_k),
\end{equation}
then
\begin{equation}
\mathbb E\left[\Delta_k^{\mathrm M}\,\mid\,\mathcal G_k\right]=0
\end{equation}
and
\begin{equation}
\lVert \Delta_k^{\mathrm M} \rVert_{2,\infty} \le 2\beta_{\mathrm M}\lVert W_k \rVert_{2,\infty}+B_{\mathrm M}^\star \qquad\text{a.s.}
\end{equation}
\end{proposition}

\begin{proof}
For fixed $s,s'\in\X$ and $r\in[0,1]^q$, the proof of Proposition~\ref{prop:mtd-lipschitz} together with Proposition~\ref{prop:mtd-contraction-app} gives
\begin{equation}
\lVert \widehat T_{\mathrm M}(U;s,(r,s'))-\widehat T_{\mathrm M}(U';s,(r,s')) \rVert_2 \le \beta_{\mathrm M}\lVert U-U'\rVert_{2,\infty}
\end{equation}
and
\begin{equation}
\lVert (\mathcal O_{\mathrm M}U)(s)-(\mathcal O_{\mathrm M}U')(s) \rVert_2 \le \beta_{\mathrm M}\lVert U-U'\rVert_{2,\infty}.
\end{equation}
Because $[0,1]^q$ is compact, the state set is finite, and the map $r \mapsto \widehat T_{\mathrm M}(U^\star;s,(r,s'))$ is continuous for each fixed $(s,s')$, the constant $B_{\mathrm M}^\star$ is finite. Now add and subtract the same quantities at $U^\star$:
\begin{equation}
\begin{gathered}
\left\lVert \widehat T_{\mathrm M}(U;s,(r,s'))-(\mathcal O_{\mathrm M}U)(s) \right\rVert_2 \le \left\lVert \widehat T_{\mathrm M}(U;s,(r,s'))-\widehat T_{\mathrm M}(U^\star;s,(r,s')) \right\rVert_2 \\
+ \left\lVert \widehat T_{\mathrm M}(U^\star;s,(r,s'))-(\mathcal O_{\mathrm M}U^\star)(s) \right\rVert_2+ \left\lVert (\mathcal O_{\mathrm M}U^\star)(s)-(\mathcal O_{\mathrm M}U)(s) \right\rVert_2\\ 
\le 2\beta_{\mathrm M}\lVert U-U^\star \rVert_{2,\infty}+B_{\mathrm M}^\star,
\end{gathered}
\end{equation}
which is exactly the centered target-level affine bound. Squaring and using $(a+b)^2\le 2a^2+2b^2$ with
\begin{equation}
a:=2\beta_{\mathrm M}\lVert W_k \rVert_{2,\infty}, \qquad b:=B_{\mathrm M}^\star,
\end{equation}
gives the conditional second-moment bound of \eqref{eq:2ndmoment-mtd}. Since both coefficients are dominated by $A_{\mathrm M}^{\mathrm{iid}}$, this bound is at most $A_{\mathrm M}^{\mathrm{iid}}(1+\lVert W_k\rVert_{2,\infty}^2)$. Next, by definition of $\widehat H_{\mathrm M}$ and $F_{\mathrm M}$, only the $S_k$-th block can be nonzero, so
\begin{equation}
\lVert \Delta_k^{\mathrm M} \rVert_{2,\infty} = \left\lVert \widehat T_{\mathrm M}(U_k;S_k,(R_k,S'_k))-(\mathcal O_{\mathrm M}U_k)(S_k) \right\rVert_2,
\end{equation}
and the pathwise affine bound follows immediately. Finally, the conditional mean identity is exactly Proposition~\ref{prop:mtd-unbiased} written in centered form.
\end{proof}

\subsection{Completion of the proof of Theorem~\ref{thm:mtd-main}}

\begin{proof}[Proof of Theorem~\ref{thm:mtd-main}]
We verify the hypotheses of Appendix~\ref{app:common} for discounted MTD.

Proposition~\ref{prop:mtd-domain} verifies the restricted-domain hypotheses with $\mathcal D=\mathcal D_{\mathrm M}$. In particular, the sampled recursion remains in $\mathcal D_{\mathrm M}$ whenever $U_0\in\mathcal D_{\mathrm M}$ and the step sizes are at most one. All contraction, Lipschitz, and perturbation estimates below are therefore required only on $\mathcal D_{\mathrm M}$.

First, Proposition~\ref{prop:mtd-contraction-app} gives the block-supremum contraction with modulus
\begin{equation}
\beta_{\mathrm M} = \gamma^{c/2}.
\end{equation}

Second, Proposition~\ref{prop:mtd-lipschitz} gives the target-level estimate
\begin{equation}
\lVert \widehat T_{\mathrm M}(U;s,(r,s')) - \widehat T_{\mathrm M}(U';s,(r,s')) \rVert_2 \le \beta_{\mathrm M}\lVert U-U' \rVert_{2,\infty} \le \lVert U-U' \rVert_{2,\infty}.
\end{equation}

Third, Proposition~\ref{prop:mtd-perturbation} gives the centered pathwise affine perturbation bound
\begin{equation}
\mathbb{E}\left[\Delta_k^{\mathrm M} \mid \mathcal G_k\right] = 0, \qquad \lVert \Delta_k^{\mathrm M} \rVert_{2,\infty} \le 2\beta_{\mathrm M}\lVert W_k \rVert_{2,\infty} + B_{\mathrm M}^\star \qquad \text{a.s.}
\end{equation}
Hence the Markovian perturbation condition in Proposition~\ref{prop:abstract-markov} holds with
\begin{equation}
A_{3,\mathrm M}=2\beta_{\mathrm M},\qquad B_{3,\mathrm M}=B_{\mathrm M}^\star.
\end{equation}

For the i.i.d.\ case, Proposition~\ref{prop:mtd-perturbation} directly gives
\begin{equation}
\mathbb{E}\bigl[ \lVert \widehat T_{\mathrm M}(U_k;S_k,(R_k,S'_k)) - (\mathcal O_{\mathrm M}U_k)(S_k) \rVert_2^2 \,\mid\, U_k,S_k \bigr] \le A_{\mathrm M}^{\mathrm{iid}}\bigl(1+\lVert W_k\rVert_{2,\infty}^2\bigr),
\end{equation}
where
\begin{equation}
A_{\mathrm M}^{\mathrm{iid}}=8\beta_{\mathrm M}^2+2\bigl(B_{\mathrm M}^\star\bigr)^2.
\end{equation}
Therefore Proposition~\ref{prop:abstract-iid} applies without any further change of coordinates.

Let $q_{\X}:=\lvert\X\rvert^{2/p^\star}$ and choose
\begin{equation}
\vartheta_{\mathrm M,\rho}:=\frac{\rho_{\min}(1-\beta_{\mathrm M})}{q_{\X}}.
\end{equation}
Applying Proposition~\ref{prop:abstract-iid} with
\begin{equation}
\beta = \beta_{\mathrm M}, \qquad A^{\mathrm{iid}} = A_{\mathrm M}^{\mathrm{iid}}, \qquad \vartheta = \vartheta_{\mathrm M,\rho}, \qquad p^\star = \max\{2,\lceil \log \lvert \X \rvert \rceil\},
\end{equation}
yields Theorem~\ref{thm:mtd-main}(i) with
\begin{equation}
a_{\mathrm M,1}^{\mathrm{iid}} := \frac{1 + \vartheta_{\mathrm M,\rho}\lvert \X \rvert^{2/p^\star}}{1+\vartheta_{\mathrm M,\rho}}, \qquad a_{\mathrm M,2}^{\mathrm{iid}} := 1-\bar\beta_{\mathrm M,\rho}\sqrt{\frac{1 + \vartheta_{\mathrm M,\rho}\lvert \X \rvert^{2/p^\star}}{1+\vartheta_{\mathrm M,\rho}}},
\end{equation}
where
\begin{equation}
\bar\beta_{\mathrm M,\rho} := 1-\rho_{\min}(1-\beta_{\mathrm M}),
\end{equation}
and
\begin{equation}
\begin{gathered}
a_{\mathrm M,3}^{\mathrm{iid}} := \frac{4(p^\star-1)\lvert \X \rvert^{2/p^\star}(A_{\mathrm M}^{\mathrm{iid}}+2)(1+\vartheta_{\mathrm M,\rho}\lvert \X \rvert^{2/p^\star})}{\vartheta_{\mathrm M,\rho}},\\
a_{\mathrm M,4}^{\mathrm{iid}} := \frac{2(p^\star-1)\lvert \X \rvert^{2/p^\star}A_{\mathrm M}^{\mathrm{iid}}(1+\vartheta_{\mathrm M,\rho}\lvert \X \rvert^{2/p^\star})}{\vartheta_{\mathrm M,\rho}}.
\end{gathered}
\end{equation}
The compact constants and thresholds in Theorem~\ref{thm:mtd-main}(i) can therefore be chosen as
\begin{equation}
\begin{gathered}
c_{\mathrm M}^{\mathrm{iid}}:=a_{\mathrm M,2}^{\mathrm{iid}},\qquad
\bar\alpha_{\mathrm M}^{\mathrm{iid}}:=\min\left\{1,\frac{a_{\mathrm M,2}^{\mathrm{iid}}}{a_{\mathrm M,3}^{\mathrm{iid}}}\right\},\\
C_{\mathrm M}^{\mathrm{iid}}:=\max\left\{a_{\mathrm M,1}^{\mathrm{iid}},4e a_{\mathrm M,4}^{\mathrm{iid}},\frac{2a_{\mathrm M,4}^{\mathrm{iid}}}{a_{\mathrm M,2}^{\mathrm{iid}}}\right\},\\
h_{\mathrm M}^{\mathrm{iid}}(\alpha):=\left\lceil\max\left\{1,\alpha,\frac{\alpha a_{\mathrm M,3}^{\mathrm{iid}}}{a_{\mathrm M,2}^{\mathrm{iid}}}\right\}\right\rceil,\\
h_{\mathrm M}^{\mathrm{iid}}(\alpha,z):=\left\lceil\max\left\{1,\alpha^{1/z},\left(\frac{\alpha a_{\mathrm M,3}^{\mathrm{iid}}}{a_{\mathrm M,2}^{\mathrm{iid}}}\right)^{1/z},\left(\frac{2z}{a_{\mathrm M,2}^{\mathrm{iid}}\alpha}\right)^{1/(1-z)}\right\}\right\rceil.
\end{gathered}
\end{equation}

For the Markovian case, Proposition~\ref{prop:abstract-markov} applies with
\begin{equation}
A_{3,\mathrm M}=2\beta_{\mathrm M},\qquad B_{3,\mathrm M}=B_{\mathrm M}^\star,\qquad \mathsf A_{\mathrm M}=(3+2\beta_{\mathrm M})^2,\qquad \mathsf B_{\mathrm M}=2(B_{\mathrm M}^\star)^2.
\end{equation}
Choose
\begin{equation}
\vartheta_{\mathrm M,\mu}:=\frac{\mu_{\min}(1-\beta_{\mathrm M})}{q_{\X}},
\end{equation}
and let $\eta_{\mathrm M,\mu}$, $K_{\mathrm M,0}$, and $K_{\mathrm M,1}$ be the explicit constants in Proposition~\ref{prop:abstract-markov} under these substitutions. The theorem drift constant and admissible constant step size can be taken as
\begin{equation}
c_{\mathrm M}^{\mathrm{mk}}:=\frac{\eta_{\mathrm M,\mu}}{2},\qquad \bar\alpha_{\mathrm M}^{\mathrm{mk}}:=\min\left\{1,\frac{\eta_{\mathrm M,\mu}}{\mathsf A_{\mathrm M}(5+2\eta_{\mathrm M,\mu})K_{\mathrm M,1}}\right\}.
\end{equation}
For the diminishing schedules, one may take
\begin{equation}
h_{\mathrm M}^{\mathrm{mk}}(\alpha):=\left\lceil\max\left\{\mathsf A_{\mathrm M}\alpha(5\alpha+8)K_{\mathrm M,1},2\right\}\right\rceil
\end{equation}
and, with $\chi_{\mathrm M,\mu}:=z/\alpha+\eta_{\mathrm M,\mu}$,
\begin{equation}
h_{\mathrm M}^{\mathrm{mk}}(\alpha,z):=\left\lceil\max\left\{\left(\frac{2\mathsf A_{\mathrm M}\alpha(5+2\chi_{\mathrm M,\mu})K_{\mathrm M,1}}{\eta_{\mathrm M,\mu}}\right)^{1/z},2\right\}\right\rceil.
\end{equation}
Set $C_{\mathrm M}^{\mathrm{mk}}:=C_\mu^{\mathrm{mk}}$ with $\mathsf A=\mathsf A_{\mathrm M}$, $\mathsf B=\mathsf B_{\mathrm M}$, $\eta_\mu=\eta_{\mathrm M,\mu}$, and $K_{\mu,1}=K_{\mathrm M,1}$ in the explicit definition following Proposition~\ref{prop:abstract-markov}. This coefficient is independent of $k$ and of the step size.

The constant, linearly-diminishing, and polynomially-diminishing step size bounds in Theorem~\ref{thm:mtd-main} now follow from the corresponding specializations of Proposition~\ref{prop:abstract-iid} and Proposition~\ref{prop:abstract-markov}.
\end{proof}

\subsection{Proof of Corollary~\ref{cor:disc-mtd}}

\begin{proof}
We use the linearly-diminishing step size bounds in Theorem~\ref{thm:mtd-main}.

For part \textup{(i)}, the linearly-diminishing i.i.d.\ bound in Theorem~\ref{thm:mtd-main}(i) gives
\begin{equation}
\mathbb E\left[\ell_{\mathrm M,\infty}(\eta_k,\eta^\star)^2\right] \le a_{\mathrm M,1}^{\mathrm{iid}}\ell_{\mathrm M,\infty}(\eta_0,\eta^\star)^2 \left(\frac{h}{k+h}\right)^{a_{\mathrm M,2}^{\mathrm{iid}}\alpha} + \frac{4e\alpha^2 a_{\mathrm M,4}^{\mathrm{iid}}}{a_{\mathrm M,2}^{\mathrm{iid}}\alpha-1}\cdot \frac{1}{k+h}.
\end{equation}
Since $\alpha>1/a_{\mathrm M,2}^{\mathrm{iid}}$, the first term is also $O((k+h)^{-1})$. Therefore
\begin{equation}
\mathbb E\left[\ell_{\mathrm M,\infty}(\eta_k,\eta^\star)^2\right] = O\left(\frac{1}{k+h}\right),
\end{equation}
and hence $\mathbb E[\ell_{\mathrm M,\infty}(\eta_k,\eta^\star)]\le\varepsilon$ is guaranteed for $k=O(\varepsilon^{-2})$.

For part \textup{(ii)}, the linearly-diminishing Markovian bound in Theorem~\ref{thm:mtd-main}(ii) gives
\begin{equation}
\mathbb E\left[\ell_{\mathrm M,\infty}(\eta_k,\eta^\star)^2\right] \le C_{\mathrm M}^{\mathrm{mk}}\bigl(1+\ell_{\mathrm M,\infty}(\eta_0,\eta^\star)^2\bigr)\left(\frac{h}{k+h}\right)^{c_{\mathrm M}^{\mathrm{mk}}\alpha}+\frac{C_{\mathrm M}^{\mathrm{mk}}\alpha^2}{c_{\mathrm M}^{\mathrm{mk}}\alpha-1}\cdot\frac{1}{k+h}.
\end{equation}
Since $\alpha>1/c_{\mathrm M}^{\mathrm{mk}}$, the first term is $O((k+h)^{-1})$. Therefore
\begin{equation}
\mathbb E\left[\ell_{\mathrm M,\infty}(\eta_k,\eta^\star)^2\right]=O\left(\frac{1}{k+h}\right),
\end{equation}
which implies $\mathbb E[\ell_{\mathrm M,\infty}(\eta_k,\eta^\star)]\le \varepsilon$ for $k=O(\varepsilon^{-2})$.
\end{proof}

\section{Undiscounted fixed-horizon CTD}
\label{app:fh-ctd}

This appendix records the undiscounted fixed-horizon CTD ingredients culminating in the proof of Theorem~\ref{thm:undisc-ctd}.

\subsection{Flattened weighted horizon-state space}

To apply the framework of Appendix~\ref{app:common}, we rewrite the weighted fixed-horizon recursion on a flattened product space indexed by horizon-state pairs. Let
\begin{equation}
\mathcal S_H := \{(h,s): h\in\{1,\dots,H\},\ s\in \X\}, \qquad \lvert \mathcal S_H \rvert = H\lvert \X \rvert.
\end{equation}
Let
\begin{equation}
V_{H, \mathrm C} := \prod_{(h, s) \in \mathcal S_H} \R^d
\end{equation}
denote the corresponding flattened Euclidean product space. For a horizon-stacked iterate $U=(U^h(s))_{1\le h\le H,\ s\in \X}$, define the weighted flattening
\begin{equation}
\bar{U}(h,s):=\lambda^h U^h(s), \qquad (h,s)\in\mathcal S_H.
\end{equation}
Then
\begin{equation}
\lVert U \rVert_{H,2,\infty} = \max_{(h,s)\in\mathcal S_H}\lVert \bar{U}(h,s) \rVert_2, \qquad \lVert U \rVert_{H,2,p} = \left( \sum_{(h,s)\in\mathcal S_H}\lVert \bar{U}(h,s) \rVert_2^p \right)^{1/p}.
\end{equation}
For a single horizon stack $U=(U^1,\dots,U^H)$, also define
\begin{equation}
\lVert U \rVert_{H,2} := \max_{1\le h\le H}\lambda^h\lVert U^h \rVert_2.
\end{equation}
Then
\begin{equation}
\lVert U \rVert_{H,2,\infty} = \max_{s\in \X}\lVert U(s) \rVert_{H,2}.
\end{equation}
Thus the fixed-horizon weighted norm is an ordinary block-supremum norm on a finite product of Euclidean blocks, and Appendix~\ref{app:common} applies with $\lvert \mathcal S_H \rvert$ in place of $\lvert \X \rvert$. In particular, throughout this appendix we set
\begin{equation}
p_H^\star := \max\{2,\lceil \log \lvert \mathcal S_H \rvert\rceil\}.
\end{equation}

\subsection{Fixed-horizon CTD residual map}

For each state $s\in \X$, let
\begin{equation}
U(s):=\bigl(U^1(s),\dots,U^H(s)\bigr)
\end{equation}
denote the horizon stack at state $s$, and let $P_s^H$ denote the coordinate projector onto that stack. Define
\begin{equation}
F_{H,\mathrm C}(U;s,(r,s')) := P_s^H\bigl(\widehat T_{H,\mathrm C}(U;s,(r,s'))-U(s)\bigr),
\end{equation}
where $\widehat T_{H,\mathrm C}(U;s,(r,s'))$ is the full horizon-stacked CTD target from Section~\ref{sec:undisc}. Then the online fixed-horizon CTD recursion is
\begin{equation}
U_{k+1}=U_k+\alpha_k F_{H,\mathrm C}(U_k;S_k,(R_k,S_{k+1})).
\end{equation}
Define also the CTD embedded averaged operator by
\begin{equation}
\mathcal O_{H,\mathrm C} := I_{H,\mathrm C}\circ \Pi_{H,\mathrm C}^{\Theta} T_H^\pi \circ I_{H,\mathrm C}^{-1}.
\end{equation}
Define
\begin{equation}
\mathcal D_{H,\mathrm C}:=I_{H,\mathrm C}(\mathcal F_{H,\mathrm C,\Theta}^{\X}).
\end{equation}
This is a compact convex subset of $V_{H,\mathrm C}$. Both $\mathcal O_{H,\mathrm C}$ and the full-sample replacement $U+F_{H,\mathrm C}(U;s,(r,s'))$ map $\mathcal D_{H,\mathrm C}$ into itself. Hence Lemma~\ref{lem:invariant-domain} shows that the online recursion remains in $\mathcal D_{H,\mathrm C}$ whenever $U_0\in\mathcal D_{H,\mathrm C}$ and $0\le\alpha_k\le1$. All fixed-horizon CTD operator estimates below are imposed only on this invariant domain.

\begin{proposition}[Fixed-horizon CTD weighted contraction]
\label{prop:fh-ctd-contract-app}
For all $U,U' \in I_{H,\mathrm C}(\mathcal F_{H,\mathrm C,\Theta}^\X)$,
\begin{equation}
\lVert \mathcal O_{H,\mathrm C}U-\mathcal O_{H,\mathrm C}U'\rVert_{H,2,\infty} \le \lambda \lVert U-U'\rVert_{H,2,\infty}.
\end{equation}
\end{proposition}

\begin{proof}
Let $\eta:=I_{H,\mathrm C}^{-1}(U)$ and $\eta':=I_{H,\mathrm C}^{-1}(U')$. Fix $(h,s)\in\mathcal S_H$ with $h \ge 1$. By the statewise CTD isometry and nonexpansiveness of the categorical projection in the Cram\'er metric,
\begin{equation}
\begin{aligned}
\lambda^h \lVert (\mathcal O_{H,\mathrm C}U)^h(s)-(\mathcal O_{H,\mathrm C}U')^h(s) \rVert_2 &= \lambda^h \ell_{\mathrm C}\left( (\Pi_{H,\mathrm C}^{\Theta}T_H^\pi\eta)^h(s), (\Pi_{H,\mathrm C}^{\Theta}T_H^\pi\eta')^h(s) \right) \\ 
&\le \lambda^h \ell_{\mathrm C}\left( (T_H^\pi\eta)^h(s), (T_H^\pi\eta')^h(s) \right).
\end{aligned}
\end{equation}
Now
\begin{equation}
(T_H^\pi\eta)^h(s) = \sum_{a\in\A}\pi(a \mid s)\sum_{s'\in\X}P(s' \mid s,a)\, (f_{R(s,a),1})_\#\eta^{h-1}(s'),
\end{equation}
and the same formula holds for $\eta'$. Since translation preserves the Cram\'er metric and the Cram\'er metric is convex under mixtures,
\begin{equation}
\begin{aligned}
\ell_{\mathrm C}\left( (T_H^\pi\eta)^h(s), (T_H^\pi\eta')^h(s) \right) &\le \sum_{a\in\A}\pi(a \mid s)\sum_{s'\in\X}P(s' \mid s,a)\, \ell_{\mathrm C}\left(\eta^{h-1}(s'),{\eta'}^{h-1}(s')\right) \\ 
&\le \max_{x\in\X} \ell_{\mathrm C}\left(\eta^{h-1}(x),{\eta'}^{h-1}(x)\right).
\end{aligned}
\end{equation}
Therefore
\begin{equation}
\begin{aligned}
\lambda^h \lVert (\mathcal O_{H,\mathrm C}U)^h(s)-(\mathcal O_{H,\mathrm C}U')^h(s) \rVert_2 &\le \lambda \max_{x\in\X} \lambda^{h-1} \ell_{\mathrm C}\left(\eta^{h-1}(x),{\eta'}^{h-1}(x)\right) \\ 
&\le \lambda \ell_{H,\mathrm C,\infty}(\eta,\eta') = \lambda \lVert U-U'\rVert_{H,2,\infty}.
\end{aligned}
\end{equation}
Taking the maximum over $(h,s)\in\mathcal S_H$ proves the claim.
\end{proof}

\subsection{One-step bounds}

\begin{proposition}[Fixed-horizon CTD samplewise Lipschitz continuity]
\label{prop:fh-ctd-lipschitz}
For every $s \in \X$, $(r, s') \in [0,1] \times \X$ and $U, U' \in I_{H,\mathrm C}(\mathcal F_{H,\mathrm C,\Theta}^\X)$,
\begin{equation}
\lVert F_{H,\mathrm C}(U;s,(r,s'))-F_{H,\mathrm C}(U';s,(r,s')) \rVert_{H,2,\infty} \le 2\lVert U-U'\rVert_{H,2,\infty}.
\end{equation}
\end{proposition}

\begin{proof}
For $h=1$, the local target $\widehat T_{H,\mathrm C}^h(U;s,(r,s'))$ is independent of $U$ because the layer-zero law is fixed. For $h\ge2$, it depends only on the $(h-1)$st horizon block at $s'$. The Cram\'er isometry, nonexpansiveness of the categorical projection, and translation invariance of the Cram\'er metric show directly that this local target is $1$-Lipschitz in the Euclidean block norm. Multiplying by $\lambda^h$ gives
\begin{equation}
\resizebox{\linewidth}{!}{$
\lambda^h \lVert \widehat T_{H,\mathrm C}^h(U;s,(r,s'))-\widehat T_{H,\mathrm C}^h(U';s,(r,s')) \rVert_2 \le \lambda^{h-1}\lVert U^{h-1}(s')-U'^{\,h-1}(s') \rVert_2 \le \lVert U-U'\rVert_{H,2,\infty}.
$}
\end{equation}
Taking the maximum over $h$ yields
\begin{equation}
\lVert \widehat T_{H,\mathrm C}(U;s,(r,s'))-\widehat T_{H,\mathrm C}(U';s,(r,s')) \rVert_{H,2} \le \lVert U-U'\rVert_{H,2,\infty}.
\end{equation}
Since $F_{H,\mathrm C}$ is the projected residual
\begin{equation}
F_{H,\mathrm C}(U;s,(r,s')) = P_s^H\bigl(\widehat T_{H,\mathrm C}(U;s,(r,s')) - U(s)\bigr),
\end{equation}
the difference $F_{H,\mathrm C}(U)-F_{H,\mathrm C}(U')$ splits into a projected target difference and a projected current-stack difference. Each is bounded by $\lVert U-U' \rVert_{H,2,\infty}$, yielding the factor $2$ through triangle inequality.
\end{proof}

\begin{proposition}[Fixed-horizon CTD pathwise boundedness]
\label{prop:fh-ctd-bounded}
For every $s \in \X$, $(r, s') \in [0,1] \times \X$ and $U \in I_{H,\mathrm C}(\mathcal F_{H,\mathrm C,\Theta}^\X)$,
\begin{equation}
\lVert F_{H,\mathrm C}(U;s,(r,s')) \rVert_{H,2,\infty} \le 2B_{H,\mathrm C},
\end{equation}
where
\begin{equation}
B_{H,\mathrm C}:= \max_{1\le h\le H,\ s\in \X} \lambda^h\sqrt{\theta_{h,d}(s)-\theta_{h,1}(s)}.
\end{equation}
\end{proposition}

\begin{proof}
Fix $h$ and $s$. Both the local target $\widehat T_{H,\mathrm C}^h(U;s,(r,s'))$ and the current iterate $U^h(s)$ are embedded categorical laws supported on $\Theta_h(s)$. By the statewise CTD support radius bound, each has Euclidean norm at most $\sqrt{\theta_{h,d}(s)-\theta_{h,1}(s)}$. After weighting by $\lambda^h$, both are bounded by $B_{H,\mathrm C}$. The triangle inequality then yields the claim.
\end{proof}

\subsection{Phasewise averaged maps}
Recall the definition of phase distributions from Section~\ref{sec:undisc}:
\begin{equation}
\rho_t(s):=\Pr(S_t=s), \qquad 0\le t\le H-1,\ s\in \X.
\end{equation}
For each phase $t\in\{0,\dots,H-1\}$, define
\begin{equation}
\Gamma_{t,\mathrm C}(U) := \sum_{s\in \X}\rho_t(s)\, P_s^H\bigl((\mathcal O_{H,\mathrm C}U)(s)-U(s)\bigr), \qquad G_{t,\mathrm C}(U):=U+\Gamma_{t,\mathrm C}(U).
\end{equation}

\begin{proposition}[Episode-averaged CTD contraction]
\label{prop:fh-ctd-avg-contract}
Suppose $\bar\alpha_m>0$ and define the step-size-weighted episode distribution
\begin{equation}
\rho_m^\alpha(s):=\frac{1}{\bar\alpha_m}\sum_{t=0}^{H-1}\alpha_{mH+t}\rho_t(s).
\end{equation}
If the step sizes are nonincreasing within episode $m$ and satisfy
\begin{equation}
\alpha_{mH+H-1}\ge\frac{1}{2}\alpha_{mH},
\label{eq:fh-step-regularity}
\end{equation}
then
\begin{equation}
\min_{s\in\X}\rho_m^\alpha(s)\ge\frac{\rho_{H,\min}}{2}.
\end{equation}
Define
\begin{equation}
G_{m,\mathrm C}^\alpha(U):=U+\sum_{s\in\X}\rho_m^\alpha(s)P_s^H\bigl((\mathcal O_{H,\mathrm C}U)(s)-U(s)\bigr).
\end{equation}
Then
\begin{equation}
\lVert G_{m,\mathrm C}^\alpha(U)-G_{m,\mathrm C}^\alpha(U')\rVert_{H,2,\infty}\le\bar\beta_{H}^{\mathrm{occ}}\lVert U-U'\rVert_{H,2,\infty},
\end{equation}
where
\begin{equation}
\bar\beta_{H}^{\mathrm{occ}}:=1-\frac{\rho_{H,\min}}{2}(1-\lambda),\qquad \kappa_H^{\mathrm{occ}}:=1-\bar\beta_H^{\mathrm{occ}}=\frac{\rho_{H,\min}}{2(H+1)}.
\end{equation}
Moreover, the frozen episode map
\begin{equation}
A_{m,\mathrm C}^{\mathrm{fr}}(U):=U+\sum_{t=0}^{H-1}\alpha_{mH+t}\Gamma_{t,\mathrm C}(U)=(1-\bar\alpha_m)U+\bar\alpha_mG_{m,\mathrm C}^\alpha(U)
\end{equation}
fixes $U_{H,\mathrm C}^\star$ and is a $(1-\kappa_H^{\mathrm{occ}}\bar\alpha_m)$-contraction whenever $\bar\alpha_m\le1$.
\end{proposition}

\begin{proof}
The denominator in $\rho_m^\alpha$ satisfies $\bar\alpha_m\le H\alpha_{mH}$, while its numerator satisfies
\begin{equation}
\sum_{t=0}^{H-1}\alpha_{mH+t}\rho_t(s)\ge\alpha_{mH+H-1}\sum_{t=0}^{H-1}\rho_t(s)\ge\frac{H\alpha_{mH}}{2}\bar\rho_H(s).
\end{equation}
This proves the lower bound on $\rho_m^\alpha$. For each state block $s$,
\begin{equation}
\begin{gathered}
\left\lVert\bigl(G_{m,\mathrm C}^\alpha(U)-G_{m,\mathrm C}^\alpha(U')\bigr)(s)\right\rVert_{H,2}\\
\le\bigl(1-\rho_m^\alpha(s)+\lambda\rho_m^\alpha(s)\bigr)\lVert U-U'\rVert_{H,2,\infty}\le\bar\beta_H^{\mathrm{occ}}\lVert U-U'\rVert_{H,2,\infty},
\end{gathered}
\end{equation}
where Proposition~\ref{prop:fh-ctd-contract-app} gives the first inequality. Taking the maximum over $s$ proves the contraction of $G_{m,\mathrm C}^\alpha$. The identity for $A_{m,\mathrm C}^{\mathrm{fr}}$ follows from the definition of $\rho_m^\alpha$. Convexity of the norm then gives its contraction factor, and $\Gamma_{t,\mathrm C}(U_{H,\mathrm C}^\star)=0$ for every phase proves the fixed-point claim.
\end{proof}

\subsection{Auxiliary phasewise bounds}

For episode $m$, write
\begin{equation}
U_{m,t}:=U_{mH+t},\qquad t=0,\dots,H.
\end{equation}
Also define the phasewise filtration
\begin{equation}
\mathcal F_{m,t} := \sigma\bigl( U_{m,0}, (S_u^m,R_u^m,S_{u+1}^m)_{0 \le u < t} \bigr), \qquad 0 \le t \le H,
\end{equation}
where the superscript $m$ indicates transitions within episode $m$. Then $U_{m,t}$ is $\mathcal F_{m,t}$-measurable for every $t$.

\begin{lemma}[Within-episode deviation]
\label{lem:fh-ctd-within}
For all $m$ and all $t\in\{0,\dots,H\}$,
\begin{equation}
\lVert U_{m,t}-U_{m,0} \rVert_{H,2,\infty} \le 2B_{H,\mathrm C}\sum_{u=0}^{t-1}\alpha_{mH+u}.
\end{equation}
\end{lemma}

\begin{proof}
From the recursion,
\begin{equation}
U_{m,t}-U_{m,0} = \sum_{u=0}^{t-1} \alpha_{mH+u} F_{H,\mathrm C}\bigl(U_{m,u};S_u^m,(R_u^m,S_{u+1}^m)\bigr).
\end{equation}
Taking $\lVert \cdot \rVert_{H,2,\infty}$ and using Proposition~\ref{prop:fh-ctd-bounded} term by term gives
\begin{equation}
\lVert U_{m,t}-U_{m,0} \rVert_{H,2,\infty} \le \sum_{u=0}^{t-1}\alpha_{mH+u}\, \lVert F_{H,\mathrm C}(U_{m,u};S_u^m,(R_u^m,S_{u+1}^m)) \rVert_{H,2,\infty} \le 2B_{H,\mathrm C}\sum_{u=0}^{t-1}\alpha_{mH+u}.
\end{equation}
\end{proof}

\begin{lemma}[Frozen residual decomposition]
\label{lem:fh-ctd-bias}
For each phase $t \in \{0,\dots,H-1\}$, define
\begin{equation}
\zeta_{m,t}^{\mathrm C} := F_{H,\mathrm C}(U_{m,0};S_t^m,(R_t^m,S_{t+1}^m)) - \Gamma_{t,\mathrm C}(U_{m,0})
\end{equation}
and
\begin{equation}
\chi_{m,t}^{\mathrm C} := F_{H,\mathrm C}(U_{m,t};S_t^m,(R_t^m,S_{t+1}^m)) - F_{H,\mathrm C}(U_{m,0};S_t^m,(R_t^m,S_{t+1}^m)).
\end{equation}
Then
\begin{equation}
\mathbb E\left[\zeta_{m,t}^{\mathrm C}\,\mid\,U_{m,0}\right]=0, \qquad \lVert \zeta_{m,t}^{\mathrm C} \rVert_{H,2,\infty} \le 4B_{H,\mathrm C},
\end{equation}
and
\begin{equation}
\lVert \chi_{m,t}^{\mathrm C} \rVert_{H,2,\infty} \le 2\lVert U_{m,t}-U_{m,0} \rVert_{H,2,\infty} \le 4B_{H,\mathrm C}\bar\alpha_m.
\end{equation}
If
\begin{equation}
D_m^{\mathrm C}:=U_{m,H}-A_{m,\mathrm C}^{\mathrm{fr}}(U_{m,0}),
\end{equation}
then
\begin{equation}
D_m^{\mathrm C}=\sum_{t=0}^{H-1}\alpha_{mH+t}\bigl(\zeta_{m,t}^{\mathrm C}+\chi_{m,t}^{\mathrm C}\bigr).
\end{equation}
Consequently, whenever $\bar\alpha_m\le1$,
\begin{equation}
\left\lVert\mathbb E\left[D_m^{\mathrm C}\mid U_{m,0}\right]\right\rVert_{H,2,\infty}\le4B_{H,\mathrm C}\bar\alpha_m^2,
\end{equation}
and
\begin{equation}
\lVert D_m^{\mathrm C}\rVert_{H,2,\infty}\le8B_{H,\mathrm C}\bar\alpha_m
\end{equation}
pathwise.
\end{lemma}

\begin{proof}
Since episode $m$ is independent of $U_{m,0}$ and its phase-$t$ sample has marginal law $\rho_t$,
\begin{equation}
\mathbb E\left[\zeta_{m,t}^{\mathrm C}\,\mid\,U_{m,0}\right]=0.
\end{equation}
\begin{equation}
\lVert \zeta_{m,t}^{\mathrm C} \rVert_{H,2,\infty} \le \lVert F_{H,\mathrm C}(U_{m,0};S_t^m,(R_t^m,S_{t+1}^m)) \rVert_{H,2,\infty} + \lVert \Gamma_{t,\mathrm C}(U_{m,0}) \rVert_{H,2,\infty}.
\end{equation}
By Proposition~\ref{prop:fh-ctd-bounded},
\begin{equation}
\lVert F_{H,\mathrm C}(U_{m,0};S_t^m,(R_t^m,S_{t+1}^m)) \rVert_{H,2,\infty} \le 2B_{H,\mathrm C},
\end{equation}
and, since $\Gamma_{t,\mathrm C}(U_{m,0})$ is a convex combination of such residuals,
\begin{equation}
\lVert \Gamma_{t,\mathrm C}(U_{m,0}) \rVert_{H,2,\infty} \le 2B_{H,\mathrm C}.
\end{equation}
This gives $\lVert \zeta_{m,t}^{\mathrm C} \rVert_{H,2,\infty} \le 4B_{H,\mathrm C}$.

Proposition~\ref{prop:fh-ctd-lipschitz} gives
\begin{equation}
\left\lVert F_{H,\mathrm C}(U;s,(r,s')) - F_{H,\mathrm C}(U';s,(r,s')) \right\rVert_{H,2,\infty} \le 2\lVert U-U' \rVert_{H,2,\infty}
\end{equation}
for every admissible $(s,r,s')$. Lemma~\ref{lem:fh-ctd-within} therefore gives
\begin{equation}
\lVert\chi_{m,t}^{\mathrm C}\rVert_{H,2,\infty}\le2\lVert U_{m,t}-U_{m,0}\rVert_{H,2,\infty}\le4B_{H,\mathrm C}\bar\alpha_m.
\end{equation}
Summing the online recursion through the episode and subtracting the definition of $A_{m,\mathrm C}^{\mathrm{fr}}$ gives the displayed identity for $D_m^{\mathrm C}$. The frozen terms have conditional mean zero separately, so
\begin{equation}
\left\lVert\mathbb E\left[D_m^{\mathrm C}\mid U_{m,0}\right]\right\rVert_{H,2,\infty}\le\sum_{t=0}^{H-1}\alpha_{mH+t}\mathbb E\left[\lVert\chi_{m,t}^{\mathrm C}\rVert_{H,2,\infty}\mid U_{m,0}\right]\le4B_{H,\mathrm C}\bar\alpha_m^2.
\end{equation}
Finally, the pathwise bounds on $\zeta_{m,t}^{\mathrm C}$ and $\chi_{m,t}^{\mathrm C}$ give
\begin{equation}
\lVert D_m^{\mathrm C}\rVert_{H,2,\infty}\le4B_{H,\mathrm C}\bar\alpha_m+4B_{H,\mathrm C}\bar\alpha_m^2\le8B_{H,\mathrm C}\bar\alpha_m.
\end{equation}
\end{proof}

\subsection{Episode-level finite-iteration drift}
\label{sec:fh-ctd-drift-app}


Let
\begin{equation}
W_{m,0}:=U_{m,0}-U_{H,\mathrm C}^\star,\qquad q_H:=\lvert\mathcal S_H\rvert^{2/p_H^\star},\qquad \kappa_H:=\kappa_H^{\mathrm{occ}}=\frac{\rho_{H,\min}}{2(H+1)}.
\end{equation}
Choose the explicit smoothing parameter
\begin{equation}
\vartheta_{H,\mathrm C}:=\frac{\kappa_H}{q_H}
\end{equation}
and define
\begin{equation}
\begin{gathered}
r_{H,\mathrm C}^{\mathrm{fh}}:=\frac{1+\vartheta_{H,\mathrm C}q_H}{1+\vartheta_{H,\mathrm C}},\qquad L_{H,\mathrm C}^{\mathrm{fh}}:=\frac{p_H^\star-1}{\vartheta_{H,\mathrm C}},\\
\omega_{H,\mathrm C}^{\mathrm{fh}}:=1-(1-\kappa_H)\sqrt{r_{H,\mathrm C}^{\mathrm{fh}}},\qquad d_{H,\mathrm C}^{\mathrm{fh}}:=8L_{H,\mathrm C}^{\mathrm{fh}}q_H(1+\vartheta_{H,\mathrm C}q_H).
\end{gathered}
\end{equation}
The same calculation as in Proposition~\ref{prop:explicit-discounted-constants} gives
\begin{equation}
r_{H,\mathrm C}^{\mathrm{fh}}\le2,\qquad \omega_{H,\mathrm C}^{\mathrm{fh}}\ge\frac{\kappa_H}{2},\qquad d_{H,\mathrm C}^{\mathrm{fh}}\le\frac{16(p_H^\star-1)q_H^2}{\kappa_H}.
\end{equation}
Define
\begin{equation}
\bar\alpha_{H,\mathrm C}:=\min\left\{\frac{1}{H},\frac{\omega_{H,\mathrm C}^{\mathrm{fh}}}{Hd_{H,\mathrm C}^{\mathrm{fh}}}\right\},
\end{equation}
\begin{equation}
c_{H,\mathrm C,1}^{\mathrm{fh}}:=\frac{\omega_{H,\mathrm C}^{\mathrm{fh}}}{2},\qquad c_{H,\mathrm C,2}^{\mathrm{fh}}:=32HL_{H,\mathrm C}^{\mathrm{fh}}q_HB_{H,\mathrm C}^2\left(1+\frac{1}{\omega_{H,\mathrm C}^{\mathrm{fh}}}\right).
\end{equation}

\begin{proposition}[Episodewise CTD potential drift]
\label{prop:fh-ctd-drift}
Suppose that $(\alpha_k)_{k\ge0}$ is nonincreasing, $\alpha_0\le\bar\alpha_{H,\mathrm C}$, and the within-episode regularity condition \eqref{eq:fh-step-regularity} holds. Then
\begin{equation}
\mathbb E\left[M_{H,\mathrm C}(W_{m+1,0})\mid U_{m,0}\right]\le\left(1-c_{H,\mathrm C,1}^{\mathrm{fh}}\bar\alpha_m\right)M_{H,\mathrm C}(W_{m,0})+c_{H,\mathrm C,2}^{\mathrm{fh}}\bar\alpha_m^{(2)},
\end{equation}
where
\begin{equation}
M_{H,\mathrm C}(W):=\inf_{Z\in V_{H,\mathrm C}}\left\{\frac{1}{2}\lVert Z\rVert_{H,2,\infty}^2+\frac{1}{2\vartheta_{H,\mathrm C}q_H}\lVert W-Z\rVert_{H,2,p_H^\star}^2\right\}.
\end{equation}
Moreover,
\begin{equation}
c_{H,\mathrm C,1}^{\mathrm{fh}}\ge\frac{\kappa_H}{4},\qquad c_{H,\mathrm C,2}^{\mathrm{fh}}\le\frac{96H(p_H^\star-1)q_H^2B_{H,\mathrm C}^2}{\kappa_H^2}.
\end{equation}
\end{proposition}

\begin{proof}
Set
\begin{equation}
\widetilde U_{m,H}:=A_{m,\mathrm C}^{\mathrm{fr}}(U_{m,0}),\qquad \widetilde W_{m,H}:=\widetilde U_{m,H}-U_{H,\mathrm C}^\star.
\end{equation}
Proposition~\ref{prop:fh-ctd-avg-contract} writes the frozen map as a $\bar\alpha_m$-relaxation of a $(1-\kappa_H)$-contraction. The Moreau-envelope smoothness inequality therefore gives
\begin{equation}
M_{H,\mathrm C}(\widetilde W_{m,H})\le\left(1-2\omega_{H,\mathrm C}^{\mathrm{fh}}\bar\alpha_m+d_{H,\mathrm C}^{\mathrm{fh}}\bar\alpha_m^2\right)M_{H,\mathrm C}(W_{m,0}).
\end{equation}
Indeed, the envelope defines a squared smooth norm, so the contraction and the envelope comparison give
\begin{equation}
\left\langle\nabla M_{H,\mathrm C}(W),G_{m,\mathrm C}^\alpha(U_{H,\mathrm C}^\star+W)-U_{H,\mathrm C}^\star-W\right\rangle\le-2\omega_{H,\mathrm C}^{\mathrm{fh}}M_{H,\mathrm C}(W),
\end{equation}
while the quadratic smoothness remainder is at most $d_{H,\mathrm C}^{\mathrm{fh}}\bar\alpha_m^2M_{H,\mathrm C}(W)$. Since $\bar\alpha_m\le H\alpha_0\le\omega_{H,\mathrm C}^{\mathrm{fh}}/d_{H,\mathrm C}^{\mathrm{fh}}$,
\begin{equation}
M_{H,\mathrm C}(\widetilde W_{m,H})\le\left(1-\omega_{H,\mathrm C}^{\mathrm{fh}}\bar\alpha_m\right)M_{H,\mathrm C}(W_{m,0}).
\end{equation}

By Lemma~\ref{lem:fh-ctd-bias}, $W_{m+1,0}=\widetilde W_{m,H}+D_m^{\mathrm C}$ and
\begin{equation}
\left\lVert\mathbb E[D_m^{\mathrm C}\mid U_{m,0}]\right\rVert_{H,2,\infty}\le4B_{H,\mathrm C}\bar\alpha_m^2,\qquad \mathbb E\left[\lVert D_m^{\mathrm C}\rVert_{H,2,\infty}^2\mid U_{m,0}\right]\le64B_{H,\mathrm C}^2\bar\alpha_m^2.
\end{equation}
Smoothness at $\widetilde W_{m,H}$, the self-bounding estimate $\lVert\nabla M(W)\rVert_*^2\le2L_{H,\mathrm C}^{\mathrm{fh}}M(W)$, Young's inequality with parameter $\omega_{H,\mathrm C}^{\mathrm{fh}}\bar\alpha_m/(4L_{H,\mathrm C}^{\mathrm{fh}})$, and $\lVert Z\rVert_{H,2,p_H^\star}^2\le q_H\lVert Z\rVert_{H,2,\infty}^2$ yield
\begin{equation}
\begin{gathered}
\mathbb E\left[M_{H,\mathrm C}(W_{m+1,0})\mid U_{m,0}\right]\le\left(1+\frac{\omega_{H,\mathrm C}^{\mathrm{fh}}\bar\alpha_m}{4}\right)M_{H,\mathrm C}(\widetilde W_{m,H})\\
+32L_{H,\mathrm C}^{\mathrm{fh}}q_HB_{H,\mathrm C}^2\left(1+\frac{1}{\omega_{H,\mathrm C}^{\mathrm{fh}}}\right)\bar\alpha_m^2.
\end{gathered}
\end{equation}
Because $(1+x/4)(1-x)\le1-x/2$ for $x\in[0,1]$ and $\bar\alpha_m^2\le H\bar\alpha_m^{(2)}$, the stated drift follows. Finally, substituting $L_{H,\mathrm C}^{\mathrm{fh}}=(p_H^\star-1)q_H/\kappa_H$ and $\omega_{H,\mathrm C}^{\mathrm{fh}}\ge\kappa_H/2$ gives the explicit polynomial upper bound.
\end{proof}

\begin{proposition}[Fixed-horizon CTD finite-iteration bound]\label{prop:fh-ctd-abstract}
Define
\begin{equation}
\resizebox{\linewidth}{!}{$
a^{\mathrm{fh}}_{H, \mathrm C, 1}:=e^2r^{\mathrm{fh}}_{H, \mathrm C},\quad a^{\mathrm{fh}}_{H, \mathrm C, 2}:=c_{H,\mathrm C,1}^{\mathrm{fh}},\quad a^{\mathrm{fh}}_{H, \mathrm C,3}:= \frac{c_{H,\mathrm C,1}^{\mathrm{fh}}}{\bar\alpha_{H, \mathrm C}}, \quad a^{\mathrm{fh}}_{H, \mathrm C, 4}:= 2e^2(1+\vartheta_{H, \mathrm C}q_H)c^{\mathrm{fh}}_{H, \mathrm C,2}.
$}
\end{equation}
These constants obey
\begin{equation}
\begin{gathered}
a_{H,\mathrm C,1}^{\mathrm{fh}}\le2e^2,\qquad a_{H,\mathrm C,2}^{\mathrm{fh}}\ge\frac{\kappa_H}{4},\qquad \bar\alpha_{H,\mathrm C}\ge\min\left\{\frac{1}{H},\frac{\kappa_H^2}{32H(p_H^\star-1)q_H^2}\right\},\\
\frac{a_{H,\mathrm C,4}^{\mathrm{fh}}}{a_{H,\mathrm C,2}^{\mathrm{fh}}}\le\frac{1536e^2H(p_H^\star-1)q_H^2B_{H,\mathrm C}^2}{\kappa_H^3}.
\end{gathered}
\end{equation}
Since $q_H\le e^2$ and $\kappa_H=\rho_{H,\min}/(2(H+1))$, every displayed dependence is polynomial in $H$, $1/\rho_{H,\min}$, $1+\log(H\lvert\X\rvert)$, and $B_{H,\mathrm C}$.
If $(\alpha_k)_{k\ge0}$ is nonincreasing, the within-episode regularity condition \eqref{eq:fh-step-regularity} holds, and
\begin{equation}
\alpha_0 \le \frac{a^{\mathrm{fh}}_{H, \mathrm C,2}}{a^{\mathrm{fh}}_{H, \mathrm C, 3}},
\end{equation}
then, for all episodes $m \ge 0$,
\begin{equation}
\begin{gathered}
\E\bigl[\ell_{H, \mathrm C, \infty}^{\mathrm{unw}}(\eta_{mH}, \eta^\star_{H, \mathrm C})^2 \bigr] \le a^{\mathrm{fh}}_{H, \mathrm C,1} \ell_{H, \mathrm C,\infty}^{\mathrm{unw}}(\eta_0, \eta^\star_{H, \mathrm C})^2 \prod_{j=0}^{m-1}(1-a^{\mathrm{fh}}_{H, \mathrm C,2}\bar\alpha_j) \\+ a^{\mathrm{fh}}_{H, \mathrm C,4}\sum_{i=0}^{m-1}\bar \alpha^{(2)}_i \prod_{j=i+1}^{m-1}(1-a^{\mathrm{fh}}_{H, \mathrm C,2}\bar\alpha_j).
\end{gathered}
\end{equation}
\end{proposition}
\begin{proof}
Iterating Proposition~\ref{prop:fh-ctd-drift} gives
\begin{equation}
\E\bigl[M_{H, \mathrm C}(W_{m, 0}) \bigr] \le M_{H, \mathrm C}(W_{0, 0}) \prod_{j=0}^{m-1}(1- a^{\mathrm{fh}}_{H, \mathrm C,2}\bar\alpha_j) + c^{\mathrm{fh}}_{H, \mathrm C,2}\sum_{i=0}^{m-1}\bar\alpha^{(2)}_i \prod_{j=i+1}^{m-1}(1- a^{\mathrm{fh}}_{H, \mathrm C,2} \bar\alpha_j).
\end{equation}
By the Moreau envelope comparison,
\begin{equation}
\E[\lVert W_{m, 0}\rVert^2_{H, 2, \infty}] \le 2(1+\vartheta_{H, \mathrm C}q_H)\E\bigl[ M_{H, \mathrm C}(W_{m, 0})\bigr]
\end{equation}
and
\begin{equation}
M_{H, \mathrm C}(W_{0, 0}) \le \frac{1}{2(1+\vartheta_{H, \mathrm C})}\lVert W_{0, 0}\rVert^2_{H, 2, \infty}.
\end{equation}
The embedding isometric identity and the choice $\lambda=H/(H+1)$ give
\begin{equation}
\ell_{H,\mathrm C,\infty}^{\mathrm{unw}}(\eta_{mH},\eta_{H,\mathrm C}^\star)^2\le e^2\lVert W_{m,0}\rVert_{H,2,\infty}^2,
\end{equation}
while $\lVert W_{0,0}\rVert_{H,2,\infty}\le\ell_{H,\mathrm C,\infty}^{\mathrm{unw}}(\eta_0,\eta_{H,\mathrm C}^\star)$. These inequalities provide the claim.
\end{proof}

\begin{corollary}[Boundary-iterate step size consequences]
\label{cor:fh-ctd-steps}
Under the hypotheses of Proposition~\ref{prop:fh-ctd-abstract}:

\textup{(a)} if $\alpha_k \equiv \alpha$ and
\begin{equation}
\alpha \le \frac{a_{H,\mathrm C,2}^{\mathrm{fh}}}{a_{H,\mathrm C,3}^{\mathrm{fh}}},
\end{equation}
then for all episodes $m \ge 0$,
\begin{equation}
\mathbb E\left[ \ell_{H,\mathrm C,\infty}^{\mathrm{unw}}(\eta_{mH},\eta_{H,\mathrm C}^\star)^2 \right] \le a_{H,\mathrm C,1}^{\mathrm{fh}} \ell_{H,\mathrm C,\infty}^{\mathrm{unw}}(\eta_0,\eta_{H,\mathrm C}^\star)^2 \bigl(1-a_{H,\mathrm C,2}^{\mathrm{fh}}H\alpha\bigr)^m + \frac{a_{H,\mathrm C,4}^{\mathrm{fh}}}{a_{H,\mathrm C,2}^{\mathrm{fh}}}\alpha.
\end{equation}

For the two diminishing-step cases below, where $g$ is the step-size offset, write $\tau_m:=mH+g+H-1$, so $\tau_0=g+H-1$.

\textup{(b)} if $\alpha_k = \alpha/(k+g)$, $\alpha > 1/a_{H,\mathrm C,2}^{\mathrm{fh}}$, and
\begin{equation}
g \ge \max\left\{H-1,1,\alpha,\frac{\alpha a_{H,\mathrm C,3}^{\mathrm{fh}}}{a_{H,\mathrm C,2}^{\mathrm{fh}}}\right\},
\end{equation}
then the boundary iterates satisfy
\begin{equation}
\mathbb E\left[ \ell_{H,\mathrm C,\infty}^{\mathrm{unw}}(\eta_{mH},\eta_{H,\mathrm C}^\star)^2 \right] \le a_{H,\mathrm C,1}^{\mathrm{fh}}\ell_{H,\mathrm C,\infty}^{\mathrm{unw}}(\eta_0,\eta_{H,\mathrm C}^\star)^2 \left(\frac{\tau_0}{\tau_m}\right)^{a_{H,\mathrm C,2}^{\mathrm{fh}}\alpha}
+\frac{a_{H,\mathrm C,4}^{\mathrm{fh}}H^2\alpha^2}{a_{H,\mathrm C,2}^{\mathrm{fh}}\alpha-1}\cdot\frac{1}{\tau_m}.
\end{equation}

\textup{(c)} if $\alpha_k = \alpha/(k+g)^z$ with $z \in (0,1)$ and
\begin{equation}
g \ge \max\left\{ H-1,1,\alpha^{1/z}, \left(\frac{\alpha a_{H,\mathrm C,3}^{\mathrm{fh}}}{a_{H,\mathrm C,2}^{\mathrm{fh}}}\right)^{1/z}, \left(\frac{2z}{a_{H,\mathrm C,2}^{\mathrm{fh}}\alpha}\right)^{1/(1-z)} \right\},
\end{equation}
then the boundary iterates satisfy
\begin{equation}
\begin{gathered}
\mathbb E\left[ \ell_{H,\mathrm C,\infty}^{\mathrm{unw}}(\eta_{mH},\eta_{H,\mathrm C}^\star)^2 \right] \le a_{H,\mathrm C,1}^{\mathrm{fh}}\ell_{H,\mathrm C,\infty}^{\mathrm{unw}}(\eta_0,\eta_{H,\mathrm C}^\star)^2 \cdot\exp\left(-\frac{a_{H,\mathrm C,2}^{\mathrm{fh}}\alpha}{1-z}\bigl(\tau_m^{1-z}-\tau_0^{1-z}\bigr)\right)\\ +\frac{2a_{H,\mathrm C,4}^{\mathrm{fh}}H^2\alpha}{a_{H,\mathrm C,2}^{\mathrm{fh}}}\cdot\frac{1}{\tau_m^z}.
\end{gathered}
\end{equation}
\end{corollary}

\begin{proof}
For a constant step size, condition \eqref{eq:fh-step-regularity} is immediate. For either diminishing schedule, the requirement $g\ge H-1$ gives
\begin{equation}
\frac{\alpha_{mH+H-1}}{\alpha_{mH}}=\left(\frac{mH+g}{mH+g+H-1}\right)^z\ge\left(\frac{g}{g+H-1}\right)^z\ge\frac{1}{2},
\end{equation}
where $z=1$ in part \textup{(b)}. Thus the regularity hypothesis of Proposition~\ref{prop:fh-ctd-abstract} holds in all three cases. For part \textup{(a)}, substitution of $\bar\alpha_m = H\alpha$ and $\bar\alpha_m^{(2)} = H\alpha^2$ into that proposition yields the claim.

For part \textup{(b)}, set $q=\tau_0/H$ and $\lambda=a_{H,\mathrm C,2}^{\mathrm{fh}}\alpha$. The bounds
\begin{equation}
\frac{H\alpha}{\tau_m} \le \bar\alpha_m \le \frac{H\alpha}{mH+g}, \qquad \bar\alpha_m^{(2)} \le \frac{H\alpha^2}{(mH+g)^2}.
\end{equation}
imply
\begin{equation}
\prod_{j=0}^{m-1}(1-a_{H,\mathrm C,2}^{\mathrm{fh}}\bar\alpha_j)\le \left(\frac{q}{m+q}\right)^\lambda.
\end{equation}
The same elementary product-sum estimate gives
\begin{equation}
\sum_{i=0}^{m-1}\bar\alpha_i^{(2)}\prod_{j=i+1}^{m-1}(1-a_{H,\mathrm C,2}^{\mathrm{fh}}\bar\alpha_j)\le \frac{H\alpha^2}{\lambda-1}\cdot\frac{1}{m+q},
\end{equation}
and the displayed bound follows from Proposition~\ref{prop:fh-ctd-abstract} after substituting $q=\tau_0/H$ and $\tau_m=H(m+q)$.

For part \textup{(c)}, keep the same $q$ and set $A=a_{H,\mathrm C,2}^{\mathrm{fh}}\alpha H^{1-z}$. The bounds
\begin{equation}
\frac{H\alpha}{\tau_m^z} \le \bar\alpha_m \le \frac{H\alpha}{(mH+g)^z}, \qquad \bar\alpha_m^{(2)} \le \frac{H\alpha^2}{(mH+g)^{2z}},
\end{equation}
imply
\begin{equation}
\prod_{j=0}^{m-1}(1-a_{H,\mathrm C,2}^{\mathrm{fh}}\bar\alpha_j)\le \exp\left(-\frac{A}{1-z}\bigl((m+q)^{1-z}-q^{1-z}\bigr)\right).
\end{equation}
The lower bound on $g$ implies $q\ge (2z/A)^{1/(1-z)}$. The elementary polynomial product-sum estimate gives
\begin{equation}
\sum_{i=0}^{m-1}\bar\alpha_i^{(2)}\prod_{j=i+1}^{m-1}(1-a_{H,\mathrm C,2}^{\mathrm{fh}}\bar\alpha_j)\le \frac{2H\alpha^2}{A}\cdot\frac{1}{(m+q)^z}.
\end{equation}
Substituting the definitions of $A$ and $q$ into Proposition~\ref{prop:fh-ctd-abstract}, using $\tau_m=H(m+q)$, and using $H^{2z}\le H^2$, gives the displayed bound.
\end{proof}

\begin{proof}[Proof of Theorem~\ref{thm:undisc-ctd}]
Combine Proposition~\ref{prop:fh-ctd-abstract} with Corollary~\ref{cor:fh-ctd-steps}. Explicitly, take
\begin{equation}
\begin{gathered}
c_{\mathrm C}^H:=a_{H,\mathrm C,2}^{\mathrm{fh}},\qquad r_{H,\mathrm C}:=\frac{a_{H,\mathrm C,3}^{\mathrm{fh}}}{a_{H,\mathrm C,2}^{\mathrm{fh}}},\qquad \bar\alpha_{\mathrm C}^H:=\frac{1}{r_{H,\mathrm C}}, \\
C_{\mathrm C}^H:=\max\left\{a_{H,\mathrm C,1}^{\mathrm{fh}},\frac{a_{H,\mathrm C,4}^{\mathrm{fh}}}{a_{H,\mathrm C,2}^{\mathrm{fh}}},a_{H,\mathrm C,4}^{\mathrm{fh}}H^2,\frac{2a_{H,\mathrm C,4}^{\mathrm{fh}}H^2}{a_{H,\mathrm C,2}^{\mathrm{fh}}}\right\}, \\
g_{\mathrm C}^H(\alpha):=\left\lceil\max\left\{1,\alpha,\alpha r_{H,\mathrm C}\right\}\right\rceil, \\
g_{\mathrm C}^H(\alpha,z):=\left\lceil\max\left\{1,\alpha^{1/z},(\alpha r_{H,\mathrm C})^{1/z},\left(\frac{2z}{a_{H,\mathrm C,2}^{\mathrm{fh}}\alpha}\right)^{1/(1-z)}\right\}\right\rceil.
\end{gathered}
\end{equation}
Each main-text display follows directly from the corresponding part of Corollary~\ref{cor:fh-ctd-steps}.
\end{proof}

\begin{proof}[Proof of Corollary~\ref{cor:undisc-ctd}]
By Corollary~\ref{cor:fh-ctd-steps}\textup{(b)},
\begin{equation}
\mathbb E\left[ \ell_{H,\mathrm C,\infty}^{\mathrm{unw}}(\eta_{mH},\eta_{H,\mathrm C}^\star)^2 \right] = O\left(\frac{1}{m+1}\right).
\end{equation}
Jensen's inequality then gives
\begin{equation}
\mathbb E\left[ \ell_{H,\mathrm C,\infty}^{\mathrm{unw}}(\eta_{mH},\eta_{H,\mathrm C}^\star) \right] = O\left(\frac{1}{\sqrt{m+1}}\right),
\end{equation}
so $m = O(\varepsilon^{-2})$ episodes suffice, with the explicit parameter dependence inherited from Corollary~\ref{cor:fh-ctd-steps}. The transition count is $k=mH$.
\end{proof}

\section{Fixed-horizon MTD}
\label{app:fh-mtd}

This appendix records the fixed-horizon MTD ingredients and the proof of Theorem~\ref{thm:undisc-mtd}.

\subsection{Flattened weighted state--horizon space}

We use the same weighted flattening as in Appendix~\ref{app:fh-ctd}. Let
\begin{equation}
\mathcal S_H := \{(h,s): h\in\{1,\dots,H\},\ s\in \X\}, \qquad \lvert \mathcal S_H \rvert = H\lvert \X \rvert, \qquad p_H^\star := \max\{2,\lceil \log \lvert \mathcal S_H \rvert\rceil\}.
\end{equation}
Let
\begin{equation}
V_{H, \mathrm M} := \prod_{(h,s) \in \mathcal S_H} \R^d
\end{equation}
denote the corresponding flattened Euclidean product space. For a horizon-stacked iterate $U=(U^h(s))$, define
\begin{equation}
\bar{U}(h,s):=\lambda^h U^h(s).
\end{equation}
Then
\begin{equation}
\lVert U \rVert_{H,2,\infty} = \max_{(h,s)\in\mathcal S_H}\lVert \bar{U}(h,s) \rVert_2, \qquad \lVert U \rVert_{H,2,p} = \left( \sum_{(h,s)\in\mathcal S_H}\lVert \bar{U}(h,s) \rVert_2^p \right)^{1/p}.
\end{equation}
For a single horizon stack $U=(U^1,\dots,U^H)$, also set
\begin{equation}
\lVert U \rVert_{H,2} := \max_{1\le h\le H}\lambda^h\lVert U^h \rVert_2.
\end{equation}
Then
\begin{equation}
\lVert U \rVert_{H,2,\infty} = \max_{s\in \X}\lVert U(s) \rVert_{H,2}.
\end{equation}
Thus Appendix~\ref{app:common} applies directly to the flattened weighted process.

\subsection{Fixed-horizon MTD residual map}

For each state $s\in \X$, let
\begin{equation}
U(s):=\bigl(U^1(s),\dots,U^H(s)\bigr),
\end{equation}
and let $P_s^H$ denote the coordinate projector onto that stack. Define
\begin{equation}
F_{H,\mathrm M}(U;s,(r,s')) := P_s^H\bigl(\widehat T_{H,\mathrm M}(U;s,(r,s'))-U(s)\bigr),
\end{equation}
where $\widehat T_{H,\mathrm M}(U;s,(r,s'))$ is the full horizon-stacked MTD target from Section~\ref{sec:undisc}. Then the online fixed-horizon MTD recursion is
\begin{equation}
U_{k+1}=U_k+\alpha_k F_{H,\mathrm M}(U_k;S_k,(R_k,S_{k+1})).
\end{equation}
Define also the MTD embedded averaged operator by
\begin{equation}
\mathcal O_{H,\mathrm M} := I_{H,\mathrm M}\circ \Pi_{H,\mathrm M}^{\Theta} T_H^\pi \circ I_{H,\mathrm M}^{-1}.
\end{equation}

Define the invariant MTD domain
\begin{equation}
\mathcal D_{H,\mathrm M}:=I_{H,\mathrm M}(\mathcal F_{H,\mathrm M,\Theta}^{\X})\subseteq V_{H,\mathrm M}.
\end{equation}
The mass-one coefficient constraint is affine and closed at every horizon-state block, so $\mathcal D_{H,\mathrm M}$ is closed and convex. Both $\mathcal O_{H,\mathrm M}$ and every full-sample replacement $U+F_{H,\mathrm M}(U;s,(r,s'))$ preserve this domain. Lemma~\ref{lem:invariant-domain} therefore shows that the online recursion remains in $\mathcal D_{H,\mathrm M}$ whenever $U_0\in\mathcal D_{H,\mathrm M}$ and $0\le\alpha_k\le1$. All operator estimates in this appendix are required only on this invariant domain.

\begin{proposition}[Fixed-horizon MTD weighted contraction]
\label{prop:fh-mtd-contract-app}
For all $U,U' \in I_{H,\mathrm M}(\mathcal F_{H,\mathrm M,\Theta}^\X)$,
\begin{equation}
\lVert \mathcal O_{H,\mathrm M}U-\mathcal O_{H,\mathrm M}U'\rVert_{H,2,\infty} \le \lambda \lVert U-U'\rVert_{H,2,\infty}.
\end{equation}
\end{proposition}

\begin{proof}
Let $\eta:=I_{H,\mathrm M}^{-1}(U)$ and $\eta':=I_{H,\mathrm M}^{-1}(U')$. Fix $(h,s)\in\mathcal S_H$ with $h \ge 1$. By the statewise MMD isometry and nonexpansiveness of the signed-categorical projection,
\begin{equation}
\begin{aligned}
\lambda^h \lVert (\mathcal O_{H,\mathrm M}U)^h(s)-(\mathcal O_{H,\mathrm M}U')^h(s) \rVert_2 &= \lambda^h \ell_{\mathrm M}\left( (\Pi_{H,\mathrm M}^{\Theta}T_H^\pi\eta)^h(s), (\Pi_{H,\mathrm M}^{\Theta}T_H^\pi\eta')^h(s) \right) \\ 
&\le \lambda^h \ell_{\mathrm M}\left( (T_H^\pi\eta)^h(s), (T_H^\pi\eta')^h(s) \right).
\end{aligned}
\end{equation}
Now
\begin{equation}
(T_H^\pi\eta)^h(s) = \sum_{a\in\A}\pi(a \mid s)\sum_{s'\in\X}P(s' \mid s,a)\, (f_{R(s,a),1})_\#\eta^{h-1}(s'),
\end{equation}
and likewise for $\eta'$. Since $(f_{r,1})_\#$ is an isometry in the MMD metric associated with a shift-invariant kernel and MMD is convex under mixtures,
\begin{equation}
\begin{aligned}
\ell_{\mathrm M}\left( (T_H^\pi\eta)^h(s), (T_H^\pi\eta')^h(s) \right) &\le \sum_{a\in\A}\pi(a \mid s)\sum_{s'\in\X}P(s' \mid s,a)\, \ell_{\mathrm M}\left(\eta^{h-1}(s'),{\eta'}^{h-1}(s')\right) \\ 
&\le \max_{x\in\X} \ell_{\mathrm M}\left(\eta^{h-1}(x),{\eta'}^{h-1}(x)\right).
\end{aligned}
\end{equation}
Therefore
\begin{equation}
\begin{aligned}
\lambda^h \lVert (\mathcal O_{H,\mathrm M}U)^h(s)-(\mathcal O_{H,\mathrm M}U')^h(s) \rVert_2 &\le \lambda \max_{x\in\X} \lambda^{h-1} \ell_{\mathrm M}\left(\eta^{h-1}(x),{\eta'}^{h-1}(x)\right) \\ 
&\le \lambda \ell_{H,\mathrm M,\infty}(\eta,\eta') = \lambda \lVert U-U'\rVert_{H,2,\infty}.
\end{aligned}
\end{equation}
Taking the maximum over $(h,s)\in\mathcal S_H$ proves the claim.
\end{proof}

\subsection{One-step bounds}

\begin{proposition}[Fixed-horizon MTD samplewise Lipschitzness]
\label{prop:fh-mtd-lipschitz}
For every $s \in \X$, $(r, s') \in [0,1]^q \times \X$ and $U, U' \in I_{H,\mathrm M}(\mathcal F_{H,\mathrm M,\Theta}^\X)$,
\begin{equation}
\lVert F_{H,\mathrm M}(U;s,(r,s'))-F_{H,\mathrm M}(U';s,(r,s')) \rVert_{H,2,\infty} \le 2\lVert U-U'\rVert_{H,2,\infty}.
\end{equation}
\end{proposition}

\begin{proof}
For $h=1$, the local target $\widehat T_{H,\mathrm M}^h(U;s,(r,s'))$ is independent of $U$ because the layer-zero law is fixed. For $h\ge2$, it depends only on the $(h-1)$st horizon block at $s'$. The MMD isometry, the Euclidean projection representation in Proposition~\ref{prop:mtd-proj-affine}, and translation invariance of the shift-invariant-kernel MMD give
\begin{equation}
\left\lVert \widehat T_{H,\mathrm M}^h(U;s,(r,s')) - \widehat T_{H,\mathrm M}^h(U';s,(r,s')) \right\rVert_2 \le \lVert U^{h-1}(s')-{U'}^{h-1}(s') \rVert_2.
\end{equation}
Multiplying by $\lambda^h$ yields
\begin{equation}
\resizebox{\linewidth}{!}{$
\lambda^h \left\lVert \widehat T_{H,\mathrm M}^h(U;s,(r,s')) - \widehat T_{H,\mathrm M}^h(U';s,(r,s')) \right\rVert_2 \le \lambda^{h-1}\lVert U^{h-1}(s')-{U'}^{h-1}(s') \rVert_2 \le \lVert U-U'\rVert_{H,2,\infty}.
$}
\end{equation}
Taking the maximum over $h$ gives
\begin{equation}
\lVert \widehat T_{H,\mathrm M}(U;s,(r,s'))-\widehat T_{H,\mathrm M}(U';s,(r,s')) \rVert_{H,2} \le \lVert U-U'\rVert_{H,2,\infty}.
\end{equation}
Since
\begin{equation}
F_{H,\mathrm M}(U;s,(r,s')) = P_s^H\bigl(\widehat T_{H,\mathrm M}(U;s,(r,s'))-U(s)\bigr),
\end{equation}
the residual difference splits into a projected target difference and a projected current-stack difference. Bounding both by $\lVert U-U'\rVert_{H,2,\infty}$ gives the stated factor $2$.
\end{proof}

\begin{proposition}[Fixed-horizon MTD affine perturbation bound]
\label{prop:fh-mtd-perturb}
Let
\begin{equation}
U^\star_{H,\mathrm M} := I_{H,\mathrm M}(\eta_{H,\mathrm M}^\star).
\end{equation}
Here $\eta_{H,\mathrm M}^\star$ denotes the unique fixed point of $\Pi_{H,\mathrm M}^{\Theta}T_H^\pi$. Define
\begin{equation}
B^{\star,\mathrm{tar}}_{H,\mathrm M}:=\max_{s,s'\in\X}\sup_{r\in[0,1]^q}\left\lVert\widehat T_{H,\mathrm M}(U^\star_{H,\mathrm M};s,(r,s'))-(\mathcal O_{H,\mathrm M}U^\star_{H,\mathrm M})(s)\right\rVert_{H,2},
\end{equation}
\begin{equation}
B^{\mathrm{tar}}_{H,\mathrm M}:=2\lambda\lVert U^\star_{H,\mathrm M}\rVert_{H,2,\infty}+B^{\star,\mathrm{tar}}_{H,\mathrm M},
\end{equation}
\begin{equation}
B^{\star,\mathrm{res}}_{H,\mathrm M}:=\max_{s,s'\in\X}\sup_{r\in[0,1]^q}\lVert F_{H,\mathrm M}(U^\star_{H,\mathrm M};s,(r,s'))\rVert_{H,2,\infty},
\end{equation}
and
\begin{equation}
B^{\mathrm{res}}_{H,\mathrm M}:=2\lVert U^\star_{H,\mathrm M}\rVert_{H,2,\infty}+B^{\star,\mathrm{res}}_{H,\mathrm M}.
\end{equation}
These constants are finite. For every admissible sample $(s,(r,s'))$ and every admissible iterate $U$,
\begin{equation}
\left\lVert \widehat T_{H,\mathrm M}(U; s, (r,s')) - (\mathcal O_{H,\mathrm M}U)(s) \right\rVert_{H,2} \le 2 \lambda \lVert U \rVert_{H,2,\infty} + B^{\mathrm{tar}}_{H,\mathrm M},
\end{equation}
and
\begin{equation}
\lVert F_{H,\mathrm M}(U; s, (r,s')) \rVert_{H,2,\infty} \le 2 \lVert U \rVert_{H,2,\infty} + B^{\mathrm{res}}_{H,\mathrm M}.
\end{equation}
Consequently,
\begin{equation}
\mathbb E \Bigl[ \left\lVert \widehat T_{H,\mathrm M}(U_k; S_k, (R_k,S_{k+1})) - (\mathcal O_{H,\mathrm M}U_k)(S_k) \right\rVert_{H,2}^2 \mid U_k, S_k \Bigr] \le 2 \bigl(B^{\mathrm{tar}}_{H,\mathrm M}\bigr)^2 + 8 \lambda^2 \lVert U_k \rVert_{H,2,\infty}^2.
\end{equation}
\end{proposition}

\begin{proof}
For fixed $s,s' \in \X$ and $r \in [0,1]^q$, the target-level estimate established in the proof of Proposition~\ref{prop:fh-mtd-lipschitz} gives, for each horizon $h$,
\begin{equation}
\left\lVert \widehat T_{H,\mathrm M}^h(U; s, (r,s')) - \widehat T_{H,\mathrm M}^h(U'; s, (r,s')) \right\rVert_2 \le \lVert U^{h-1}(s')-{U'}^{h-1}(s') \rVert_2.
\end{equation}
After multiplying by $\lambda^h$ and taking the maximum over $h$, the horizon shift contributes exactly one factor $\lambda$. Therefore
\begin{equation}
\left\lVert \widehat T_{H,\mathrm M}(U; s, (r,s')) - \widehat T_{H,\mathrm M}(U'; s, (r,s')) \right\rVert_{H,2} \le \lambda \lVert U-U' \rVert_{H,2,\infty}.
\end{equation}
Likewise, Proposition~\ref{prop:fh-mtd-contract-app} yields
\begin{equation}
\lVert (\mathcal O_{H,\mathrm M}U)(s) - (\mathcal O_{H,\mathrm M}U')(s) \rVert_{H,2} \le \lambda \lVert U-U' \rVert_{H,2,\infty}.
\end{equation}

The constant $B^{\star,\mathrm{tar}}_{H,\mathrm M}$ is finite because $[0,1]^q$ is compact, there are only finitely many state pairs, and the embedded projected target depends continuously on $r$.

Then
\begin{equation}
\begin{gathered}
\left\lVert \widehat T_{H,\mathrm M}(U; s, (r,s')) - (\mathcal O_{H,\mathrm M}U)(s) \right\rVert_{H,2} \le \left\lVert \widehat T_{H,\mathrm M}(U; s, (r,s')) - \widehat T_{H,\mathrm M}(U^\star_{H,\mathrm M}; s, (r,s')) \right\rVert_{H,2} \\ 
+ \left\lVert \widehat T_{H,\mathrm M}(U^\star_{H,\mathrm M}; s, (r,s')) - (\mathcal O_{H,\mathrm M}U^\star_{H,\mathrm M})(s) \right\rVert_{H,2} + \left\lVert (\mathcal O_{H,\mathrm M}U^\star_{H,\mathrm M})(s) - (\mathcal O_{H,\mathrm M}U)(s) \right\rVert_{H,2} \\ 
\le 2 \lambda \lVert U-U^\star_{H,\mathrm M} \rVert_{H,2,\infty} + B^{\star,\mathrm{tar}}_{H,\mathrm M} \\
\le 2 \lambda \lVert U \rVert_{H,2,\infty} + \Bigl( 2 \lambda \lVert U^\star_{H,\mathrm M} \rVert_{H,2,\infty} + B^{\star,\mathrm{tar}}_{H,\mathrm M} \Bigr).
\end{gathered}
\end{equation}
Hence the target-level affine bound holds with the stated $B^{\mathrm{tar}}_{H,\mathrm M}$.

For the residual map, Proposition~\ref{prop:fh-mtd-lipschitz} gives
\begin{equation}
\left\lVert F_{H,\mathrm M}(U; s, (r,s')) - F_{H,\mathrm M}(U'; s, (r,s')) \right\rVert_{H,2,\infty} \le 2 \lVert U-U' \rVert_{H,2,\infty}.
\end{equation}
The constant $B^{\star,\mathrm{res}}_{H,\mathrm M}$ is finite by the same continuity and compactness argument. Then
\begin{equation}
\resizebox{\linewidth}{!}{$
\lVert F_{H,\mathrm M}(U; s, (r,s')) \rVert_{H,2,\infty} \le 2 \lVert U-U^\star_{H,\mathrm M} \rVert_{H,2,\infty} + B^{\star,\mathrm{res}}_{H,\mathrm M} \le 2 \lVert U \rVert_{H,2,\infty} + \Bigl( 2 \lVert U^\star_{H,\mathrm M} \rVert_{H,2,\infty} + B^{\star,\mathrm{res}}_{H,\mathrm M} \Bigr).
$}
\end{equation}
Thus the residual-level affine bound holds with the stated $B^{\mathrm{res}}_{H,\mathrm M}$.

Finally, the conditional second-moment estimate follows from the target-level affine bound by squaring and using
\begin{equation}
(a+b)^2 \le 2a^2 + 2b^2
\end{equation}
with
\begin{equation}
a := 2 \lambda \lVert U_k \rVert_{H,2,\infty}, \qquad b := B^{\mathrm{tar}}_{H,\mathrm M}.
\end{equation}
\end{proof}


\subsection{Phasewise averaged maps}

For each phase $t\in\{0,\dots,H-1\}$, define
\begin{equation}
\Gamma_{t,\mathrm M}(U) := \sum_{s\in \X}\rho_t(s)\, P_s^H\bigl((\mathcal O_{H,\mathrm M}U)(s)-U(s)\bigr), \qquad G_{t,\mathrm M}(U):=U+\Gamma_{t,\mathrm M}(U).
\end{equation}

\begin{proposition}[Episode-averaged MTD contraction]
\label{prop:fh-mtd-avg-contract}
Suppose $\bar\alpha_m>0$ and define
\begin{equation}
\bar\alpha_m:=\sum_{t=0}^{H-1}\alpha_{mH+t},\qquad \rho_m^\alpha(s):=\frac{1}{\bar\alpha_m}\sum_{t=0}^{H-1}\alpha_{mH+t}\rho_t(s).
\end{equation}
If the step sizes are nonincreasing within episode $m$ and satisfy \eqref{eq:fh-step-regularity}, then
\begin{equation}
\min_{s\in\X}\rho_m^\alpha(s)\ge\frac{\rho_{H,\min}}{2}.
\end{equation}
Define
\begin{equation}
G_{m,\mathrm M}^\alpha(U):=U+\sum_{s\in\X}\rho_m^\alpha(s)P_s^H\bigl((\mathcal O_{H,\mathrm M}U)(s)-U(s)\bigr).
\end{equation}
Then
\begin{equation}
\lVert G_{m,\mathrm M}^\alpha(U)-G_{m,\mathrm M}^\alpha(U')\rVert_{H,2,\infty}\le\bar\beta_H^{\mathrm{occ}}\lVert U-U'\rVert_{H,2,\infty},
\end{equation}
where
\begin{equation}
\bar\beta_H^{\mathrm{occ}}:=1-\frac{\rho_{H,\min}}{2}(1-\lambda),\qquad \kappa_H^{\mathrm{occ}}:=1-\bar\beta_H^{\mathrm{occ}}=\frac{\rho_{H,\min}}{2(H+1)}.
\end{equation}
Moreover, the frozen episode map
\begin{equation}
A_{m,\mathrm M}^{\mathrm{fr}}(U):=U+\sum_{t=0}^{H-1}\alpha_{mH+t}\Gamma_{t,\mathrm M}(U)=(1-\bar\alpha_m)U+\bar\alpha_mG_{m,\mathrm M}^\alpha(U)
\end{equation}
fixes $U_{H,\mathrm M}^\star$ and is a $(1-\kappa_H^{\mathrm{occ}}\bar\alpha_m)$-contraction whenever $\bar\alpha_m\le1$.
\end{proposition}

\begin{proof}
The same weighted-occupancy calculation as in Proposition~\ref{prop:fh-ctd-avg-contract} gives
\begin{equation}
\rho_m^\alpha(s)\ge\frac{\alpha_{mH+H-1}}{H\alpha_{mH}}\sum_{t=0}^{H-1}\rho_t(s)\ge\frac{\rho_{H,\min}}{2}.
\end{equation}
For each state block $s$, Proposition~\ref{prop:fh-mtd-contract-app} gives
\begin{equation}
\left\lVert\bigl(G_{m,\mathrm M}^\alpha(U)-G_{m,\mathrm M}^\alpha(U')\bigr)(s)\right\rVert_{H,2}\le\bigl(1-\rho_m^\alpha(s)(1-\lambda)\bigr)\lVert U-U'\rVert_{H,2,\infty}\le\bar\beta_H^{\mathrm{occ}}\lVert U-U'\rVert_{H,2,\infty}.
\end{equation}
Taking the maximum over $s$ proves the contraction. The identity for $A_{m,\mathrm M}^{\mathrm{fr}}$ follows from the definition of $\rho_m^\alpha$. Convexity of the norm gives its contraction factor, and $\Gamma_{t,\mathrm M}(U_{H,\mathrm M}^\star)=0$ for every phase gives the fixed-point claim.
\end{proof}

\subsection{Auxiliary phasewise bounds}

For episode $m$, write
\begin{equation}
U_{m,t}:=U_{mH+t},\qquad t=0,\dots,H.
\end{equation}

\begin{lemma}[Within-episode growth and movement]
\label{lem:fh-mtd-within}
Suppose $\bar\alpha_m\le1$. Define
\begin{equation}
\begin{gathered}
C_{H,\mathrm M}^{\mathrm{grow}}:=e^2\bigl(1+B_{H,\mathrm M}^{\mathrm{res}}\bigr),\\
C_{H,\mathrm M}^{\mathrm{move}}:=2C_{H,\mathrm M}^{\mathrm{grow}}+B_{H,\mathrm M}^{\mathrm{res}}.
\end{gathered}
\end{equation}
Then, pathwise,
\begin{equation}
\max_{0\le t\le H}\lVert U_{m,t}\rVert_{H,2,\infty}\le C_{H,\mathrm M}^{\mathrm{grow}}\bigl(1+\lVert U_{m,0}\rVert_{H,2,\infty}\bigr)
\end{equation}
and
\begin{equation}
\max_{0\le t\le H}\lVert U_{m,t}-U_{m,0}\rVert_{H,2,\infty}\le C_{H,\mathrm M}^{\mathrm{move}}\bar\alpha_m\bigl(1+\lVert U_{m,0}\rVert_{H,2,\infty}\bigr).
\end{equation}
\end{lemma}

\begin{proof}
Proposition~\ref{prop:fh-mtd-perturb} gives
\begin{equation}
\lVert U_{m,t+1}\rVert_{H,2,\infty}\le(1+2\alpha_{mH+t})\lVert U_{m,t}\rVert_{H,2,\infty}+\alpha_{mH+t}B_{H,\mathrm M}^{\mathrm{res}}.
\end{equation}
Adding $B_{H,\mathrm M}^{\mathrm{res}}/2$ to both sides and iterating gives
\begin{equation}
\lVert U_{m,t}\rVert_{H,2,\infty}+\frac{B_{H,\mathrm M}^{\mathrm{res}}}{2}\le\exp\left(2\sum_{u=0}^{t-1}\alpha_{mH+u}\right)\left(\lVert U_{m,0}\rVert_{H,2,\infty}+\frac{B_{H,\mathrm M}^{\mathrm{res}}}{2}\right).
\end{equation}
The first claim follows from $\bar\alpha_m\le1$. Summing the residuals in the online recursion and applying that claim gives
\begin{equation}
\begin{gathered}
\lVert U_{m,t}-U_{m,0}\rVert_{H,2,\infty}\le\sum_{u=0}^{t-1}\alpha_{mH+u}\bigl(2\lVert U_{m,u}\rVert_{H,2,\infty}+B_{H,\mathrm M}^{\mathrm{res}}\bigr)\\
\le C_{H,\mathrm M}^{\mathrm{move}}\bar\alpha_m\bigl(1+\lVert U_{m,0}\rVert_{H,2,\infty}\bigr).
\end{gathered}
\end{equation}
\end{proof}

\begin{lemma}[Frozen residual decomposition]
\label{lem:fh-mtd-bias}
For each phase $t\in\{0,\dots,H-1\}$, define
\begin{equation}
\zeta_{m,t}^{\mathrm M}:=F_{H,\mathrm M}(U_{m,0};S_t^m,(R_t^m,S_{t+1}^m))-\Gamma_{t,\mathrm M}(U_{m,0})
\end{equation}
and
\begin{equation}
\chi_{m,t}^{\mathrm M}:=F_{H,\mathrm M}(U_{m,t};S_t^m,(R_t^m,S_{t+1}^m))-F_{H,\mathrm M}(U_{m,0};S_t^m,(R_t^m,S_{t+1}^m)).
\end{equation}
Then
\begin{equation}
\mathbb E\left[\zeta_{m,t}^{\mathrm M}\mid U_{m,0}\right]=0,\qquad
\lVert\zeta_{m,t}^{\mathrm M}\rVert_{H,2,\infty}\le\bigl(4\lVert U_{m,0}\rVert_{H,2,\infty}+2B_{H,\mathrm M}^{\mathrm{res}}\bigr),
\end{equation}
and, whenever $\bar\alpha_m\le1$,
\begin{equation}
\lVert\chi_{m,t}^{\mathrm M}\rVert_{H,2,\infty}\le2C_{H,\mathrm M}^{\mathrm{move}}\bar\alpha_m\bigl(1+\lVert U_{m,0}\rVert_{H,2,\infty}\bigr).
\end{equation}
Define
\begin{equation}
\begin{gathered}
C_{H,\mathrm M}^{\mathrm{noise}}:=4+2B_{H,\mathrm M}^{\mathrm{res}},\qquad
C_{H,\mathrm M}^{\mathrm{end}}:=C_{H,\mathrm M}^{\mathrm{noise}}+2C_{H,\mathrm M}^{\mathrm{move}},\\
D_m^{\mathrm M}:=U_{m,H}-A_{m,\mathrm M}^{\mathrm{fr}}(U_{m,0}).
\end{gathered}
\end{equation}
Then
\begin{equation}
D_m^{\mathrm M}=\sum_{t=0}^{H-1}\alpha_{mH+t}\bigl(\zeta_{m,t}^{\mathrm M}+\chi_{m,t}^{\mathrm M}\bigr).
\end{equation}
Consequently,
\begin{equation}
\left\lVert\mathbb E[D_m^{\mathrm M}\mid U_{m,0}]\right\rVert_{H,2,\infty}\le2C_{H,\mathrm M}^{\mathrm{move}}\bar\alpha_m^2\bigl(1+\lVert U_{m,0}\rVert_{H,2,\infty}\bigr)
\end{equation}
and
\begin{equation}
\lVert D_m^{\mathrm M}\rVert_{H,2,\infty}\le C_{H,\mathrm M}^{\mathrm{end}}\bar\alpha_m\bigl(1+\lVert U_{m,0}\rVert_{H,2,\infty}\bigr)
\end{equation}
pathwise.
\end{lemma}

\begin{proof}
The episode is independent of $U_{m,0}$, and its phase-$t$ sample has the marginal law defining $\Gamma_{t,\mathrm M}$. This proves the conditional mean identity. Proposition~\ref{prop:fh-mtd-perturb} gives
\begin{equation}
\lVert F_{H,\mathrm M}(U_{m,0};S_t^m,(R_t^m,S_{t+1}^m))\rVert_{H,2,\infty}\le2\lVert U_{m,0}\rVert_{H,2,\infty}+B_{H,\mathrm M}^{\mathrm{res}}.
\end{equation}
The same bound holds for $\Gamma_{t,\mathrm M}(U_{m,0})$ by averaging, which proves the bound on $\zeta_{m,t}^{\mathrm M}$. Proposition~\ref{prop:fh-mtd-lipschitz} and Lemma~\ref{lem:fh-mtd-within} give the bound on $\chi_{m,t}^{\mathrm M}$.

Summing the online recursion and subtracting $A_{m,\mathrm M}^{\mathrm{fr}}(U_{m,0})$ gives the identity for $D_m^{\mathrm M}$. The frozen terms have conditional mean zero, so summing the bound on $\chi_{m,t}^{\mathrm M}$ proves the conditional-mean estimate. Finally,
\begin{equation}
\lVert\zeta_{m,t}^{\mathrm M}\rVert_{H,2,\infty}\le C_{H,\mathrm M}^{\mathrm{noise}}\bigl(1+\lVert U_{m,0}\rVert_{H,2,\infty}\bigr),
\end{equation}
and $\bar\alpha_m\le1$ give the pathwise endpoint estimate.
\end{proof}

\subsection{Episode-level finite-iteration drift}
\label{sec:fh-mtd-drift-app}

Let
\begin{equation}
W_{m,0}:=U_{m,0}-U_{H,\mathrm M}^\star,\qquad \bar\alpha_m^{(2)}:=\sum_{t=0}^{H-1}\alpha_{mH+t}^2,
\end{equation}
and set
\begin{equation}
q_H:=\lvert\mathcal S_H\rvert^{2/p_H^\star},\qquad \kappa_H:=\kappa_H^{\mathrm{occ}}=\frac{\rho_{H,\min}}{2(H+1)}.
\end{equation}
Choose
\begin{equation}
\vartheta_{H,\mathrm M}:=\frac{\kappa_H}{q_H}
\end{equation}
and define
\begin{equation}
M_{H,\mathrm M}(W):=\inf_{Z\in V_{H,\mathrm M}}\left\{\frac{1}{2}\lVert Z\rVert_{H,2,\infty}^2+\frac{1}{2\vartheta_{H,\mathrm M}q_H}\lVert W-Z\rVert_{H,2,p_H^\star}^2\right\}.
\end{equation}
The Moreau-envelope comparison gives
\begin{equation}
(1+\vartheta_{H,\mathrm M})M_{H,\mathrm M}(W)\le\frac{1}{2}\lVert W\rVert_{H,2,\infty}^2\le(1+\vartheta_{H,\mathrm M}q_H)M_{H,\mathrm M}(W).
\end{equation}
Define
\begin{equation}
\begin{gathered}
r_{H,\mathrm M}^{\mathrm{fh}}:=\frac{1+\vartheta_{H,\mathrm M}q_H}{1+\vartheta_{H,\mathrm M}},\qquad
L_{H,\mathrm M}^{\mathrm{fh}}:=\frac{p_H^\star-1}{\vartheta_{H,\mathrm M}},\\
\omega_{H,\mathrm M}^{\mathrm{fh}}:=1-(1-\kappa_H)\sqrt{r_{H,\mathrm M}^{\mathrm{fh}}},\qquad
d_{H,\mathrm M}^{\mathrm{fh}}:=8L_{H,\mathrm M}^{\mathrm{fh}}q_H(1+\vartheta_{H,\mathrm M}q_H).
\end{gathered}
\end{equation}
As in Proposition~\ref{prop:explicit-discounted-constants},
\begin{equation}
r_{H,\mathrm M}^{\mathrm{fh}}\le2,\qquad
\omega_{H,\mathrm M}^{\mathrm{fh}}\ge\frac{\kappa_H}{2},\qquad
d_{H,\mathrm M}^{\mathrm{fh}}\le\frac{16(p_H^\star-1)q_H^2}{\kappa_H}.
\end{equation}
The endpoint error coefficient and a convenient polynomial upper bound are
\begin{equation}
\begin{gathered}
C_{H,\mathrm M}^{\mathrm{pot}}:=L_{H,\mathrm M}^{\mathrm{fh}}q_H\left(\frac{8(C_{H,\mathrm M}^{\mathrm{move}})^2}{\omega_{H,\mathrm M}^{\mathrm{fh}}}+\frac{(C_{H,\mathrm M}^{\mathrm{end}})^2}{2}\right),\\
\overline C_{H,\mathrm M}^{\mathrm{pot}}:=(p_H^\star-1)q_H^2\left(\frac{16(C_{H,\mathrm M}^{\mathrm{move}})^2}{\kappa_H^2}+\frac{(C_{H,\mathrm M}^{\mathrm{end}})^2}{2\kappa_H}\right).
\end{gathered}
\end{equation}
Then $C_{H,\mathrm M}^{\mathrm{pot}}\le\overline C_{H,\mathrm M}^{\mathrm{pot}}$. Define
\begin{equation}
\bar\alpha_{H,\mathrm M}:=\min\left\{\frac{1}{H},\frac{\omega_{H,\mathrm M}^{\mathrm{fh}}}{Hd_{H,\mathrm M}^{\mathrm{fh}}},\frac{\omega_{H,\mathrm M}^{\mathrm{fh}}}{32H C_{H,\mathrm M}^{\mathrm{pot}}(1+\vartheta_{H,\mathrm M}q_H)}\right\},
\end{equation}
and
\begin{equation}
c_{H,\mathrm M,1}^{\mathrm{fh}}:=\frac{\omega_{H,\mathrm M}^{\mathrm{fh}}}{4},\qquad
c_{H,\mathrm M,2}^{\mathrm{fh}}:=2H C_{H,\mathrm M}^{\mathrm{pot}}\left(1+2\lVert U_{H,\mathrm M}^\star\rVert_{H,2,\infty}^2\right).
\end{equation}

\begin{proposition}[Episodewise MTD potential drift]
\label{prop:fh-mtd-drift}
Suppose that $(\alpha_k)_{k\ge0}$ is nonincreasing, $\alpha_0\le\bar\alpha_{H,\mathrm M}$, and the within-episode regularity condition \eqref{eq:fh-step-regularity} holds. Then
\begin{equation}
\mathbb E\left[M_{H,\mathrm M}(W_{m+1,0})\mid U_{m,0}\right]\le\left(1-c_{H,\mathrm M,1}^{\mathrm{fh}}\bar\alpha_m\right)M_{H,\mathrm M}(W_{m,0})+c_{H,\mathrm M,2}^{\mathrm{fh}}\bar\alpha_m^{(2)}.
\end{equation}
Moreover,
\begin{equation}
c_{H,\mathrm M,1}^{\mathrm{fh}}\ge\frac{\kappa_H}{8}
\end{equation}
and
\begin{equation}
\bar\alpha_{H,\mathrm M}\ge\min\left\{\frac{1}{H},\frac{\kappa_H^2}{32H(p_H^\star-1)q_H^2},\frac{\kappa_H}{128H\overline C_{H,\mathrm M}^{\mathrm{pot}}}\right\}.
\end{equation}
\end{proposition}

\begin{proof}
Set
\begin{equation}
\widetilde U_{m,H}:=A_{m,\mathrm M}^{\mathrm{fr}}(U_{m,0}),\qquad
\widetilde W_{m,H}:=\widetilde U_{m,H}-U_{H,\mathrm M}^\star.
\end{equation}
Proposition~\ref{prop:fh-mtd-avg-contract} writes the frozen map as a $\bar\alpha_m$-relaxation of a $(1-\kappa_H)$-contraction. The Moreau-envelope smoothness inequality gives
\begin{equation}
M_{H,\mathrm M}(\widetilde W_{m,H})\le\left(1-2\omega_{H,\mathrm M}^{\mathrm{fh}}\bar\alpha_m+d_{H,\mathrm M}^{\mathrm{fh}}\bar\alpha_m^2\right)M_{H,\mathrm M}(W_{m,0}).
\end{equation}
The linear term follows from the dual-norm bound
\begin{equation}
\left\langle\nabla M_{H,\mathrm M}(W),G_{m,\mathrm M}^\alpha(U_{H,\mathrm M}^\star+W)-U_{H,\mathrm M}^\star-W\right\rangle\le-2\omega_{H,\mathrm M}^{\mathrm{fh}}M_{H,\mathrm M}(W),
\end{equation}
and the quadratic smoothness remainder is bounded by $d_{H,\mathrm M}^{\mathrm{fh}}\bar\alpha_m^2M_{H,\mathrm M}(W)$. Since $\bar\alpha_m\le H\alpha_0\le\omega_{H,\mathrm M}^{\mathrm{fh}}/d_{H,\mathrm M}^{\mathrm{fh}}$,
\begin{equation}
M_{H,\mathrm M}(\widetilde W_{m,H})\le\left(1-\omega_{H,\mathrm M}^{\mathrm{fh}}\bar\alpha_m\right)M_{H,\mathrm M}(W_{m,0}).
\end{equation}

Lemma~\ref{lem:fh-mtd-bias} gives $W_{m+1,0}=\widetilde W_{m,H}+D_m^{\mathrm M}$ together with
\begin{equation}
\left\lVert\mathbb E[D_m^{\mathrm M}\mid U_{m,0}]\right\rVert_{H,2,\infty}\le2C_{H,\mathrm M}^{\mathrm{move}}\bar\alpha_m^2Y_m,\qquad
\mathbb E\left[\lVert D_m^{\mathrm M}\rVert_{H,2,\infty}^2\mid U_{m,0}\right]\le(C_{H,\mathrm M}^{\mathrm{end}})^2\bar\alpha_m^2Y_m^2,
\end{equation}
where $Y_m:=1+\lVert U_{m,0}\rVert_{H,2,\infty}$. Smoothness, the self-bounding gradient estimate, Young's inequality with parameter $\omega_{H,\mathrm M}^{\mathrm{fh}}\bar\alpha_m/(4L_{H,\mathrm M}^{\mathrm{fh}})$, and $\lVert Z\rVert_{H,2,p_H^\star}^2\le q_H\lVert Z\rVert_{H,2,\infty}^2$ give
\begin{equation}
\mathbb E\left[M_{H,\mathrm M}(W_{m+1,0})\mid U_{m,0}\right]\le\left(1+\frac{\omega_{H,\mathrm M}^{\mathrm{fh}}\bar\alpha_m}{4}\right)M_{H,\mathrm M}(\widetilde W_{m,H})+C_{H,\mathrm M}^{\mathrm{pot}}\bar\alpha_m^2Y_m^2.
\end{equation}
Because $(1+x/4)(1-x)\le1-x/2$ for $x\in[0,1]$,
\begin{equation}
\mathbb E\left[M_{H,\mathrm M}(W_{m+1,0})\mid U_{m,0}\right]\le\left(1-\frac{\omega_{H,\mathrm M}^{\mathrm{fh}}}{2}\bar\alpha_m\right)M_{H,\mathrm M}(W_{m,0})+C_{H,\mathrm M}^{\mathrm{pot}}\bar\alpha_m^2Y_m^2.
\end{equation}
The envelope comparison implies
\begin{equation}
Y_m^2\le8(1+\vartheta_{H,\mathrm M}q_H)M_{H,\mathrm M}(W_{m,0})+2\left(1+2\lVert U_{H,\mathrm M}^\star\rVert_{H,2,\infty}^2\right).
\end{equation}
Using $\bar\alpha_m^2\le H\bar\alpha_m^{(2)}$ and $\bar\alpha_m^{(2)}\le\alpha_0\bar\alpha_m$, the last condition in the definition of $\bar\alpha_{H,\mathrm M}$ absorbs the state-dependent term into one half of the negative drift. This proves the stated recursion. The explicit lower bounds follow from $\omega_{H,\mathrm M}^{\mathrm{fh}}\ge\kappa_H/2$, $d_{H,\mathrm M}^{\mathrm{fh}}\le16(p_H^\star-1)q_H^2/\kappa_H$, $1+\vartheta_{H,\mathrm M}q_H=1+\kappa_H\le2$, and $C_{H,\mathrm M}^{\mathrm{pot}}\le\overline C_{H,\mathrm M}^{\mathrm{pot}}$.
\end{proof}


\begin{proposition}[Fixed-horizon MTD finite-iteration bound]
\label{prop:fh-mtd-abstract}
Define
\begin{equation}
\begin{gathered}
a_{H,\mathrm M,1}^{\mathrm{fh}}:=e^2r_{H,\mathrm M}^{\mathrm{fh}},\qquad
a_{H,\mathrm M,2}^{\mathrm{fh}}:=c_{H,\mathrm M,1}^{\mathrm{fh}},\qquad
a_{H,\mathrm M,3}^{\mathrm{fh}}:=\frac{c_{H,\mathrm M,1}^{\mathrm{fh}}}{\bar\alpha_{H,\mathrm M}},\\
a_{H,\mathrm M,4}^{\mathrm{fh}}:=2e^2(1+\vartheta_{H,\mathrm M}q_H)c_{H,\mathrm M,2}^{\mathrm{fh}}.
\end{gathered}
\end{equation}
These constants obey
\begin{equation}
\begin{gathered}
a_{H,\mathrm M,1}^{\mathrm{fh}}\le2e^2,\qquad
a_{H,\mathrm M,2}^{\mathrm{fh}}\ge\frac{\kappa_H}{8},\\
\bar\alpha_{H,\mathrm M}\ge\min\left\{\frac{1}{H},\frac{\kappa_H^2}{32H(p_H^\star-1)q_H^2},\frac{\kappa_H}{128H\overline C_{H,\mathrm M}^{\mathrm{pot}}}\right\},\\
\frac{a_{H,\mathrm M,4}^{\mathrm{fh}}}{a_{H,\mathrm M,2}^{\mathrm{fh}}}\le\frac{64e^2H\overline C_{H,\mathrm M}^{\mathrm{pot}}}{\kappa_H}\left(1+2\lVert U_{H,\mathrm M}^\star\rVert_{H,2,\infty}^2\right).
\end{gathered}
\end{equation}
Since $q_H\le e^2$ and $\kappa_H=\rho_{H,\min}/(2(H+1))$, the displayed dependence is polynomial in $H$, $1/\rho_{H,\min}$, $1+\log(H\lvert\X\rvert)$, $B_{H,\mathrm M}^{\mathrm{res}}$, and $\lVert U_{H,\mathrm M}^\star\rVert_{H,2,\infty}$. If $(\alpha_k)_{k\ge0}$ is nonincreasing, \eqref{eq:fh-step-regularity} holds, and
\begin{equation}
\alpha_0\le\frac{a_{H,\mathrm M,2}^{\mathrm{fh}}}{a_{H,\mathrm M,3}^{\mathrm{fh}}},
\end{equation}
then, for all episodes $m\ge0$,
\begin{equation}
\begin{gathered}
\mathbb E\left[\ell_{H,\mathrm M,\infty}^{\mathrm{unw}}\bigl(\eta_{mH},\eta_{H,\mathrm M}^\star\bigr)^2\right]\le
a_{H,\mathrm M,1}^{\mathrm{fh}}\ell_{H,\mathrm M,\infty}^{\mathrm{unw}}\bigl(\eta_0,\eta_{H,\mathrm M}^\star\bigr)^2
\prod_{j=0}^{m-1}\bigl(1-a_{H,\mathrm M,2}^{\mathrm{fh}}\bar\alpha_j\bigr)\\
+a_{H,\mathrm M,4}^{\mathrm{fh}}\sum_{i=0}^{m-1}\bar\alpha_i^{(2)}
\prod_{j=i+1}^{m-1}\bigl(1-a_{H,\mathrm M,2}^{\mathrm{fh}}\bar\alpha_j\bigr).
\end{gathered}
\end{equation}
\end{proposition}

\begin{proof}
Iterating Proposition~\ref{prop:fh-mtd-drift} gives
\begin{equation}
\mathbb E\left[M_{H,\mathrm M}(W_{m,0})\right]\le
M_{H,\mathrm M}(W_{0,0})\prod_{j=0}^{m-1}\bigl(1-a_{H,\mathrm M,2}^{\mathrm{fh}}\bar\alpha_j\bigr)
+c_{H,\mathrm M,2}^{\mathrm{fh}}\sum_{i=0}^{m-1}\bar\alpha_i^{(2)}
\prod_{j=i+1}^{m-1}\bigl(1-a_{H,\mathrm M,2}^{\mathrm{fh}}\bar\alpha_j\bigr).
\end{equation}
The envelope comparison and $\lambda=H/(H+1)$ give
\begin{equation}
\ell_{H,\mathrm M,\infty}^{\mathrm{unw}}(\eta_{mH},\eta_{H,\mathrm M}^\star)^2
\le e^2\lVert W_{m,0}\rVert_{H,2,\infty}^2
\le2e^2(1+\vartheta_{H,\mathrm M}q_H)M_{H,\mathrm M}(W_{m,0}),
\end{equation}
while
\begin{equation}
M_{H,\mathrm M}(W_{0,0})\le\frac{1}{2(1+\vartheta_{H,\mathrm M})}
\ell_{H,\mathrm M,\infty}^{\mathrm{unw}}(\eta_0,\eta_{H,\mathrm M}^\star)^2.
\end{equation}
Substitution proves the claim.
\end{proof}

\begin{corollary}[Boundary-iterate step size consequences]
\label{cor:fh-mtd-steps}
Under the hypotheses of Proposition~\ref{prop:fh-mtd-abstract}:

\textup{(a)} if $\alpha_k \equiv \alpha$ and
\begin{equation}
\alpha \le \frac{a_{H,\mathrm M,2}^{\mathrm{fh}}}{a_{H,\mathrm M,3}^{\mathrm{fh}}},
\end{equation}
then for all episodes $m \ge 0$,
\begin{equation}
\mathbb E\left[ \ell_{H,\mathrm M,\infty}^{\mathrm{unw}}(\eta_{mH},\eta_{H,\mathrm M}^\star)^2 \right] \le a_{H,\mathrm M,1}^{\mathrm{fh}} \ell_{H,\mathrm M,\infty}^{\mathrm{unw}}(\eta_0,\eta_{H,\mathrm M}^\star)^2 \bigl(1-a_{H,\mathrm M,2}^{\mathrm{fh}}H\alpha\bigr)^m + \frac{a_{H,\mathrm M,4}^{\mathrm{fh}}}{a_{H,\mathrm M,2}^{\mathrm{fh}}}\alpha.
\end{equation}

For the two diminishing-step cases below, where $g$ is the step-size offset, write $\tau_m:=mH+g+H-1$, so $\tau_0=g+H-1$.

\textup{(b)} if $\alpha_k = \alpha/(k+g)$, $\alpha > 1/a_{H,\mathrm M,2}^{\mathrm{fh}}$, and
\begin{equation}
g \ge \max\left\{H-1,1,\alpha,\frac{\alpha a_{H,\mathrm M,3}^{\mathrm{fh}}}{a_{H,\mathrm M,2}^{\mathrm{fh}}}\right\},
\end{equation}
then the boundary iterates satisfy
\begin{equation}
\mathbb E\left[ \ell_{H,\mathrm M,\infty}^{\mathrm{unw}}(\eta_{mH},\eta_{H,\mathrm M}^\star)^2 \right] \le a_{H,\mathrm M,1}^{\mathrm{fh}}\ell_{H,\mathrm M,\infty}^{\mathrm{unw}}(\eta_0,\eta_{H,\mathrm M}^\star)^2 \left(\frac{\tau_0}{\tau_m}\right)^{a_{H,\mathrm M,2}^{\mathrm{fh}}\alpha}+\frac{a_{H,\mathrm M,4}^{\mathrm{fh}}H^2\alpha^2}{a_{H,\mathrm M,2}^{\mathrm{fh}}\alpha-1}\cdot\frac{1}{\tau_m}.
\end{equation}

\textup{(c)} if $\alpha_k = \alpha/(k+g)^z$ with $z \in (0,1)$ and
\begin{equation}
g \ge \max\left\{H-1,1,\alpha^{1/z}, \left(\frac{\alpha a_{H,\mathrm M,3}^{\mathrm{fh}}}{a_{H,\mathrm M,2}^{\mathrm{fh}}}\right)^{1/z}, \left(\frac{2z}{a_{H,\mathrm M,2}^{\mathrm{fh}}\alpha}\right)^{1/(1-z)} \right\},
\end{equation}
then the boundary iterates satisfy
\begin{equation}
\begin{gathered}
\mathbb E\left[ \ell_{H,\mathrm M,\infty}^{\mathrm{unw}}(\eta_{mH},\eta_{H,\mathrm M}^\star)^2 \right] \le a_{H,\mathrm M,1}^{\mathrm{fh}}\ell_{H,\mathrm M,\infty}^{\mathrm{unw}}(\eta_0,\eta_{H,\mathrm M}^\star)^2\cdot\exp\left(-\frac{a_{H,\mathrm M,2}^{\mathrm{fh}}\alpha}{1-z}\bigl(\tau_m^{1-z}-\tau_0^{1-z}\bigr)\right)\\+\frac{2a_{H,\mathrm M,4}^{\mathrm{fh}}H^2\alpha}{a_{H,\mathrm M,2}^{\mathrm{fh}}}\cdot\frac{1}{\tau_m^z}.
\end{gathered}
\end{equation}
\end{corollary}

\begin{proof}
For a constant step size, condition \eqref{eq:fh-step-regularity} is immediate. For either diminishing schedule, the requirement $g\ge H-1$ gives
\begin{equation}
\frac{\alpha_{mH+H-1}}{\alpha_{mH}}=\left(\frac{mH+g}{mH+g+H-1}\right)^z\ge\left(\frac{g}{g+H-1}\right)^z\ge\frac{1}{2},
\end{equation}
where $z=1$ in part \textup{(b)}. Thus the regularity hypothesis of Proposition~\ref{prop:fh-mtd-abstract} holds in all three cases. For part \textup{(a)}, substitution of $\bar\alpha_m = H\alpha$ and $\bar\alpha_m^{(2)} = H\alpha^2$ into Proposition~\ref{prop:fh-mtd-abstract} yields part \textup{(a)}.

For part \textup{(b)}, set $q=\tau_0/H$ and $\lambda=a_{H,\mathrm M,2}^{\mathrm{fh}}\alpha$. The bounds
\begin{equation}
\frac{H\alpha}{\tau_m} \le \bar\alpha_m \le \frac{H\alpha}{mH+g}, \qquad \bar\alpha_m^{(2)} \le \frac{H\alpha^2}{(mH+g)^2}.
\end{equation}
imply
\begin{equation}
\prod_{j=0}^{m-1}(1-a_{H,\mathrm M,2}^{\mathrm{fh}}\bar\alpha_j)\le \left(\frac{q}{m+q}\right)^\lambda.
\end{equation}
The same elementary product-sum estimate gives
\begin{equation}
\sum_{i=0}^{m-1}\bar\alpha_i^{(2)}\prod_{j=i+1}^{m-1}(1-a_{H,\mathrm M,2}^{\mathrm{fh}}\bar\alpha_j)\le \frac{H\alpha^2}{\lambda-1}\cdot\frac{1}{m+q},
\end{equation}
and the displayed bound follows from Proposition~\ref{prop:fh-mtd-abstract} after substituting $q=\tau_0/H$ and $\tau_m=H(m+q)$.

For part \textup{(c)}, keep the same $q$ and set $A=a_{H,\mathrm M,2}^{\mathrm{fh}}\alpha H^{1-z}$. The bounds
\begin{equation}
\frac{H\alpha}{\tau_m^z} \le \bar\alpha_m \le \frac{H\alpha}{(mH+g)^z}, \qquad \bar\alpha_m^{(2)} \le \frac{H\alpha^2}{(mH+g)^{2z}}.
\end{equation}
imply
\begin{equation}
\prod_{j=0}^{m-1}(1-a_{H,\mathrm M,2}^{\mathrm{fh}}\bar\alpha_j)\le \exp\left(-\frac{A}{1-z}\bigl((m+q)^{1-z}-q^{1-z}\bigr)\right).
\end{equation}
The lower bound on $g$ implies $q\ge (2z/A)^{1/(1-z)}$. The elementary polynomial product-sum estimate gives
\begin{equation}
\sum_{i=0}^{m-1}\bar\alpha_i^{(2)}\prod_{j=i+1}^{m-1}(1-a_{H,\mathrm M,2}^{\mathrm{fh}}\bar\alpha_j)\le \frac{2H\alpha^2}{A}\cdot\frac{1}{(m+q)^z}.
\end{equation}
Substituting the definitions of $A$ and $q$ into Proposition~\ref{prop:fh-mtd-abstract}, using $\tau_m=H(m+q)$, and using $H^{2z}\le H^2$, gives the displayed bound.
\end{proof}

\begin{proof}[Proof of Theorem~\ref{thm:undisc-mtd}]
Combine Proposition~\ref{prop:fh-mtd-abstract} with Corollary~\ref{cor:fh-mtd-steps}. Explicitly, take
\begin{equation}
\begin{gathered}
c_{\mathrm M}^H:=a_{H,\mathrm M,2}^{\mathrm{fh}},\qquad r_{H,\mathrm M}:=\frac{a_{H,\mathrm M,3}^{\mathrm{fh}}}{a_{H,\mathrm M,2}^{\mathrm{fh}}},\qquad \bar\alpha_{\mathrm M}^H:=\frac{1}{r_{H,\mathrm M}}, \\
C_{\mathrm M}^H:=\max\left\{a_{H,\mathrm M,1}^{\mathrm{fh}},\frac{a_{H,\mathrm M,4}^{\mathrm{fh}}}{a_{H,\mathrm M,2}^{\mathrm{fh}}},a_{H,\mathrm M,4}^{\mathrm{fh}}H^2,\frac{2a_{H,\mathrm M,4}^{\mathrm{fh}}H^2}{a_{H,\mathrm M,2}^{\mathrm{fh}}}\right\}, \\
g_{\mathrm M}^H(\alpha):=\left\lceil\max\left\{1,\alpha,\alpha r_{H,\mathrm M}\right\}\right\rceil, \\
g_{\mathrm M}^H(\alpha,z):=\left\lceil\max\left\{1,\alpha^{1/z},(\alpha r_{H,\mathrm M})^{1/z},\left(\frac{2z}{a_{H,\mathrm M,2}^{\mathrm{fh}}\alpha}\right)^{1/(1-z)}\right\}\right\rceil.
\end{gathered}
\end{equation}
Each main-text display follows directly from the corresponding part of Corollary~\ref{cor:fh-mtd-steps}.
\end{proof}

\begin{proof}[Proof of Corollary~\ref{cor:undisc-mtd}]
By Corollary~\ref{cor:fh-mtd-steps}\textup{(b)},
\begin{equation}
\mathbb E\left[ \ell_{H,\mathrm M,\infty}^{\mathrm{unw}}(\eta_{mH},\eta_{H,\mathrm M}^\star)^2 \right] = O\left(\frac{1}{m+1}\right).
\end{equation}
By Jensen's inequality,
\begin{equation}
\mathbb E\left[ \ell_{H,\mathrm M,\infty}^{\mathrm{unw}}(\eta_{mH},\eta_{H,\mathrm M}^\star) \right] = O\left(\frac{1}{\sqrt{m+1}}\right).
\end{equation}
Thus $m = O(\varepsilon^{-2})$ episodes suffice, with the explicit dependence on $H$ and coverage inherited from Corollary~\ref{cor:fh-mtd-steps}. The transition count is $k=mH$.
\end{proof}

\section{Further discussion of the results}\label{app:discus}

\noindent\textbf{Bounded and affine perturbation regimes.}
After the statewise isometric embeddings, both CTD and MTD fit the same asynchronous contractive stochastic-approximation template. The main difference is the geometry of the sampled perturbation. In the scalar categorical case, bounded supports together with the Cram\'er geometry yield uniform samplewise bounds in both the discounted and fixed-horizon undiscounted settings. In the discounted setting,
\begin{equation}
\lVert \widehat T_{\mathrm C}(U;s,(r,s'))-(\mathcal O_{\mathrm C}U)(s) \rVert_2 \le 2B_{\mathrm C}.
\end{equation}
In the fixed-horizon setting,
\begin{equation}
\lVert F_{H,\mathrm C}(U;s,(r,s')) \rVert_{H,2,\infty} \le 2B_{H,\mathrm C}.
\end{equation}
Thus CTD falls into a bounded-noise regime throughout the work. In particular, the i.i.d.\ conditional second moment is uniformly bounded in the discounted analysis, and the fixed-horizon episodewise argument inherits only support-radius constants. By contrast, in the multivariate signed-categorical case one obtains a centered affine perturbation bound. In the discounted setting,
\begin{equation}
\lVert \widehat T_{\mathrm M}(U;s,(r,s'))-(\mathcal O_{\mathrm M}U)(s) \rVert_2 \le 2\beta_{\mathrm M}\lVert U-U^\star \rVert_{2,\infty}+B_{\mathrm M}^\star,
\end{equation}
together with a local second-moment estimate
\begin{equation}
\mathbb E\left[ \lVert \widehat T_{\mathrm M}(U_k;S_k,(R_k,S'_k))-(\mathcal O_{\mathrm M}U_k)(S_k) \rVert_2^2 \,\mid\, U_k,S_k \right] \le A_{\mathrm M}^{\mathrm{iid}}\bigl(1+\lVert U_k-U^\star \rVert_{2,\infty}^2\bigr).
\end{equation}
In the fixed-horizon setting, the analogous residual bound is
\begin{equation}
\lVert F_{H,\mathrm M}(U;s,(r,s')) \rVert_{H,2,\infty} \le 2 \lVert U \rVert_{H,2,\infty}+B^{\mathrm{res}}_{H,\mathrm M},
\end{equation}
again with an affine conditional second-moment estimate. Consequently, the MTD theorem constants contain additional growth, bias, and absorption quantities, whereas the CTD bounds depend only on support-radius constants.

\noindent\textbf{Smoothing constants and the number of asynchronous blocks.}
The finite-iteration analysis is driven by a block-supremum contraction norm and a smoothed block-$\ell_p$ potential. The only explicit dimension loss in this smoothing step comes from the comparison between $\lVert \cdot \rVert_{2,\infty}$ and $\lVert \cdot \rVert_{2,p}$, and therefore depends only on the number of asynchronously updated blocks. In the discounted appendices, those blocks are indexed only by the state variable and the relevant count is $\lvert \X \rvert$. In the fixed-horizon appendices, the stacked process is first flattened over horizon and state, so the relevant count becomes $\lvert \mathcal S_H \rvert = H\lvert \X \rvert$. The inner representation dimension $d$, and in the multivariate case the reward dimension $q$, still affect problem-dependent constants, but they do not enter through the Moreau-envelope smoothing comparison itself.

\noindent\textbf{Discounted and fixed-horizon contraction mechanisms.}
The fixed-horizon undiscounted theorems results are not obtained by setting $\gamma=1$ in the discounted analysis. In the discounted setting, the contraction comes directly from the Bellman discount factor and each update modifies only one sampled state block. In the fixed-horizon setting there is no discount-driven contraction. Instead, the contraction is recovered by weighting the horizon layers and exploiting the fact that horizon $h$ bootstraps from horizon $h-1$. The corresponding online recursion is also structurally different: one sampled transition updates the entire horizon stack at the sampled state. Thus the fixed-horizon setting is different both in its contraction mechanism and in its stochastic approximation structure.

\noindent\textbf{Weighted and unweighted fixed-horizon errors.}
The fixed-horizon proofs use the weighted metric, but Theorems~\ref{thm:undisc-ctd} and~\ref{thm:undisc-mtd} are stated in the natural unweighted metric. If
\begin{equation}
\ell_{H,\infty}^{\mathrm{unw}}(\eta,\eta'):=\max_{1\le h\le H,\ s\in\X}\ell\bigl(\eta^h(s),{\eta'}^h(s)\bigr),
\end{equation}
then
\begin{equation}
\ell_{H,\infty}(\eta,\eta')\le\ell_{H,\infty}^{\mathrm{unw}}(\eta,\eta')\le\lambda^{-H}\ell_{H,\infty}(\eta,\eta').
\end{equation}
With $\lambda=H/(H+1)$ the conversion factor satisfies $\lambda^{-H}\le e$, hence its dependence on H is uniformly bounded. The contraction gap is $1-\lambda=1/(H+1)$, which enters the explicit constants polynomially through $\kappa_H=\rho_{H,\min}/(2(H+1))$. Each sampled transition updates all $H$ horizon layers, so a dense implementation has arithmetic cost proportional to $H$ times the cost of one categorical block update.

\noindent\textbf{Interpretation of the step size regimes.}
Across all theorem statements, constant step sizes yield geometric convergence to a controllable $O(\alpha)$ neighborhood of the projected fixed point. Linearly-diminishing step sizes yield $O(1/k)$ squared error and hence $O(\varepsilon^{-2})$ sample complexity for mean error at most $\varepsilon$. Polynomially-diminishing step sizes trade a weaker asymptotic decay for milder admissibility requirements. In the discounted Markovian case, the Poisson-equation argument preserves these rates without a logarithmic mixing-time factor. The constant $A_Q$ records the temporal dependence of the trajectory.

\end{document}
